\chapter{Dynamics: kinematics and Newton's laws}
\label{vol03:ch04}

\epigraph{The quantity of motion is the measure of the same, arising
from the velocity and quantity of matter conjunctly.}{Isaac Newton,
\emph{Philosophiae Naturalis Principia Mathematica} (1687), Definition
II, in the Cohen and Whitman translation \cite[p.~404]{hist:newton1687}.}

\chapmeta{Half-life: durable. Archetypes: balance, transport.
Exercise target: 35. Project track: physical observation plus
analysis (hybrid). Hazard class: low (thrown objects in controlled
space). Reader path default: standard, with sections explicitly
tagged. Named cases: none.}

\noindent
Chapters 1, 2, and 3 of this volume have treated static loading: the
forces, structures, and stresses that occur when bodies are at rest
or in steady motion. This chapter takes the next step: bodies that
accelerate. Dynamics replaces the right-hand-zero of the static
equilibrium equation with the body's mass times its acceleration, and
the moment equation with the rate of change of its angular momentum.
The replacement looks small on the page; it opens the analysis of
projectiles, vehicles, vibrations, rotating machinery, and every
mechanical system whose state changes with time. The treatment here
follows the engineering-mechanics tradition; Meriam and Kraige give
the standard fuller account at problem-set depth
\cite{text:meriam-kraige2018}.

\section{Position, velocity, acceleration}\pathtag{core}

A particle's motion in three dimensions is described by its position
vector $\vec{r}(t)$, a function of time. Velocity is the time
derivative,
\[
\vec{v}(t) = \frac{d\vec{r}}{dt},
\]
and acceleration the second time derivative,
\[
\vec{a}(t) = \frac{d\vec{v}}{dt} = \frac{d^2 \vec{r}}{dt^2}.
\]
Velocity is tangent to the trajectory at every instant; the
acceleration's direction is set by the rate at which the velocity
changes in both magnitude and direction, with no general requirement
that it be tangent, perpendicular, or in any other simple geometric
relationship to the velocity.

In Cartesian components, $\vec{r} = x\hat{\imath} + y\hat{\jmath} +
z\hat{k}$, and the time derivatives act componentwise. In polar
coordinates in two dimensions, $\vec{r} = r\hat{r}$, the velocity
has both a radial component and a transverse component because the
basis vectors $\hat{r}$ and $\hat{\theta}$ themselves rotate with the
particle:
\[
\vec{v} = \dot{r}\hat{r} + r\dot{\theta}\hat{\theta},
\]
where the dot denotes time derivative. Acceleration in polar
coordinates has four terms:
\[
\vec{a} = (\ddot{r} - r\dot{\theta}^2)\hat{r} + (r\ddot{\theta} + 2
\dot{r}\dot{\theta})\hat{\theta}.
\]
The four terms have standard names: the radial term $-r\dot{\theta}^2$
is the centripetal acceleration (directed inward for circular motion);
the transverse term $r\ddot{\theta}$ is the angular acceleration; the
$\ddot{r}$ term is the rate of change of radial speed; and the
$2\dot{r}\dot{\theta}$ term is the Coriolis acceleration, which
appears whenever a particle moves radially in a rotating frame. The
Coriolis term is real (it has the dimensions of acceleration and
contributes to the body's force balance), and its physical meaning
becomes clearest in the rotating-reference-frame discussion of
section 4.3.

Three coordinate systems handle most engineering problems: Cartesian
for trajectories in uniform fields, polar (or cylindrical) for
rotating systems and central-force problems, and natural (path)
coordinates $(\hat{t}, \hat{n})$ for curvilinear motion along a known
trajectory. In path coordinates,
\[
\vec{v} = v\hat{t}, \quad \vec{a} = \dot{v}\hat{t} + \frac{v^2}{\rho}
\hat{n},
\]
where $\rho$ is the local radius of curvature of the trajectory. The
tangential component $\dot{v}$ changes the speed; the normal component
$v^2/\rho$ changes the direction. The decomposition matters in vehicle
dynamics, where the tangential acceleration is the engine's
contribution and the normal acceleration is the cornering acceleration
that the tyres must supply.

\paragraph{Deriving the polar components.} The four-term polar
acceleration is stated above; the derivation is short and worth carrying
in full, because every term traces to a single geometric fact about the
rotating basis. The polar unit vectors are
\[
\hat{r}=\cos\theta\,\hat{\imath}+\sin\theta\,\hat{\jmath}, \qquad
\hat{\theta}=-\sin\theta\,\hat{\imath}+\cos\theta\,\hat{\jmath},
\]
and differentiating each with respect to $\theta$ rotates it a quarter
turn:
\[
\frac{d\hat{r}}{d\theta}=-\sin\theta\,\hat{\imath}+\cos\theta\,\hat{\jmath}
=\hat{\theta}, \qquad
\frac{d\hat{\theta}}{d\theta}=-\cos\theta\,\hat{\imath}-\sin\theta\,
\hat{\jmath}=-\hat{r}.
\]
Figure~\ref{fig:vol03ch04:polar-basis} draws the two relations. By the
chain rule the time derivatives are $\dot{\hat{r}}=\dot{\theta}
\hat{\theta}$ and $\dot{\hat{\theta}}=-\dot{\theta}\hat{r}$, the only
ingredients the rest of the derivation needs. Start from the position
$\vec{r}=r\hat{r}$ and differentiate once, using the product rule:
\[
\vec{v}=\dot{r}\hat{r}+r\dot{\hat{r}}=\dot{r}\hat{r}+r\dot{\theta}
\hat{\theta},
\]
the radial term from the changing length and the transverse term from
the turning direction. Differentiate again, applying the product rule to
each of the two terms and substituting $\dot{\hat{r}}=\dot{\theta}
\hat{\theta}$ and $\dot{\hat{\theta}}=-\dot{\theta}\hat{r}$:
\begin{align*}
\vec{a}&=\frac{d}{dt}\bigl(\dot{r}\hat{r}\bigr)
+\frac{d}{dt}\bigl(r\dot{\theta}\hat{\theta}\bigr)\\
&=\bigl(\ddot{r}\hat{r}+\dot{r}\dot{\theta}\hat{\theta}\bigr)
+\bigl(\dot{r}\dot{\theta}\hat{\theta}+r\ddot{\theta}\hat{\theta}
-r\dot{\theta}^{2}\hat{r}\bigr).
\end{align*}
Collecting the $\hat{r}$ and $\hat{\theta}$ parts gives the four-term
result
\[
\vec{a}=(\ddot{r}-r\dot{\theta}^{2})\hat{r}+(r\ddot{\theta}+2\dot{r}
\dot{\theta})\hat{\theta}.
\]
The Coriolis term $2\dot{r}\dot{\theta}$ arises twice, once from
differentiating $\dot{r}$ in the transverse direction and once from
differentiating $\hat{\theta}$, which is why it carries the factor of
two. The same product-rule discipline gives the cylindrical case by
appending an independent $\ddot{z}\hat{k}$, and it gives the
path-coordinate split $\vec{a}=\dot{v}\hat{t}+(v^{2}/\rho)\hat{n}$ once
the turning rate of the tangent $\hat{t}$ is written as
$d\hat{t}/ds=\hat{n}/\rho$, the Frenet relation that defines the radius
of curvature $\rho$.

% Figure: the rotating polar basis. The unit vectors r-hat and theta-hat
% at a point turn with the particle, and their time derivatives point
% along each other, which is the geometric origin of the extra terms in
% the polar-coordinate velocity and acceleration.
\begin{figure}[htbp]
\centering
\begin{tikzpicture}[
  vec/.style={-{Stealth[length=2.6mm]}, very thick},
  lbl/.style={font=\small},
]
% origin
\coordinate (O) at (0,0);
\fill[engblack] (O) circle (1.4pt);
\node[lbl, below left] at (O) {$O$};
% reference axes
\draw[enggray, -{Stealth[length=2mm]}] (O) -- (4.6,0) node[right, lbl]{$x$};
\draw[enggray, -{Stealth[length=2mm]}] (O) -- (0,4.0) node[above, lbl]{$y$};
% particle position at angle theta = 40 deg, radius 3.4
\coordinate (P) at (40:3.4);
\draw[engblue, very thick] (O) -- (P);
\node[engblue, lbl] at (20:1.9) {$r$};
% angle arc
\draw[enggray] (1.0,0) arc (0:40:1.0);
\node[lbl] at (20:1.3) {$\theta$};
% particle dot
\fill[engred] (P) circle (2.6pt);
\node[engred, lbl, above right] at (P) {particle};
% r-hat along the radius
\draw[vec, draw=englime!70!black] (P) -- ($(P)+(40:1.1)$);
\node[englime!60!black, lbl] at ($(P)+(40:1.35)$) {$\hat{r}$};
% theta-hat perpendicular, in direction of increasing theta (+90 deg)
\draw[vec, draw=engamber] (P) -- ($(P)+(130:1.1)$);
\node[engamber, lbl] at ($(P)+(130:1.35)$) {$\hat{\theta}$};
% small curved arrows showing that the basis rotates as theta grows:
% d r-hat / d theta points along + theta-hat; d theta-hat / d theta points along - r-hat
\draw[engred, -{Stealth[length=2mm]}, thick]
  ($(P)+(40:1.1)$) arc (40:70:1.1);
\node[engred, lbl] at ($(P)+(70:1.55)$) {$\frac{d\hat{r}}{d\theta}=\hat{\theta}$};
\end{tikzpicture}
\caption[The rotating polar basis]{The polar unit vectors at a moving
particle. The radial vector $\hat{r}$ points outward along the line from
$O$; the transverse vector $\hat{\theta}$ points in the direction of
increasing $\theta$, perpendicular to $\hat{r}$. As the particle moves
the angle $\theta$ changes, and the basis turns with it: differentiating
$\hat{r}$ with respect to $\theta$ rotates it a quarter turn forward into
$\hat{\theta}$, and differentiating $\hat{\theta}$ rotates it into
$-\hat{r}$. These two facts, $d\hat{r}/d\theta=\hat{\theta}$ and
$d\hat{\theta}/d\theta=-\hat{r}$, are the whole reason the polar velocity
and acceleration carry terms beyond the Cartesian $\dot{x}$ and
$\ddot{x}$: the basis itself accelerates.}
\label{fig:vol03ch04:polar-basis}
\end{figure}


\begin{principle}[Kinematics is geometry plus time]
Kinematics describes motion without reference to its cause. The
description requires only a position function and the rules of
differentiation. The choice of coordinate system is the engineer's;
the position function is the body's. The same trajectory expressed
in Cartesian, polar, or path coordinates is the same physical motion;
the components of velocity and acceleration are different in each
system because the components of any vector are coordinate-dependent.
The extra terms in the polar and path forms are not new physics; they
are the derivatives of a basis that turns with the motion.
\end{principle}

\paragraph{Reading kinematics off a graph.} Most one-dimensional
motion an engineer meets is piecewise: a vehicle accelerates, then
holds speed, then brakes; an actuator ramps, dwells, and returns. The
three quantities stack. Velocity is the running integral of
acceleration, so a step in $a$ becomes a corner in $v$; position is
the running integral of velocity, so a corner in $v$ becomes a change
of curvature in $x$. The inverse reading is equally useful: the slope
of the position curve is the velocity, the slope of the velocity curve
is the acceleration, and the area under the velocity curve over an
interval is the displacement over that interval.

Consider a car that accelerates from rest at a constant
$3\,\si{\meter\per\second\squared}$ for $10\,\si{\second}$, then brakes
at a constant $2\,\si{\meter\per\second\squared}$ until it stops. The
acceleration is a two-level step; the velocity is two straight
segments meeting in a corner at $t=10\,\si{\second}$, where the speed
peaks at $v_{\max}=3\times 10=30\,\si{\meter\per\second}$; and the
position is two parabolic arcs of opposite curvature. The braking
phase lasts $v_{\max}/2=15\,\si{\second}$, so the motion ends at
$t=25\,\si{\second}$. Total distance is the area under the velocity
triangle, $\tfrac{1}{2}\times 25\times 30=375\,\si{\meter}$.
Figure~\ref{fig:vol03ch04:kinematics-1d} shows all three panels on a
shared time axis; the same problem appears as a calculation exercise,
where the algebra reproduces the graphical reading.

% Figure: position, velocity, acceleration versus time for the two-phase
% car motion (accelerate then decelerate). Stacked panels share a time axis.
\begin{figure}[htbp]
\centering
\begin{tikzpicture}
\begin{axis}[
  name=accax,
  width=11cm, height=3.2cm,
  ylabel={$a$ (\si{\meter\per\second\squared})},
  xmin=0, xmax=25, ymin=-2.5, ymax=3.5,
  axis lines=left, xtick=\empty,
  grid=major, grid style={enggray!20},
  tick label style={font=\scriptsize}, label style={font=\small},
]
% phase 1: a = +3 for 0..10; phase 2: a = -2 for 10..25
\addplot[engred, very thick] coordinates {(0,3) (10,3)};
\addplot[engred, very thick] coordinates {(10,-2) (25,-2)};
\addplot[engred, very thick, dashed] coordinates {(10,3) (10,-2)};
\end{axis}
\begin{axis}[
  name=velax, at={(accax.below south west)}, anchor=above north west, yshift=-0.3cm,
  width=11cm, height=3.2cm,
  ylabel={$v$ (\si{\meter\per\second})},
  xmin=0, xmax=25, ymin=0, ymax=35,
  axis lines=left, xtick=\empty,
  grid=major, grid style={enggray!20},
  tick label style={font=\scriptsize}, label style={font=\small},
]
% v rises 0->30 over 0..10, falls 30->0 over 10..25
\addplot[engblue, very thick] coordinates {(0,0) (10,30) (25,0)};
\end{axis}
\begin{axis}[
  name=posax, at={(velax.below south west)}, anchor=above north west, yshift=-0.3cm,
  width=11cm, height=3.2cm,
  xlabel={time $t$ (\si{\second})}, ylabel={$x$ (\si{\meter})},
  xmin=0, xmax=25, ymin=0, ymax=400,
  axis lines=left,
  grid=major, grid style={enggray!20},
  tick label style={font=\scriptsize}, label style={font=\small},
]
% x: integral of v. Phase 1: x=1.5 t^2 (0..10 -> 150). Phase 2: x=150+30(t-10)-1(t-10)^2.
\addplot[engamber, very thick, domain=0:10, samples=40] {1.5*x^2};
\addplot[engamber, very thick, domain=10:25, samples=60] {150 + 30*(x-10) - 1*(x-10)^2};
\end{axis}
\end{tikzpicture}
\caption[Stacked kinematics of two-phase motion]{Acceleration,
velocity, and position against a shared time axis for the
two-phase car of the second calculation exercise: constant
$+3\,\si{\meter\per\second\squared}$ for the first
$10\,\si{\second}$, then constant $-2\,\si{\meter\per\second\squared}$
to rest at $t=25\,\si{\second}$. Velocity is the running integral of
acceleration, so the step in $a$ becomes a kink in $v$; position is
the running integral of velocity, so the kink in $v$ becomes a change
of curvature in $x$. The area under the velocity triangle,
$\tfrac{1}{2}(25)(30)=375\,\si{\meter}$, is the total distance, a
reading that the integral confirms. Reading these three panels
together is the core skill of one-dimensional kinematics.}
\label{fig:vol03ch04:kinematics-1d}
\end{figure}


\section{Rotational kinematics}\pathtag{core}

A rigid body's orientation is described by three angles (in three
dimensions) or one angle (in two dimensions). The time derivatives
of those angles are the body's angular velocity and angular
acceleration. In two dimensions, with $\theta$ the body's angular
position about the out-of-plane axis,
\[
\omega = \dot{\theta}, \quad \alpha = \ddot{\theta},
\]
the angular velocity and angular acceleration. The units are
$\si{\radian\per\second}$ and $\si{\radian\per\second\squared}$
respectively; in engineering practice the units of $\omega$ are often
RPM (revolutions per minute), with the conversion $\omega \,[\si{
\radian\per\second}] = \omega \,[\text{RPM}] \times 2\pi / 60$.

For a point $P$ on a rigid body rotating about a fixed axis with
angular velocity $\omega$, the linear velocity of $P$ is
\[
\vec{v}_P = \vec{\omega} \times \vec{r}_P,
\]
where $\vec{r}_P$ is the position of $P$ from any point on the axis
and $\vec{\omega}$ is the angular velocity vector (directed along the
axis by the right-hand rule, with magnitude $\omega$). The linear
acceleration has tangential and centripetal components:
\[
\vec{a}_P = \vec{\alpha} \times \vec{r}_P + \vec{\omega} \times
(\vec{\omega} \times \vec{r}_P),
\]
the first being the tangential contribution and the second the
centripetal (directed inward toward the rotation axis).

A rigid body in general planar motion combines translation of a
reference point with rotation about that point. If point $A$ has
velocity $\vec{v}_A$ and the body has angular velocity $\omega$ about
the out-of-plane axis, then any other point $B$ on the body has
velocity
\[
\vec{v}_B = \vec{v}_A + \vec{\omega} \times \vec{r}_{B/A},
\]
where $\vec{r}_{B/A}$ is the vector from $A$ to $B$. The formula
is the basis of the velocity-image method for analysing mechanisms
and is used heavily in linkage and gear analysis (chapter 14 of
Volume VIII).

The \engterm{instantaneous centre of rotation} is the point on (or
extended from) a rigid body at which the velocity is instantaneously
zero. For a body in general planar motion at an instant, every other
point on the body has velocity directed perpendicular to its line
from the instantaneous centre and with magnitude proportional to its
distance from the centre. The instantaneous centre changes its
position on the body with time; the geometric locus of these
positions is the body's centrode. Identifying the instantaneous
centre lets the engineer analyse the body's motion at that instant
as pure rotation about that centre, which is sometimes simpler than
the translation-plus-rotation form.

The cleanest instance is a wheel rolling without slipping. The no-slip
condition is exactly the statement that the contact point has zero
velocity at the instant of contact, which makes the contact point the
instantaneous centre. Every other point on the wheel then moves
perpendicular to its line from the contact, with speed proportional to
its distance from it: the centre, at distance $R$, moves at
$v_{G}=\omega R$; the top, at distance $2R$, moves at $2\omega R$; the
leading and trailing edges, at distance $\sqrt{2}\,R$, move at
$\sqrt{2}\,\omega R$. Figure~\ref{fig:vol03ch04:rolling-velocity}
draws the field. The same $v_{G}=\omega R$ relation, differentiated,
gives the rolling constraint $a_{G}=\alpha R$ used throughout the
rigid-body dynamics that follows.

% Figure: velocity field of a wheel rolling without slipping, showing the
% instantaneous centre at the contact point and velocities at key points.
\begin{figure}[htbp]
\centering
\begin{tikzpicture}[
  vel/.style={draw=engblue, -{Stealth}, very thick},
  lbl/.style={font=\small},
]
% ground
\draw[engblack, thick] (-2.2,0) -- (4.2,0);
% wheel of radius 1.4 centred at (1,1.4)
\coordinate (O) at (1,1.4);
\coordinate (P) at (1,0);     % contact (instantaneous centre)
\coordinate (T) at (1,2.8);   % top
\coordinate (F) at (2.4,1.4); % front (right)
\coordinate (B) at (-0.4,1.4);% back (left)
\draw[draw=engblack, thick, fill=enggray!10] (O) circle (1.4);
\foreach \pt in {O,T,F,B} {\fill[engblack] (\pt) circle (1.3pt);}
\fill[engred] (P) circle (1.8pt);
% velocity at centre: v = omega R, to the right
\draw[vel] (O) -- ++(1.6,0) node[above, lbl] {$\vec{v}_{G}=\omega R$};
% velocity at top: 2 v
\draw[vel] (T) -- ++(3.0,0) node[above, lbl] {$\vec{v}_{\text{top}}=2\omega R$};
% velocity at front: directed up-back at 45 (magnitude sqrt2 v), but
% draw vertical-ish for clarity: actually v_F = omega x r, r from P to F.
% From instantaneous centre P=(1,0): F=(2.4,1.4), v perpendicular to PF.
\draw[vel] (F) -- ++(0.95,-0.95) node[right, lbl] {$\sqrt{2}\,\omega R$};
% velocity at back: from P=(1,0) to B=(-0.4,1.4), perpendicular
\draw[vel] (B) -- ++(0.95,0.95) node[left, lbl] {$\sqrt{2}\,\omega R$};
% contact point: zero velocity, mark
\node[engred, font=\small, below] at (P) {$\vec{v}_{P}=\vec{0}$};
\node[engred, font=\scriptsize, below right] at ($(P)+(0.05,-0.45)$) {instantaneous centre};
% omega arc
\draw[draw=engamber, -{Stealth}, thick] (O) ++(140:0.55) arc (140:-30:0.55);
\node[engamber, font=\small] at ($(O)+(0.0,0.78)$) {$\omega$};
\end{tikzpicture}
\caption[Velocity field of a rolling wheel]{Velocity of selected
points on a wheel rolling without slipping. The contact point $P$
(red) is the instantaneous centre of rotation: its velocity is zero
at that instant, which is exactly the no-slip condition. Every other
point moves perpendicular to its line from $P$, with speed
proportional to its distance from $P$. The centre $G$, at distance
$R$ from $P$, moves at $\omega R$; the top, at distance $2R$, moves
at $2\omega R$; the leading and trailing points, at distance
$\sqrt{2}\,R$, move at $\sqrt{2}\,\omega R$ directed perpendicular to
their radius from $P$. Treating the motion as pure rotation about the
instantaneous centre reproduces the translation-plus-rotation result
$\vec{v}_{B}=\vec{v}_{G}+\vec{\omega}\times\vec{r}_{B/G}$ with less
bookkeeping.}
\label{fig:vol03ch04:rolling-velocity}
\end{figure}


\paragraph{When rolling gives way to slipping.} The no-slip relation
$v_{G}=\omega R$ is a constraint, not a law, and it holds only while the
contact patch can supply the friction the constraint demands. A wheel
driven or braked hard enough breaks the constraint and begins to skid,
and the boundary between the two regimes is a friction check. For a
wheel of radius $R$ carrying a normal load $N$ under a drive torque
$\tau$ at the axle, rolling without slip requires a contact friction
$f=\tau/R$ that the tyre must provide; the tyre can provide at most
$\mu_{s}N$, so rolling persists only while $\tau\le\mu_{s}NR$. Past that
torque the friction saturates at the kinetic value $\mu_{k}N$, the wheel
spins faster than the ground beneath it, and $v_{G}<\omega R$: the
contact point no longer has zero velocity, the instantaneous centre
leaves the contact, and the tidy velocity field of the rolling wheel is
gone. The same boundary read the other way governs a locked brake: a
wheel braked past $\mu_{s}NR$ stops turning while the car still moves, so
$\omega=0$ while $v_{G}>0$, the skidding-tyre condition that the braking
estimate and the accident reconstruction both rest on. The lesson is
that rolling is the privileged low-loss regime a designer wants to stay
inside, and the anti-lock brake exists to hold the tyre at the edge of
that regime rather than let it cross into the high-loss skid.

\paragraph{Gear trains: angular velocity in ratio.} Two gears in mesh,
or two pulleys joined by a belt, share the linear speed at the point of
contact, which ties their angular velocities in inverse proportion to
their radii. For gears of pitch radii $R_{1}$ and $R_{2}$ (tooth counts
$N_{1}$ and $N_{2}$, which are proportional to the radii), the meshing
condition $\omega_{1}R_{1}=\omega_{2}R_{2}$ gives the gear ratio
\[
\frac{\omega_{2}}{\omega_{1}}=\frac{R_{1}}{R_{2}}=\frac{N_{1}}{N_{2}}.
\]
A small gear driving a large one turns slower and, as the moment
equation will show, turns with more torque: the product $\tau\omega$,
the transmitted power, is preserved across an ideal mesh, so a ratio
that divides the speed multiplies the torque by the same factor. A
compound train stacks the ratios multiplicatively, one factor per mesh,
which is how a gearbox spans from a launch ratio that multiplies engine
torque for the standing start to an overdrive ratio that trades torque
for road speed at cruise. The angular-acceleration relation
differentiates the velocity relation unchanged, $\alpha_{2}/\alpha_{1}=
N_{1}/N_{2}$, so a train that steps speed down steps angular acceleration
down in the same ratio, the fact that lets a small motor spin up a large
inertial load through reduction gearing.

In three dimensions, angular velocity is a vector but angular position
is not (the angles do not commute under composition; finite rotations
are described by orthogonal matrices, quaternions, or Euler angles).
The non-commutation is concrete: rotate a book $90^{\circ}$ about a
vertical axis and then $90^{\circ}$ about a horizontal axis, and it
lands in a different orientation than if the two rotations are applied
in the opposite order. Infinitesimal rotations do commute to first
order, which is why their rates, the angular velocities, add as
vectors even though the finite rotations they integrate to do not. The
non-vector character of angular position is the source of the
gimbal-lock problem in three-axis gyroscopic instruments, where two of
the three gimbal axes can align and the instrument loses a degree of
freedom, and it is the mathematical reason rotation in three dimensions
is more difficult than in two. The same non-commutation forces
attitude software to track orientation with a quaternion or a rotation
matrix rather than three stored angles, because the three-angle
representation has the gimbal singularity built in. The chapter
restricts attention to two-dimensional and fixed-axis problems;
chapter 4 of Volume VIII returns to the full three-dimensional case in
the context of rotating machinery.

\section{Newton's second law applied}\pathtag{core}

The three laws work as a system, and the second is the one that carries
the arithmetic. The \engterm{first law} fixes the reference frames in
which the others hold: a body free of net force keeps its velocity, so
a frame in which a free body moves uniformly is an inertial frame, and
the laws are stated in such frames (the mastery box on non-inertial
frames below collects the correction terms that appear when the frame
accelerates). The \engterm{third law} supplies the bookkeeping for
interacting bodies: every force one body exerts on another is matched
by an equal and opposite force on itself, which is what makes the
internal forces of a system cancel in pairs and what underwrites the
conservation of momentum derived later in the chapter. The
\engterm{second law} is the quantitative core, and the rest of the
section applies it.

Newton's second law states that the net force on a particle equals
the rate of change of its linear momentum,
\[
\vec{F}_{\text{net}} = \frac{d\vec{p}}{dt}.
\]
For a particle of constant mass, $\vec{p} = m\vec{v}$ and the law
reduces to
\[
\vec{F}_{\text{net}} = m\vec{a},
\]
the familiar working form. For a system whose mass changes with time
(a rocket expelling propellant, a sled accumulating snow, a conveyor
belt picking up material), the momentum form is required and the
$m\vec{a}$ form is wrong; the difference is treated in the
impulse-momentum section below.

For a rigid body in plane motion, Newton's second law for the
centre of mass combines with the moment equation about the centre of
mass:
\begin{align*}
\sum \vec{F}_{\text{ext}} &= m\vec{a}_G, \\
\sum M_G &= I_G \alpha,
\end{align*}
where $\vec{a}_G$ is the centre-of-mass acceleration, $\sum M_G$ is
the net external moment about the centre of mass, $I_G$ is the
body's mass moment of inertia about the centre-of-mass axis, and
$\alpha$ is the angular acceleration about that axis. The mass moment
of inertia (distinct from the area moment of inertia of chapter 2)
is
\[
I_G = \int_V r^2 \,dm,
\]
the integral of $r^2$ over the mass distribution, with $r$ the
distance from the axis. Standard values for common bodies:

\begin{itemize}[leftmargin=*]
  \item Solid cylinder of mass $m$ and radius $R$ about its axis:
    $I = \tfrac{1}{2} m R^2$.
  \item Thin-walled cylinder of mass $m$ and radius $R$ about its
    axis: $I = m R^2$.
  \item Solid sphere of mass $m$ and radius $R$ about a diameter:
    $I = \tfrac{2}{5} m R^2$.
  \item Thin rod of mass $m$ and length $L$ about its centre: $I =
    \tfrac{1}{12} m L^2$.
  \item Thin rod of mass $m$ and length $L$ about one end: $I =
    \tfrac{1}{3} m L^2$.
\end{itemize}

The \engterm{parallel-axis theorem} for mass moments: $I = I_G + m
d^2$, where $I$ is the moment about an axis at perpendicular distance
$d$ from a parallel axis through the centre of mass and $I_G$ is the
moment about the centre-of-mass axis. The theorem follows from the
definition. Write a mass element's position from the offset axis as the
sum of the position of the centre of mass and the element's position
$\vec{r}'$ from the centre of mass; squaring and integrating gives $I=
\int|\vec{d}+\vec{r}'|^{2}\,dm=d^{2}\!\int dm+2\vec{d}\cdot\!\int\vec{r}'
\,dm+\int r'^{2}\,dm$. The middle term vanishes because $\int\vec{r}'\,dm
=\vec{0}$ by the definition of the centre of mass, and the surviving
terms are $md^{2}$ and $I_{G}$. The vanishing of the cross term is the
same fact that made the rigid-body kinetic energy split cleanly about the
centre of mass, and it is the reason the centre of mass is the privileged
point for both.

\paragraph{The radius of gyration collapses the mass distribution to one
length.} A body's moment of inertia is often written $I=mk^{2}$, defining
the \engterm{radius of gyration} $k=\sqrt{I/m}$, the single radius at
which the whole mass could be concentrated as a thin ring to give the
same moment of inertia. The radius of gyration is a geometric property of
the shape, independent of the mass: for a solid cylinder $k=R/\sqrt{2}$,
for a thin ring $k=R$, for a solid sphere $k=R\sqrt{2/5}$, for a rod
about its end $k=L/\sqrt{3}$. It reads directly the shape factor that
governs how a body rolls and spins: the rolling-acceleration result
$a_{G}=g\sin\beta/(1+k^{2}/R^{2})$ of the incline example depends on the
body only through $k^{2}/R^{2}$, so a hoop ($k^{2}/R^{2}=1$) rolls slower
than a sphere ($k^{2}/R^{2}=2/5$) because its mass sits farther from the
axis. The same ratio measures how much rotational energy a rolling body
hides: at rolling speed $v_{G}=\omega R$ the kinetic energy is
$\tfrac{1}{2}mv_{G}^{2}(1+k^{2}/R^{2})$, so a hoop carries half again the
translational energy in spin while a sphere carries only two-fifths.
Where a designer wants a flywheel to store the most energy per unit mass
at a given rim speed, a large $k$ is the target and the rim-loaded shape
wins; where a wheel must accelerate quickly, a small $k$ is wanted and
the mass is drawn toward the axle. The radius of gyration is the one
number that carries the mass distribution into every rotational balance
the chapter writes.

\paragraph{Free-body and kinetic diagrams.} The discipline that keeps
rigid-body dynamics from collapsing into sign errors is the paired
diagram. The free-body diagram shows every external force and moment
acting on the body; the kinetic diagram shows the inertia terms
$m\vec{a}_{G}$ and $I_{G}\alpha$. Newton's second law is the statement
that the two diagrams are equal, drawn side by side and summed
component by component. Keeping $m\vec{a}_{G}$ on the kinetic side,
never the free-body side, prevents the most common beginner's error,
which is to draw $m\vec{a}$ as though it were an applied force and then
add it twice. Figure~\ref{fig:vol03ch04:incline-fbd} shows the pair
for a cylinder rolling down an incline.

% Figure: free-body diagram plus kinetic diagram for a cylinder rolling
% down an incline. Left: forces. Right: m a_G and I alpha.
\begin{figure}[htbp]
\centering
\begin{tikzpicture}[
  force/.style={draw=engred, -{Stealth}, very thick},
  kin/.style={draw=engblue, -{Stealth}, very thick},
  surf/.style={draw=engblack, thick},
  lbl/.style={font=\small},
]
% == Left panel: free-body diagram ==
\begin{scope}[shift={(0,0)}]
% incline: angle ~ 25 deg for clarity
\draw[surf] (-0.3,0) -- (4.3,2.15);
\draw[surf] (-0.3,0) -- (-0.3,-0.2);
% angle arc
\draw (1.0,0.0) ++(0:0.9) arc (0:26.6:0.9);
\node[font=\scriptsize] at (1.95,0.22) {$\beta$};
% cylinder centre on incline
\coordinate (C) at (2.2,1.55);
\draw[draw=engblack, thick, fill=enggray!12] (C) circle (0.55);
\node[lbl] at (C) {$G$};
% weight (down)
\draw[force] (C) -- ++(0,-1.4) node[below, lbl] {$m\vec{g}$};
% normal (perpendicular to incline, pointing away)
\draw[force] (C) -- ++(116.6:1.25) node[above left, lbl] {$\vec{N}$};
% friction up the slope (prevents slipping)
\draw[force] (C) -- ++(206.6:1.25) node[below left, lbl] {$\vec{f}$};
\node[font=\scriptsize, engred] at (1.0,2.3) {free-body diagram};
\end{scope}
% == equals sign ==
\node at (5.4,0.9) {\Large $=$};
% == Right panel: kinetic diagram ==
\begin{scope}[shift={(6.0,0)}]
\draw[surf] (-0.3,0) -- (4.3,2.15);
\coordinate (C2) at (2.2,1.55);
\draw[draw=engblack, thick, dashed] (C2) circle (0.55);
\node[lbl] at (C2) {$G$};
% m a_G down the slope
\draw[kin] (C2) -- ++(206.6:1.55) node[below left, lbl] {$m\vec{a}_{G}$};
% I alpha (curved arrow, clockwise rolling)
\draw[kin, -{Stealth}] (C2) ++(40:0.85) arc (40:-150:0.85);
\node[engblue, font=\small] at ($(C2)+(-0.05,1.1)$) {$I_{G}\alpha$};
\node[font=\scriptsize, engblue] at (1.0,2.3) {kinetic diagram};
\end{scope}
\end{tikzpicture}
\caption[Free-body and kinetic diagram for a rolling cylinder]{The
paired free-body diagram (left, forces in red) and kinetic diagram
(right, the inertia terms $m\vec{a}_{G}$ and $I_{G}\alpha$ in blue)
for a solid cylinder rolling without slipping down an incline of
angle $\beta$. Setting the two diagrams equal component by component
is Newton's second law for the rigid body: $\sum\vec{F}=m\vec{a}_{G}$
along and normal to the slope, and $\sum M_{G}=I_{G}\alpha$ about the
mass centre. The friction force $\vec{f}$ provides the moment that
spins the cylinder up; the no-slip constraint $a_{G}=\alpha R$ closes
the system. The kinetic-diagram convention keeps the inertia terms
visibly separate from the forces and prevents the common error of
double-counting $m\vec{a}$ as if it were an applied force.}
\label{fig:vol03ch04:incline-fbd}
\end{figure}


\paragraph{Worked example: cylinder on an incline.} A solid cylinder
of mass $m$ and radius $R$ rolls without slipping down a plane
inclined at angle $\beta$. Take axes along and normal to the slope.
The forces are gravity $mg$, the normal force $N$, and friction $f$ up
the slope. Newton's second law along the slope and the moment equation
about the mass centre are
\[
mg\sin\beta - f = m a_{G}, \qquad fR = I_{G}\alpha,
\]
with $I_{G}=\tfrac{1}{2}mR^{2}$ for the solid cylinder. The no-slip
constraint $a_{G}=\alpha R$ closes the system. Substituting $\alpha =
a_{G}/R$ into the moment equation gives $f=\tfrac{1}{2}m a_{G}$, and
substituting that into the force equation gives $mg\sin\beta -
\tfrac{1}{2}m a_{G}=m a_{G}$, so
\[
a_{G}=\frac{g\sin\beta}{1+I_{G}/(mR^{2})}=\frac{2}{3}\,g\sin\beta.
\]
The factor $2/3$ is the cylinder's signature: a sphere
($I=\tfrac{2}{5}mR^{2}$) rolls faster at $\tfrac{5}{7}g\sin\beta$, a
ring ($I=mR^{2}$) slower at $\tfrac{1}{2}g\sin\beta$, and a
frictionless block (no rotation) fastest at $g\sin\beta$. The required
friction is $f=\tfrac{1}{3}mg\sin\beta$; for rolling to persist
without slipping it must not exceed $\mu_{s}N=\mu_{s}mg\cos\beta$,
which sets the steepest no-slip angle $\tan\beta\le 3\mu_{s}$.

\begin{archetype}[balance]
The rigid-body equations $\sum\vec{F}=m\vec{a}_{G}$ and $\sum M_{G}=
I_{G}\alpha$ are the dynamic generalisation of the static balance of
Chapters 1 through 3. Statics is the special case $\vec{a}_{G}=\vec{0}$,
$\alpha=0$, where both right-hand sides vanish and the sums of forces
and moments are separately zero. Dynamics replaces those zeros with
mass times acceleration and inertia times angular acceleration. Every
free-body diagram drawn for a static problem carries over unchanged;
only the right-hand side of the balance changes. The continuity is the
reason the same diagram discipline serves both: an engineer who can
balance a truss can write the equation of motion of a mechanism by
moving the inertia terms to the other side of the equals sign.
\end{archetype}

\begin{mastery}
\textbf{Non-inertial frames.} Newton's laws are stated in inertial
frames. In a frame accelerating with respect to inertial space (an
accelerating elevator, a turning vehicle, the rotating Earth at high
precision), Newton's second law gains pseudo-force terms. For a frame
translating with acceleration $\vec{A}_F$ relative to inertial space,
the equation of motion in the frame is
\[
\vec{F}_{\text{net}} + (-m\vec{A}_F) = m\vec{a}',
\]
where $\vec{a}'$ is the particle's acceleration in the non-inertial
frame and $-m\vec{A}_F$ is the inertial pseudo-force. For a rotating
frame with angular velocity $\vec{\Omega}$, the pseudo-forces are the
centrifugal force $-m\vec{\Omega} \times (\vec{\Omega} \times \vec{r}')$
(directed away from the rotation axis) and the Coriolis force
$-2m\vec{\Omega} \times \vec{v}'$ (perpendicular to the particle's
velocity in the rotating frame). Both pseudo-forces have engineering
consequences: the centrifugal force is what designers of centrifuges,
turbine blades, and grinding wheels must contend with; the Coriolis
force is what gives the world's atmosphere its hurricanes and what
biases long-range artillery in a predictable direction depending on
hemisphere.

The pseudo-forces are real in the accelerating frame and are felt by
every body in it: the passenger pressed into a turning car's door
feels the centrifugal pseudo-force, and the eastward deflection of a
body dropped from a tower is the Coriolis pseudo-force made visible.
They vanish only in inertial frames. Working in a rotating frame
trades the convenience of co-rotating coordinates for the cost of
carrying the extra terms; the choice of frame is the engineer's, and
Landau and Lifshitz set out the trade in full
\cite[\S{}39]{text:landau-mechanics}.
\end{mastery}

\begin{mastery}
\textbf{D'Alembert's principle: dynamics as statics.} The archetype box
treats dynamics as statics with the inertia terms moved to the other
side of the equals sign. D'Alembert's principle makes that move formal
and gives it a name. Rewrite Newton's second law as
\[
\vec{F}_{\text{net}}-m\vec{a}=\vec{0},
\]
and read the term $-m\vec{a}$ as a force, the \engterm{inertial force}
(or d'Alembert force). The equation then says the applied forces and the
inertial force are in equilibrium, so every method built for statics, the
free-body diagram, the summation of moments about any convenient point,
the principle of virtual work, carries over to dynamics unchanged once
the inertial force is added to the diagram. A pendulum bob swinging
through its arc is, in this reading, in equilibrium under gravity, the
string tension, and an inertial force $-m\vec{a}$; an engineer who can
solve a static three-force problem can solve the bob's instantaneous
balance with the same triangle.

The principle earns its keep in two places. The first is constrained
multi-body systems, where summing the virtual work of the applied and
inertial forces over a virtual displacement that respects the constraints
eliminates the constraint forces automatically, the same advantage the
Lagrangian enjoys and from the same root, since the Euler--Lagrange
equations follow from d'Alembert's principle plus the virtual-work
argument. The second is rotating machinery, where adding the inertial
force $-m\vec{a}_{G}$ at the centre of mass and the inertial couple
$-I_{G}\alpha$ to a rotor's free-body diagram turns a dynamics problem
into a statics problem that the bearing-reaction analysis of Chapter~2
already knows how to solve. The caution is the same one the non-inertial
box raised: the inertial force is a bookkeeping term, not a force any
third body exerts, so it has no third-law reaction. Treating it as real
inside the accelerating frame is correct; looking for the body that
pushes back is the error.
\end{mastery}

\section{Work-energy theorem}\pathtag{core}

The work done by a force $\vec{F}$ on a particle as it moves through
a displacement $d\vec{r}$ is
\[
dW = \vec{F} \cdot d\vec{r}.
\]
Over a finite trajectory from point $A$ to point $B$,
\[
W_{A \to B} = \int_A^B \vec{F} \cdot d\vec{r}.
\]
For a constant force, the integral is $\vec{F} \cdot (\vec{r}_B -
\vec{r}_A)$. For a force that varies along the path, the integral
depends in general on the path; for the special class of
\engterm{conservative} forces, it depends only on the endpoints.

Gravity, elastic spring forces, and electrostatic forces are
conservative. Friction, drag, and external applied forces are not.
For a conservative force there exists a scalar function (the
potential energy) such that $\vec{F} = -\nabla U$ and the work done
by the force is $W_{A \to B} = U(A) - U(B)$.

The \engterm{work-energy theorem} states that the net work done on a
particle equals the change in its kinetic energy:
\[
W_{\text{net}} = \Delta T = \tfrac{1}{2} m v_B^2 - \tfrac{1}{2} m v_A^2,
\]
where $T = \tfrac{1}{2} m v^2$ is the kinetic energy. The theorem
follows directly from $\vec{F} = m\vec{a}$ and the identity $\vec{a}
\cdot d\vec{r} = \vec{a} \cdot \vec{v} \,dt = d(\tfrac{1}{2} v^2)$.

The derivation is worth carrying in full, because it shows the theorem
to be Newton's second law integrated along the path rather than a
separate principle. Start from $\vec{F}_{\text{net}}=m\,d\vec{v}/dt$ and
take the dot product with the velocity $\vec{v}=d\vec{r}/dt$:
\[
\vec{F}_{\text{net}}\cdot\vec{v}=m\,\frac{d\vec{v}}{dt}\cdot\vec{v}.
\]
The right side is a total time derivative, because $\frac{d}{dt}\bigl(
\tfrac{1}{2}\vec{v}\cdot\vec{v}\bigr)=\vec{v}\cdot\frac{d\vec{v}}{dt}$, so
$\vec{F}_{\text{net}}\cdot\vec{v}=\frac{d}{dt}\bigl(\tfrac{1}{2}mv^{2}
\bigr)=\dot{T}$. The left side is the rate at which the net force does
work, the \engterm{power} delivered to the particle. The work-energy
theorem in its instantaneous form is therefore $P_{\text{net}}=\dot{T}$:
net power equals the rate of change of kinetic energy. Integrating over
time, with $\vec{v}\,dt=d\vec{r}$,
\[
\int_{t_{A}}^{t_{B}}\vec{F}_{\text{net}}\cdot\vec{v}\,dt
=\int_{A}^{B}\vec{F}_{\text{net}}\cdot d\vec{r}=W_{\text{net}},
\qquad
\int_{t_{A}}^{t_{B}}\dot{T}\,dt=T_{B}-T_{A}=\Delta T,
\]
which is the finite form $W_{\text{net}}=\Delta T$. The change of
variable from $dt$ to $d\vec{r}$ is the step that converts a statement
about time into a statement about path, and it is the reason the theorem
relates speeds to positions without time appearing in the answer.

\paragraph{Why some forces have a potential.} A force is conservative
when the work it does around any closed path is zero, equivalently when
the work between two points is the same for every path joining them. For
such a force the line integral $\int_{A}^{B}\vec{F}\cdot d\vec{r}$
defines a scalar function of the endpoint alone, up to the choice of a
reference point, and that function (with a sign convention) is the
potential energy
\[
U(\vec{r})=-\int_{\vec{r}_{0}}^{\vec{r}}\vec{F}\cdot d\vec{r}'.
\]
Differentiating recovers $\vec{F}=-\nabla U$, the gradient relation
stated above. For a uniform gravitational field $\vec{F}=-mg\hat{\jmath}$
the integral gives $U=mgy$; for a linear spring $\vec{F}=-kx\,\hat{\imath}$
it gives $U=\tfrac{1}{2}kx^{2}$. The test for conservativeness in
practice is whether the force field has zero curl, $\nabla\times\vec{F}=
\vec{0}$, which holds for gravity and spring forces and fails for
friction and drag, whose direction depends on the direction of motion
rather than on position alone. That failure is why a block dragged
around a closed loop on a rough table returns to its start having lost
energy to friction, while the same block on a frictionless track returns
with its kinetic energy intact.

For a rigid body in plane motion, the kinetic energy has translational
and rotational contributions:
\[
T = \tfrac{1}{2} m v_G^2 + \tfrac{1}{2} I_G \omega^2,
\]
with $v_G$ the centre-of-mass speed and $\omega$ the angular speed.
The work-energy theorem applies in the same form: $W_{\text{net}} =
\Delta T$, with work done by both forces and moments counted.

When all forces are conservative, the work-energy theorem becomes
conservation of mechanical energy: $T + U = \text{constant}$. The
conservation form is the working tool for trajectory problems: a
pendulum's swing, a roller-coaster car's speed, a projectile's range,
all admit closed-form solutions in terms of initial and final
kinetic and potential energies without solving the equation of
motion at every intermediate point.

\paragraph{Power, efficiency, and the rating of machines.} The
instantaneous form $P=\vec{F}\cdot\vec{v}$ is what an engineer sizes a
machine by, because a motor, an engine, or a muscle is rated by the rate
at which it does work rather than by the work itself. A hoist lifting a
load $W$ at steady speed $v$ demands power $Wv$; a vehicle holding speed
$v$ against a road load $R(v)$ demands $R(v)\,v$, which the traction case
study showed grows as $v^{3}$ at high speed because the drag part of
$R$ grows as $v^{2}$. The rotational counterpart is $P=\tau\omega$,
torque times angular speed, the form that prices a gearbox: an ideal
mesh preserves $\tau\omega$, so a reduction that divides speed multiplies
torque, and the power passes through unchanged. Real machines are not
ideal, and the gap is the \engterm{efficiency} $\eta=P_{\text{out}}/
P_{\text{in}}$, the fraction of the input power that leaves as useful
work rather than as the friction heat, the windage, and the electrical
loss that the work-energy bookkeeping assigns to the non-conservative
forces. A chain of stages multiplies its efficiencies, $\eta_{\text{
total}}=\eta_{1}\eta_{2}\cdots$, so a drivetrain of several modest stages
can waste a surprising fraction of its input before the load ever moves,
the reason a designer counts the stages between the prime mover and the
load and fights to remove them.

\begin{mastery}
\textbf{From energy to the Lagrangian.} The work-energy theorem hints
at a reformulation of mechanics that takes energy, not force, as the
primitive. Define the \engterm{Lagrangian} $L=T-U$, the kinetic energy
minus the potential energy, expressed in any convenient set of
generalised coordinates $q_{i}$. The equations of motion are then the
Euler--Lagrange equations
\[
\frac{d}{dt}\!\left(\frac{\partial L}{\partial \dot{q}_{i}}\right)
-\frac{\partial L}{\partial q_{i}}=0,
\]
one for each coordinate. For the mass on a spring, $L=\tfrac{1}{2}m
\dot{x}^{2}-\tfrac{1}{2}kx^{2}$, and the single Euler--Lagrange
equation returns $m\ddot{x}+kx=0$, the same result Newton's law gives.
The advantage shows when constraints are present: the Lagrangian
method works in coordinates adapted to the constraint (the angle of a
pendulum, the arc length along a track) and never requires the
constraint forces to be written down, because they do no work and so
drop out of $T-U$. The method scales to systems of many bodies where a
force-by-force accounting becomes unmanageable, and it is the natural
setting for the conserved quantities of Chapter 6 and for the control
formulations of Volume IX. Landau and Lifshitz build all of mechanics
from the action principle $\delta\!\int L\,dt=0$ that underlies the
Euler--Lagrange equations \cite[\S{}1--2]{text:landau-mechanics}.
\end{mastery}

\section{Impulse-momentum}\pathtag{core}

Integrating Newton's second law $\vec{F}_{\text{net}} = d\vec{p}/dt$
over a time interval gives the \engterm{impulse-momentum theorem}:
\[
\int_{t_1}^{t_2} \vec{F}_{\text{net}} \,dt = \Delta \vec{p} = m\vec{v}_2
- m\vec{v}_1.
\]
The left-hand side is the \engterm{impulse}; the right-hand side is
the change in linear momentum. The theorem is exactly equivalent to
Newton's second law for a constant-mass particle; for a variable-mass
particle, the equivalent statement is in terms of momentum rather
than $m\vec{a}$, as noted in section 4.3.

The impulse-momentum form is the right tool for problems in which the
force is large but acts over a short time (impacts, collisions,
explosions, impulse-driven motors), for which the detailed
time-history of the force is unknown or irrelevant but the integrated
effect is the answer required. Two important regimes:

\engterm{Impulsive forces}. A baseball struck by a bat, a hammer blow,
a vehicle collision, a piledriver impact. The contact time is short
(milliseconds to a few hundred milliseconds), the peak force is
large (hundreds of times the body's weight), and the impulse is what
the engineer can measure or compute from the change in velocity. The
detailed force history is rarely needed.

\engterm{Variable-mass systems}. A rocket loses mass at rate $\dot{m}$
through its exhaust at velocity $\vec{v}_{\text{exh}}$ relative to the
rocket. The rocket's equation of motion is the Tsiolkovsky rocket
equation,
\[
m\frac{d\vec{v}}{dt} = \vec{F}_{\text{ext}} + \dot{m}_{\text{loss}}
\vec{v}_{\text{exh,rel}},
\]
where $\dot{m}_{\text{loss}}$ is the mass-loss rate (positive) and
$\vec{v}_{\text{exh,rel}}$ is the exhaust velocity relative to the
rocket. The second term is the \engterm{thrust}: a force generated by
expelling mass. The same principle drives jet engines, water rockets,
and squids.

\engterm{Steady momentum flux}. A continuous stream carries momentum at
a steady rate, and the force to deflect or stop it follows from the same
theorem written per unit time rather than per collision. A stream of
mass flow $\dot{m}=\rho A v$ (density $\rho$, cross-section $A$, speed
$v$) carries momentum past a plane at the rate $\dot{m}v$. Turning that
stream through an angle, or arresting it, changes its momentum flux, and
by $\vec{F}=d\vec{p}/dt$ the change equals the force. A jet brought to
rest against a flat plate delivers a force $\dot{m}v=\rho A v^{2}$; a jet
turned back on itself by a cupped vane delivers $2\dot{m}v$, twice as
much, because the stream leaves with its momentum reversed rather than
merely removed. The $v^{2}$ dependence, hidden inside $\dot{m}v=\rho A
v^{2}$, is the reason a fire hose kicks hard, a pressure washer bites,
and a river in flood shoves a bridge pier with a force that grows as the
square of the current. The same accounting sizes the reaction on a
pipe bend, where the fluid's change of direction loads the anchor
flanges, and it is the working principle of the impulse water turbine
treated in the vane vignette below.

The angular analogue of impulse-momentum is the angular impulse-angular
momentum theorem:
\[
\int_{t_1}^{t_2} \vec{M}_O \,dt = \Delta \vec{L}_O,
\]
where $\vec{L}_O = \int \vec{r} \times \vec{v} \,dm$ is the angular
momentum about point $O$. For a rigid body rotating about a
centre-of-mass axis, $\vec{L}_G = I_G \vec{\omega}$.

\paragraph{Impulsive forces separate from finite ones.} Over the short
time of an impact the impulse-momentum theorem admits a simplification
that makes hard problems tractable. Split the forces on a body into two
kinds: impulsive forces, whose magnitude is enormous during the contact
(the bat on the ball, the hammer on the anvil), and finite forces such as
gravity and ordinary spring loads, whose magnitude is bounded. The
impulse of a finite force over a contact time $\Delta t$ is at most
$F_{\max}\Delta t$, which vanishes as $\Delta t\to 0$, while the impulse
of an impulsive force stays finite because its magnitude grows as the
time shrinks. The working rule follows: over a collision, keep only the
impulsive impulses and drop the finite ones. A ball struck by a bat
changes its momentum by the bat's impulse alone; gravity, acting for the
millisecond of contact, contributes an impulse of order $mg\,\Delta t$
that is negligible against it. The same rule is why the exploding-shell
example could ignore gravity across the burst and why the accident
reconstruction treats the crush as conserving momentum despite the road
friction acting throughout: friction is finite, the crush forces are
impulsive, and only the impulsive forces survive the short-time limit.

\paragraph{Impulsive tension in connected bodies.} When a taut,
inextensible cord suddenly jerks a system into motion, the tension during
the jerk is impulsive and the constraint it enforces is a relation among
the bodies' velocity changes, not their accelerations. Consider two
masses $m_{1}$ and $m_{2}$ joined by a slack cord that snaps taut when
$m_{1}$ has fallen to speed $v_{0}$ while $m_{2}$ is still at rest. During
the vanishingly short snatch the cord carries an impulsive tension of
magnitude $\hat{T}=\int T\,dt$, equal and opposite on the two bodies. The
impulse-momentum theorem for each body, dropping the finite gravity
impulse over the snatch, gives $-\hat{T}=m_{1}(v-v_{0})$ for the falling
mass and $+\hat{T}=m_{2}v$ for the mass jerked into motion, where the
inextensible cord forces a common post-snatch speed $v$. Adding the two
eliminates the impulsive tension and gives $v=m_{1}v_{0}/(m_{1}+m_{2})$,
the same reduced form the inelastic collision produced, with the cord
snatch playing the role of the sticking contact. The impulsive tension is
then $\hat{T}=m_{2}v=m_{1}m_{2}v_{0}/(m_{1}+m_{2})$, and the energy the
snatch throws away is the same $\tfrac{1}{2}\mu v_{0}^{2}$ with reduced
mass $\mu=m_{1}m_{2}/(m_{1}+m_{2})$ that a perfectly inelastic collision
loses. A cord that snaps a system taut is an inelastic collision
conducted through a rope, which is why a tow line jerked from slack loses
energy and loads its anchors impulsively, the quantitative backing for the
snatch worked example above.

\begin{estimation}
\textbf{Estimate the time and distance for a passenger car to brake
from $100\,\si{\kilo\meter\per\hour}$ to rest on dry asphalt.}

\smallskip
\textit{Estimate, before reading on.}

\smallskip
At rest, the only horizontal force on a braking car is the friction
between tyres and road. The friction is at most $\mu_s m g$ when the
tyres do not skid (anti-lock braking limits the tyres at the
static-friction edge; modern passenger cars stop better with anti-lock
than without). For tyres on dry asphalt, $\mu_s \approx 0.8$.

The deceleration is therefore at most $\mu_s g \approx 0.8 \times
9.81 \approx 7.8\,\si{\meter\per\second\squared}$. Initial speed is
$v_0 = 100\,\si{\kilo\meter\per\hour} = 27.8\,\si{\meter\per\second}$.

Time to stop: $t = v_0 / a = 27.8 / 7.8 \approx 3.6\,\si{\second}$.
Distance: $d = v_0^2 / (2a) = (27.8)^2 / (2 \times 7.8) \approx
49\,\si{\meter}$.

Estimate: about $3.5\,\si{\second}$ and $50\,\si{\meter}$.

The postmortem: the dominant uncertainty is the coefficient of
friction, which drops to $\mu_s \approx 0.3$ to $0.4$ on wet asphalt
and to $\mu_s \approx 0.05$ to $0.15$ on snow or ice. On wet roads
the stopping distance roughly doubles ($\sim 100\,\si{\meter}$); on
ice it grows by a factor of five or more ($\sim 300\,\si{\meter}$).
The factor of $0.8$ for dry road is a working approximation;
specific tyre compounds, road textures, and tread conditions move
the value over a $\pm 20\,\%$ range. Vehicle stopping tables
published by transport agencies are time-stamped to the road
condition and are the right source for design and forensic work.

The estimate explains why the highway code requires longer following
distances in rain: at $100\,\si{\kilo\meter\per\hour}$ a driver who
reacts in $1\,\si{\second}$ has already travelled $28\,\si{\meter}$
before the brake engages; the total stopping distance on dry road is
$\sim 80\,\si{\meter}$ from the moment a hazard appears, and on wet
road $\sim 130\,\si{\meter}$.
\end{estimation}

\section{Conservation of momentum}\pathtag{core}

A system of particles with no external forces conserves total linear
momentum:
\[
\vec{P}_{\text{system}} = \sum_i m_i \vec{v}_i = \text{constant}.
\]
The statement is a consequence of Newton's third law: the internal
forces between any pair of particles in the system come in equal-
and-opposite pairs, so they sum to zero in the system's total force
balance. With no external forces, the system's centre of mass moves
at constant velocity, and the total momentum (which is the
centre-of-mass momentum) is constant.

The conservation law applies to collisions: even when the contact
forces during the collision are large and unknown in detail, the
total momentum before equals the total momentum after, provided no
external impulse acts during the contact time. For a one-dimensional
two-particle collision:
\[
m_1 v_{1i} + m_2 v_{2i} = m_1 v_{1f} + m_2 v_{2f}.
\]
This gives one equation for two unknowns; a second equation is
required, usually a kinematic constraint (the particles stick
together, the collision is elastic, the coefficient of restitution is
specified).

% Figure: before/after schematic of a one-dimensional perfectly inelastic
% collision (the two-car problem), with momentum bars.
\begin{figure}[htbp]
\centering
\begin{tikzpicture}[
  car/.style={draw=engblack, thick, rounded corners=2pt},
  vel/.style={draw=engblue, -{Stealth}, very thick},
  lbl/.style={font=\small},
]
% == BEFORE ==
\node[lbl, anchor=west] at (-0.3,2.3) {\textbf{before}};
% car 1: moving right at 20 m/s, m=1500
\draw[car, fill=engblue!12] (0,1.4) rectangle (1.4,2.0);
\node[font=\scriptsize] at (0.7,1.7) {$1500$};
\draw[vel] (1.45,1.7) -- ++(1.0,0) node[above, font=\scriptsize] {$20\,\si{\meter\per\second}$};
% car 2: stationary, m=1000
\draw[car, fill=engamber!15] (3.4,1.4) rectangle (4.6,2.0);
\node[font=\scriptsize] at (4.0,1.7) {$1000$};
\node[font=\scriptsize, below] at (4.0,1.35) {$0\,\si{\meter\per\second}$};
% == AFTER ==
\node[lbl, anchor=west] at (-0.3,0.3) {\textbf{after}};
% combined wreck moving right at 12 m/s
\draw[car, fill=engblue!12] (1.6,-0.6) rectangle (3.0,0.0);
\draw[car, fill=engamber!15] (3.0,-0.6) rectangle (4.2,0.0);
\node[font=\scriptsize] at (2.9,-0.3) {$2500$};
\draw[vel] (4.25,-0.3) -- ++(0.85,0) node[above, font=\scriptsize] {$12\,\si{\meter\per\second}$};
% momentum bar (conserved): before = 1500*20 = 30000; after = 2500*12 = 30000
\draw[draw=engred, very thick, -{Stealth}] (6.0,2.0) -- ++(2.5,0)
  node[midway, above, font=\scriptsize, engred] {$p=30\,000\,\si{\kilogram\meter\per\second}$};
\draw[draw=engred, very thick, -{Stealth}] (6.0,0.0) -- ++(2.5,0)
  node[midway, above, font=\scriptsize, engred] {$p=30\,000\,\si{\kilogram\meter\per\second}$};
\node[engred, font=\scriptsize] at (7.25,1.0) {momentum conserved};
\end{tikzpicture}
\caption[Perfectly inelastic two-car collision]{The momentum
bookkeeping of the perfectly inelastic two-car collision (sixth
calculation exercise). A $1{,}500\,\si{\kilogram}$ car at
$20\,\si{\meter\per\second}$ strikes a stationary
$1{,}000\,\si{\kilogram}$ car; they lock together. Total momentum is
conserved across the contact because the impulsive contact forces are
internal to the two-car system, so the common final speed is
$v_{f}=(1500)(20)/2500=12\,\si{\meter\per\second}$ (red bars, equal
before and after). Kinetic energy is not conserved: it falls from
$300\,\si{\kilo\joule}$ to $180\,\si{\kilo\joule}$, a $40\,\%$ loss
into crush, heat, and sound. The split is the defining contrast
between the two conservation tools: momentum survives the collision,
mechanical energy need not.}
\label{fig:vol03ch04:collision}
\end{figure}


The perfectly inelastic case is the cleanest illustration of the
contrast between the two conservation tools.
Figure~\ref{fig:vol03ch04:collision} works the two-car collision of
the calculation set: momentum survives the contact intact, because the
crushing forces are internal to the two-car system, while kinetic
energy falls by $40\,\%$ into crush, heat, and sound. An engineer
analysing a crash uses momentum to find the post-impact motion and
the energy balance to find how much was absorbed; the two answers are
different numbers from the same event, and confusing them is a
recurring forensic error.

The \engterm{coefficient of restitution} $e$ is the ratio of the
relative speed after impact to the relative speed before:
\[
e = -\frac{v_{1f} - v_{2f}}{v_{1i} - v_{2i}}.
\]
For a perfectly elastic collision, $e = 1$ and kinetic energy is
conserved. For a perfectly inelastic collision, $e = 0$ and the
particles stick together (kinetic energy is not conserved; the lost
energy goes into deformation, heat, and sound). Real collisions have
$0 < e < 1$ depending on the materials and the geometry; a steel
ball-bearing on a steel anvil has $e \approx 0.8$; a tennis ball on
a tennis racket strung at $25\,\text{kp}$ has $e \approx 0.7$; a
beanbag on a concrete floor has $e \approx 0.05$.

\paragraph{Solving a collision in closed form.} Momentum and restitution
together close the one-dimensional two-body collision exactly, and the
algebra is worth carrying once so that the special cases fall out as
limits. Conservation of momentum and the restitution definition are the
two equations
\[
m_{1}v_{1i}+m_{2}v_{2i}=m_{1}v_{1f}+m_{2}v_{2f}, \qquad
v_{2f}-v_{1f}=e\,(v_{1i}-v_{2i}).
\]
Treat them as a linear system in the two unknowns $v_{1f}$ and $v_{2f}$.
Solving gives
\[
v_{1f}=\frac{(m_{1}-e\,m_{2})v_{1i}+(1+e)m_{2}v_{2i}}{m_{1}+m_{2}},
\qquad
v_{2f}=\frac{(m_{2}-e\,m_{1})v_{2i}+(1+e)m_{1}v_{1i}}{m_{1}+m_{2}}.
\]
Three checks read the result. With $e=1$ (elastic) and a stationary
target $v_{2i}=0$, the expressions reduce to $v_{1f}=(m_{1}-m_{2})v_{1i}/
(m_{1}+m_{2})$ and $v_{2f}=2m_{1}v_{1i}/(m_{1}+m_{2})$, the familiar
elastic-collision pair: equal masses exchange velocities ($v_{1f}=0$,
$v_{2f}=v_{1i}$), a heavy striker barely slows while the light target
flies off near $2v_{1i}$, and a light striker rebounds off a heavy target
with reversed velocity. With $e=0$ both final velocities collapse to the
common $v_{f}=(m_{1}v_{1i}+m_{2}v_{2i})/(m_{1}+m_{2})$, the stuck-together
case. The kinetic energy lost in a collision of restitution $e$, written
in the centre-of-mass frame, is
\[
\Delta T=-\tfrac{1}{2}\,\frac{m_{1}m_{2}}{m_{1}+m_{2}}\,(1-e^{2})\,(v_{1i}
-v_{2i})^{2},
\]
which is negative for every $e<1$, vanishes only at $e=1$, and reaches
its largest magnitude at $e=0$. The combination $\mu=m_{1}m_{2}/(m_{1}+
m_{2})$ is the \engterm{reduced mass}: the collision behaves, in the
relative coordinate, like a single particle of mass $\mu$ losing the
fraction $(1-e^{2})$ of its relative kinetic energy.

\begin{mastery}
\textbf{What the coefficient of restitution is measuring.} The number
$e$ is treated as a material-and-geometry constant, and for a given pair
of bodies over a modest speed range it nearly is, but it packs a good
deal of physics into a single figure. The clean way to see what it holds
is to split a collision into two phases. During \emph{compression} the
two bodies deform and the contact force does negative work on their
relative motion until the approach speed reaches zero at maximum
deformation; during \emph{restitution} the stored elastic energy pushes
them apart and the force does positive work as they separate. If the
contact were perfectly elastic every joule stored in compression would
be returned in restitution and the separation speed would equal the
approach speed, giving $e=1$. Real contacts return less: some energy
goes into plastic deformation of the material, some into internal
vibration and sound, some into heat, and $e$ is the square root of the
fraction of the relative kinetic energy that survives, since energy
scales as speed squared. A hard steel-on-steel contact stays nearly
elastic ($e\approx 0.8$) because the stresses stay below yield and the
material springs back; a soft or ductile contact ($e\approx 0.05$ for a
beanbag) yields plastically and returns almost nothing. Two further
facts matter in practice. First, $e$ falls as the impact speed rises,
because a faster impact drives more of the material past yield, so a
restitution measured at low speed overpredicts the bounce of a hard
hit. Second, $e$ is a property of the contact, not of either body alone,
which is why the same ball reports different values off a wooden floor,
a concrete slab, and a taut string. An engineer who quotes a single $e$
should record the speed and the surface pair it was measured on, the
same discipline the friction coefficient demanded.
\end{mastery}

\paragraph{The centre of mass carries the system momentum.} The reason
internal forces drop out is the centre-of-mass theorem. Define the
centre of mass of a system by $\vec{R}=\bigl(\sum_{i}m_{i}\vec{r}_{i}
\bigr)/M$ with $M=\sum_{i}m_{i}$ the total mass. Differentiating twice,
$M\ddot{\vec{R}}=\sum_{i}m_{i}\ddot{\vec{r}}_{i}=\sum_{i}\vec{F}_{i}$,
the sum of all forces on all particles. Split each force into external
and internal parts; the internal forces cancel in equal-and-opposite
pairs by Newton's third law, leaving
\[
M\ddot{\vec{R}}=\sum_{i}\vec{F}_{i}^{\text{ext}}=\vec{F}_{\text{ext}}.
\]
The centre of mass moves as though the whole mass were concentrated
there and all external forces applied there, whatever the internal
forces do. With no external force the centre of mass keeps a constant
velocity, which is the conservation of total momentum $\vec{P}=M
\dot{\vec{R}}$ restated. A shell that bursts in flight scatters its
fragments along paths whose centre of mass continues on the original
parabola until a fragment strikes the ground, because the explosion is
internal and changes no external force.

The angular analogue is conservation of angular momentum: a system
with no external moment about a point conserves the total angular
momentum about that point. Conservation of angular momentum is the
mechanism behind a figure skater's spin (pulling in arms reduces
$I$, so $\omega$ must increase to keep $L = I\omega$ constant), a
satellite's attitude control (reaction wheels exchange angular
momentum between the spacecraft and the wheels), and the gyroscopic
stability of a bicycle wheel. The skater's spin carries a second
lesson worth the arithmetic: $L$ is conserved but the rotational
kinetic energy $T=\tfrac{1}{2}I\omega^{2}=L^{2}/(2I)$ rises as $I$
falls, so pulling the arms in does mechanical work. The skater's
muscles supply that work against the centripetal force needed to draw
the arms inward against their rotation, which is why a fast spin feels
effortful to tighten and why the energy and momentum accounts give
different answers from the same manoeuvre, the same split that
separates the momentum and energy analyses of a collision. A reaction
wheel makes the exchange explicit and reversible: spinning the wheel
one way turns the spacecraft the other, with the total angular momentum
fixed, and the wheel can be spun back down to undo the turn, which is
how a spacecraft points without expending propellant.

\begin{mastery}
\textbf{A decision map: which conservation tool to reach for.} Three
integrals of Newton's second law solve most particle and rigid-body
problems, and choosing the right one at the start saves the most work.
The choice turns on what the problem gives and what it asks.

Reach for the \emph{work-energy theorem} $W_{\text{net}}=\Delta T$ when
the question relates speeds to positions and the forces are known as
functions of position. A block sliding down a ramp of known height, a
pendulum swinging through a known angle, a car braking over a known
distance: each pairs a position change with a speed change, and energy
connects them without time entering at all. Energy is the tool when
time is neither given nor wanted.

Reach for the \emph{impulse-momentum theorem} $\int\vec{F}\,dt=\Delta
\vec{p}$ when the question relates velocities to time and the force
acts over a known interval, especially a short and violent one. A bat
striking a ball, a rocket burning for a known duration, a jet of water
hitting a vane: the force history may be unknown in detail, but its
time integral is the change in momentum. Momentum is the tool when time
is in the problem and the detailed force is not.

Reach for \emph{conservation of momentum} when two or more bodies
interact through internal forces and no external impulse acts during
the interaction. A collision, an explosion, a recoil, a stage
separation: the internal forces cancel in pairs, so the system's total
momentum is the same before and after even though no single force is
known. The same logic in rotation gives conservation of angular
momentum for a system free of external moments.

Two cautions keep the map honest. Energy is conserved only when the
forces are conservative; an inelastic collision conserves momentum but
not kinetic energy, which is why the accident reconstruction uses
momentum to find the post-impact motion and the energy balance only to
report how much was lost. And momentum is a vector while energy is a
scalar, so a two-dimensional collision gives two momentum equations
(one per axis) but a single energy equation. The hardest problems use
two of the three tools in sequence, as the ballistic pendulum does:
momentum across the fast embedding, energy across the slow swing.
\end{mastery}

\section{Worked examples}\pathtag{standard}

\subsection{Projectile in a uniform gravitational field}

A projectile is launched from ground level with speed $v_0$ at angle
$\theta$ above horizontal. Air resistance is neglected. Find the
maximum height, the range, and the time of flight.

\textbf{Equations of motion.} With horizontal axis $x$ and vertical
axis $y$, the acceleration is $(0, -g)$ and the initial velocity is
$(v_0 \cos\theta, v_0 \sin\theta)$. Integrating,
\[
x(t) = v_0 \cos\theta \cdot t, \quad y(t) = v_0 \sin\theta \cdot t -
\tfrac{1}{2} g t^2.
\]

\textbf{Time of flight.} The projectile returns to $y = 0$ at $t = 2
v_0 \sin\theta / g$.

\textbf{Range.} Substituting into $x(t)$:
\[
R = \frac{v_0^2 \sin 2\theta}{g}.
\]
The maximum range is at $\theta = 45^{\circ}$, where $R_{\max} = v_0^2
/ g$. The result is exact for the no-air-resistance model; real
projectiles peak well below $45^{\circ}$ because drag is larger at
higher trajectory angles where flight time is longer.

\textbf{Maximum height.} At the apex, $\dot{y} = 0$:
\[
h_{\max} = \frac{v_0^2 \sin^2\theta}{2g}.
\]

% Figure: projectile trajectories in a uniform field (no drag) at three
% launch angles, same launch speed, showing the 45-degree range maximum
% and the symmetry of complementary angles.
\begin{figure}[htbp]
\centering
\begin{tikzpicture}
\begin{axis}[
  width=12cm, height=7cm,
  xlabel={horizontal distance $x$ (\si{\meter})},
  ylabel={height $y$ (\si{\meter})},
  xmin=0, xmax=95, ymin=0, ymax=48,
  axis lines=left,
  samples=200,
  legend pos=north east,
  legend cell align=left,
  grid=major, grid style={enggray!20},
  tick label style={font=\scriptsize},
  label style={font=\small},
  legend style={font=\scriptsize},
]
% v0 = 30 m/s, g = 9.81. Trajectory y(x) = x tan(theta) - g x^2 / (2 v0^2 cos^2 theta).
% theta = 30: tan=0.5774, cos^2=0.75, g/(2 v0^2 cos^2) = 9.81/(2*900*0.75)=0.007267
\addplot[engblue, very thick, domain=0:79.4] {0.5774*x - 0.007267*x^2};
\addlegendentry{$\theta=30^{\circ}$}
% theta = 45: tan=1, cos^2=0.5, coeff = 9.81/(2*900*0.5)=0.0109
\addplot[engred, very thick, domain=0:91.7] {1.0*x - 0.0109*x^2};
\addlegendentry{$\theta=45^{\circ}$ (max range)}
% theta = 60: tan=1.7321, cos^2=0.25, coeff = 9.81/(2*900*0.25)=0.0218
\addplot[engamber, very thick, domain=0:79.4] {1.7321*x - 0.0218*x^2};
\addlegendentry{$\theta=60^{\circ}$}
\end{axis}
\end{tikzpicture}
\caption[Projectile trajectories at three launch angles]{Drag-free
projectile trajectories for a fixed launch speed $v_{0}=30\,\si{\meter
\per\second}$ at three launch angles. The range $R=v_{0}^{2}\sin
2\theta/g$ is maximised at $\theta=45^{\circ}$ (red), where it reaches
$R_{\max}=v_{0}^{2}/g\approx 91.7\,\si{\meter}$. The complementary
angles $30^{\circ}$ (blue) and $60^{\circ}$ (amber) share the same
range because $\sin 2\theta$ takes the same value at $\theta$ and
$90^{\circ}-\theta$; the steeper launch trades range for height and
flight time. Real projectiles peak well below $45^{\circ}$ once drag
is admitted, because the longer flight time of a high launch exposes
the body to drag for longer.}
\label{fig:vol03ch04:projectile-angles}
\end{figure}


Figure~\ref{fig:vol03ch04:projectile-angles} plots the trajectory for
three launch angles at fixed speed. Complementary angles share a
range, because $\sin 2\theta$ takes the same value at $\theta$ and at
$90^{\circ}-\theta$; the choice between them trades range against
flight time and apex height. The drag-free model is the right first
cut for a dense, slow, compact projectile (a shot put, a thrown
stone). For a light or fast one it is a poor predictor: the file
\texttt{projectile\_drag.py} integrates the same launch under linear
and quadratic drag and shows the range of a tennis ball
($v_{0}=30\,\si{\meter\per\second}$, $45^{\circ}$) falling from
$91.7\,\si{\meter}$ to $38.9\,\si{\meter}$ once quadratic drag is
admitted, with the trajectory losing its fore-aft symmetry
(Figure~\ref{fig:vol03ch04:projectile-drag}). The asymmetry is the
diagnostic the chapter project asks the reader to find in video data.

% Figure: comparison of drag-free, linear-drag, and quadratic-drag
% trajectories from the same launch condition, computed numerically.
% Curves are precomputed sample points (no closed form for drag cases).
\begin{figure}[htbp]
\centering
\begin{tikzpicture}
\begin{axis}[
  width=12cm, height=7cm,
  xlabel={horizontal distance $x$ (\si{\meter})},
  ylabel={height $y$ (\si{\meter})},
  xmin=0, xmax=95, ymin=0, ymax=26,
  axis lines=left,
  legend pos=north east,
  legend cell align=left,
  grid=major, grid style={enggray!20},
  tick label style={font=\scriptsize},
  label style={font=\small},
  legend style={font=\scriptsize},
]
% Drag-free: v0=30, theta=45, coeff 0.0109 (from previous figure).
\addplot[engblue, very thick, domain=0:91.7, samples=120] {1.0*x - 0.0109*x^2};
\addlegendentry{no drag}
% Linear drag, sampled from projectile_drag.py (K_LIN = 0.02 N.s/m).
\addplot[engamber, very thick, dashed] coordinates {
  (0,0) (4.1,3.9) (7.9,7.2) (11.5,9.8) (14.8,11.9) (17.9,13.5)
  (20.8,14.6) (23.5,15.2) (26.0,15.5) (28.3,15.3) (30.5,14.8)
  (32.5,13.9) (34.4,12.7) (36.2,11.2) (37.8,9.4) (39.4,7.4)
  (40.8,5.1) (42.1,2.6) (43.3,0)};
\addlegendentry{linear drag}
% Quadratic drag, sampled from projectile_drag.py (tennis ball).
\addplot[engred, very thick] coordinates {
  (0,0) (4.0,3.8) (7.5,6.8) (10.7,9.1) (13.7,10.9) (16.4,12.2)
  (19.0,13.0) (21.4,13.4) (23.7,13.4) (25.9,13.0) (28.0,12.2)
  (30.0,11.1) (31.8,9.7) (33.6,8.0) (35.3,6.0) (36.8,3.7)
  (38.3,1.2) (38.9,0)};
\addlegendentry{quadratic drag}
\end{axis}
\end{tikzpicture}
\caption[Effect of drag on a projectile trajectory]{Trajectory of a
projectile launched at $45^{\circ}$ with $v_{0}=30\,\si{\meter\per
\second}$, computed under three force models by the solver in
\texttt{projectile\_drag.py}. The drag-free parabola (blue) is
symmetric about its apex and reaches $91.7\,\si{\meter}$. Quadratic
drag $F=-b\,|\vec{v}|\vec{v}$ (red), the correct model for a tennis
ball in this speed range ($\mathrm{Re}\sim 10^{4}$ to $10^{5}$), cuts
the range to $38.9\,\si{\meter}$ and breaks the symmetry: the
descending branch is steeper than the ascending branch because the
body sheds horizontal speed throughout the flight. Linear drag (amber
dashed), valid only in the creeping-flow regime, lands at
$43.3\,\si{\meter}$; it is the wrong model for a ball of this size and
speed, included to show how far the regime assumption can move the
answer. The asymmetry is the visible signature of drag in the
project's video data.}
\label{fig:vol03ch04:projectile-drag}
\end{figure}


\subsection{An oscillating cart on a spring}

A cart of mass $m = 2\,\si{\kilogram}$ on a frictionless horizontal
track is connected to a wall by a spring of stiffness $k = 50\,\si{
\newton\per\meter}$. The cart is displaced by $x_0 = 0.1\,\si{\meter}$
from equilibrium and released from rest. Find the cart's motion.

\textbf{Equation of motion.} The spring force is $-kx$; Newton's
second law gives
\[
m\ddot{x} = -kx, \quad \ddot{x} + \omega_n^2 x = 0,
\]
with $\omega_n = \sqrt{k/m}$ the natural angular frequency. The
solution is
\[
x(t) = x_0 \cos(\omega_n t),
\]
since the cart starts from rest at $x_0$. With $\omega_n = \sqrt{50/2}
= 5\,\si{\radian\per\second}$, the period is $T = 2\pi/\omega_n
\approx 1.26\,\si{\second}$. The cart's maximum speed is $v_{\max} =
x_0 \omega_n = 0.5\,\si{\meter\per\second}$ at the equilibrium
position; the maximum acceleration is $a_{\max} = x_0 \omega_n^2 =
2.5\,\si{\meter\per\second\squared}$ at the extreme positions.

The energy of the oscillator: at any instant $E = \tfrac{1}{2} m
\dot{x}^2 + \tfrac{1}{2} k x^2$, constant in time. At the
equilibrium position the energy is all kinetic, $E = \tfrac{1}{2} m
v_{\max}^2 = 0.25\,\si{\joule}$; at the extremes it is all potential,
$E = \tfrac{1}{2} k x_0^2 = 0.25\,\si{\joule}$. Consistency.

Chapter 6 of this volume develops the mass-spring system into the
full theory of vibrations, including damping, forced response,
and resonance.

\subsection{A spinning top}

A symmetric top of mass $m$ and moment of inertia $I$ about its
spinning axis rotates with angular velocity $\omega$ about an axis
inclined at angle $\theta$ from vertical. The top is supported at a
single point on its axis. Find the rate at which the spinning axis
precesses about the vertical.

\textbf{Angular-momentum analysis.} The top's angular momentum about
the support point is approximately $\vec{L} = I\vec{\omega}$,
directed along the spinning axis (the approximation is excellent
when the spin rate is large compared to the precession rate). The
external moment about the support is $\vec{M} = \vec{r}_G \times m
\vec{g}$, with $\vec{r}_G$ from the support to the centre of mass.
The moment is horizontal (perpendicular to the vertical $\vec{g}$
and to the inclined spin axis), and the angular momentum changes
direction at rate
\[
\frac{d\vec{L}}{dt} = \vec{M}.
\]
Since $\vec{L}$ has approximately constant magnitude (the spin is
preserved), the moment must rotate $\vec{L}$ in the horizontal
plane, which corresponds to a precession of the spin axis about
the vertical. The precession rate is
\[
\Omega_p = \frac{|\vec{M}|}{|\vec{L}| \sin\theta} = \frac{m g r_G}{I
\omega},
\]
where $r_G$ is the distance from the support to the centre of mass.
A typical mechanical gyroscope toy with $I\omega^2 \gg m g r_G$
precesses slowly compared to its spin rate; the precession is
counterintuitive in direction (perpendicular to the applied moment)
and is what makes gyroscopes useful for attitude reference in
aircraft and spacecraft.

\subsection{A planet in a central force}

A body of mass $m$ moves under the gravitational pull of a much larger
mass $M$ fixed at the origin. The force is central (directed along the
line to the origin) and inverse-square,
\[
\vec{F}=-\frac{GMm}{r^{2}}\hat{r}.
\]
Because the force is central, it exerts no moment about the origin, so
the body's angular momentum $L=mr^{2}\dot{\theta}$ about the origin is
conserved. The conservation step is worth showing: the moment about the
origin is $\vec{M}=\vec{r}\times\vec{F}$, and a central force is parallel
to $\vec{r}$, so the cross product vanishes and $d\vec{L}/dt=\vec{M}=
\vec{0}$. The angular momentum of a particle in polar coordinates is
$\vec{L}=m\vec{r}\times\vec{v}=m\,r\hat{r}\times(\dot{r}\hat{r}+r
\dot{\theta}\hat{\theta})=mr^{2}\dot{\theta}\,\hat{k}$, since
$\hat{r}\times\hat{r}=\vec{0}$ and $\hat{r}\times\hat{\theta}=\hat{k}$,
which identifies the conserved scalar as $L=mr^{2}\dot{\theta}$.
Conservation of angular momentum is Kepler's second law in disguise. The
area swept by the radius vector in a time $dt$ is the area of the thin
triangle with sides $r$ and $r+dr$ subtending the angle $d\theta$, which
to first order is $dA=\tfrac{1}{2}r^{2}\,d\theta$; dividing by $dt$ gives
the \engterm{areal velocity}
\[
\frac{dA}{dt}=\tfrac{1}{2}r^{2}\dot{\theta}=\frac{L}{2m},
\]
constant in time. The radius vector sweeps equal areas in equal times,
so the body moves fast near the centre and slow far from it
(Figure~\ref{fig:vol03ch04:areal-velocity}). Because gravity is
conservative, mechanical energy is conserved too, and the combination of
the two conservation laws forces the bound orbit to be a closed ellipse
with the centre at one focus (Kepler's first law). The orbital period
and the semi-major axis obey Kepler's third law $T^{2}\propto a^{3}$;
the full derivation of the orbit shape from the inverse-square force is
carried in the Kepler mastery box of the Hohmann case study.

% Figure: Kepler's second law as conserved areal velocity. The radius
% vector sweeps equal areas in equal times, so the body moves fast near
% the focus and slow far from it. Two shaded sweeps of equal area subtend
% a wide arc near periapsis and a narrow arc near apoapsis.
\begin{figure}[htbp]
\centering
\begin{tikzpicture}[lbl/.style={font=\small}]
% ellipse: semi-major a=3.4, semi-minor b=2.4, focus offset c=sqrt(a^2-b^2)
% c = sqrt(11.56-5.76)=sqrt(5.8)=2.408
% the orbit
\draw[engblue, very thick] (0,0) ellipse (3.4 and 2.4);
% the focus (the attracting centre), to the right of geometric centre
\coordinate (F) at (2.408,0);
\fill[engamber] (F) circle (2.6pt);
\node[engamber, lbl, below right] at (F) {focus (centre of force)};
% periapsis sweep (near focus): wide arc on the near side
% near vertex at (a,0)=(3.4,0); sweep a thin wedge of large angular extent
\fill[engred!22]
  (F) -- (3.4,0)
  .. controls (3.25,0.95) and (2.7,1.5) .. (2.0,1.85)
  -- (F);
\draw[engred, thick] (F) -- (3.4,0);
\draw[engred, thick] (F) -- (2.0,1.85);
\node[engred, lbl] at (3.0,0.55) {fast};
% apoapsis sweep (far from focus): narrow arc, equal area, far side
% far vertex at (-a,0)=(-3.4,0)
\fill[englime!28]
  (F) -- (-3.4,0)
  .. controls (-3.3,-0.7) and (-3.05,-1.1) .. (-2.75,-1.45)
  -- (F);
\draw[englime!65!black, thick] (F) -- (-3.4,0);
\draw[englime!65!black, thick] (F) -- (-2.75,-1.45);
\node[englime!55!black, lbl] at (-2.0,-0.55) {slow};
% the body on each arc
\fill[engblack] (3.4,0) circle (2.0pt);
\fill[engblack] (-3.4,0) circle (2.0pt);
\end{tikzpicture}
\caption[Equal areas in equal times]{Kepler's second law as the
conservation of areal velocity. A body under a central force has no
moment about the centre of force, so its angular momentum $L=m
r^2\dot{\theta}$ is constant, and the rate at which the radius vector
sweeps area, $\tfrac{1}{2}r^2\dot{\theta}=L/(2m)$, is constant in time.
The two shaded wedges have equal area and are swept in equal times. Near
the focus the radius is short, so the body must travel through a wide
angle, and a long arc, to sweep that area; far from the focus the radius
is long, so a small angle, and a short arc, suffices. The body therefore
runs fast at periapsis and slow at apoapsis, the qualitative fact that
the areal-velocity law makes quantitative.}
\label{fig:vol03ch04:areal-velocity}
\end{figure}


% Figure: an elliptical orbit under an inverse-square central force, with
% the focus at the attracting body and equal-area sweeps marked.
\begin{figure}[htbp]
\centering
\begin{tikzpicture}[
  rad/.style={draw=enggray, thin},
  lbl/.style={font=\small},
]
% ellipse: a = 3.2, b = 2.4, c = sqrt(a^2-b^2)=sqrt(10.24-5.76)=sqrt(4.48)=2.117
% centre at origin; focus (Sun) at (c,0) = (2.117,0). Place ellipse centred (0,0).
\draw[draw=engblue, very thick] (0,0) ellipse (3.2 and 2.4);
% Sun at right focus
\coordinate (S) at (2.117,0);
\fill[engamber] (S) circle (4pt);
\node[engamber, lbl, below right] at (S) {Sun (focus)};
% perihelion and aphelion
\coordinate (peri) at (3.2,0);
\coordinate (apo) at (-3.2,0);
\fill[engblack] (peri) circle (1.3pt);
\fill[engblack] (apo) circle (1.3pt);
\node[lbl, right] at (peri) {perihelion};
\node[lbl, left] at (apo) {aphelion};
% Equal-area sweep near perihelion (small fast arc) shaded
\fill[engred!25] (S) -- (3.2,0) arc (0:34:3.2 and 2.4) -- cycle;
% the arc command above uses ellipse param; approximate with two radii + chord
\draw[rad] (S) -- (peri);
\coordinate (p2) at (2.654,1.342); % point on ellipse near peri
\draw[rad] (S) -- (p2);
% Equal-area sweep near aphelion (wide slow arc)
\fill[engred!25] (S) -- (-3.2,0) arc (180:120:3.2 and 2.4) -- cycle;
\coordinate (a2) at (-1.6,2.078); % point on ellipse near apo
\draw[rad] (S) -- (apo);
\draw[rad] (S) -- (a2);
% planet markers
\fill[engblue] (peri) circle (2.2pt);
\fill[engblue] (apo) circle (2.2pt);
% annotation
\node[engred, font=\scriptsize] at (3.0,1.0) {fast};
\node[engred, font=\scriptsize] at (-2.6,1.3) {slow};
\end{tikzpicture}
\caption[Elliptical orbit and Kepler's equal-area law]{An orbit under
an inverse-square central force is a closed ellipse with the
attracting body at one focus (Kepler's first law). Angular momentum
about the focus is conserved because the force is central, so the
radius vector sweeps equal areas in equal times (Kepler's second law,
the two shaded sectors): the orbiting body moves fast near
perihelion, where the radius is short, and slow near aphelion, where
the radius is long. The two-body solver \texttt{two\_body.py} produces
this curve numerically and recovers Kepler's third law $T^{2}\propto
a^{3}$ from the simulated period, the verification asked of the
central-force simulation exercise.}
\label{fig:vol03ch04:central-orbit}
\end{figure}


No closed-form expression gives the position as a function of time
(the relation between time and angle is Kepler's transcendental
equation), so the orbit is usually integrated numerically. The file
\texttt{two\_body.py} integrates Earth's orbit with a velocity-Verlet
step, a \engterm{symplectic} integrator whose energy error stays
bounded over many orbits rather than drifting as a Runge--Kutta scheme
would. Over $1.2$ years the relative energy drift is of order
$10^{-9}$, the orbit closes, and the period recovered from the energy
and the semi-major axis matches Kepler's third law to four figures.
The data file \texttt{orbit\_period.csv} tabulates $T^{2}/a^{3}$ for
the eight planets, a constant to three figures across a span of
semi-major axes from Mercury to Neptune. The central-force simulation
exercise asks the reader to reproduce this check.

\subsection{Maximum speed on a banked curve}

A car of mass $m$ rounds a curve of radius $r$ on a road banked at
angle $\theta$ to the horizontal. The tyre-road friction coefficient
is $\mu$. The fastest speed at which the car holds the curve without
sliding up and off the bank follows from the force balance in the
plane of the corner. Figure~\ref{fig:vol03ch04:banked-curve} draws the
free-body diagram.

\textbf{Force balance.} The acceleration is horizontal, of magnitude
$v^{2}/r$, directed toward the centre of the curve. Resolve the
weight $mg$, the normal force $N$, and the friction $f$ into horizontal
and vertical components. At the slide-up limit the friction acts
down-slope at its maximum $f=\mu N$, because the car tends to climb the
bank. The vertical balance (no vertical acceleration) and the
horizontal balance (centripetal) are
\[
N\cos\theta - \mu N\sin\theta - mg = 0, \qquad
N\sin\theta + \mu N\cos\theta = \frac{m v^{2}}{r}.
\]
Solve the first for $N=mg/(\cos\theta-\mu\sin\theta)$ and substitute
into the second:
\[
v_{\max}=\sqrt{g r\,\frac{\sin\theta+\mu\cos\theta}{\cos\theta-\mu
\sin\theta}}=\sqrt{g r\,\frac{\tan\theta+\mu}{1-\mu\tan\theta}}.
\]
The expression collects the three special cases cleanly. With no bank
($\theta=0$) it reduces to the flat-corner result $v_{\max}=\sqrt{\mu
g r}$. With no friction ($\mu=0$) it gives the design speed
$v_{0}=\sqrt{g r\tan\theta}$, the speed at which the bank alone turns
the car and the tyres carry no side load. With both present the two
contributions combine, and the denominator warns that the model breaks
when $\mu\tan\theta\to 1$ (the friction and bank together would let the
car corner at any speed, which signals that the small-angle, rigid-tyre
idealisation has been pushed past its range).

\textbf{Numbers.} Take a highway ramp of radius $r=120\,\si{\meter}$
banked at $\theta=8^{\circ}$ ($\tan\theta=0.1405$), dry asphalt
$\mu=0.75$. The design speed is $v_{0}=\sqrt{9.81\times 120\times
0.1405}=12.9\,\si{\meter\per\second}$, about $46\,\si{\kilo\meter\per
\hour}$, the speed at which a sheet of ice on the ramp would still let
the car corner. The friction-limited maximum is
\[
v_{\max}=\sqrt{9.81\times 120\,\frac{0.1405+0.75}{1-0.75\times 0.1405}}
=\sqrt{9.81\times 120\times 0.992}=34.2\,\si{\meter\per\second},
\]
about $123\,\si{\kilo\meter\per\hour}$. The bank buys little at high
grip and a great deal at low grip: on packed snow ($\mu=0.18$) the
same ramp limits the car to $v_{\max}=\sqrt{9.81\times 120\times
(0.1405+0.18)/(1-0.18\times 0.1405)}=20.2\,\si{\meter\per\second}$
($73\,\si{\kilo\meter\per\hour}$), so the posted advisory speed on a
banked ramp is set by the worst grip the road will see, not the dry
maximum.

\begin{mastery}
\textbf{The friction circle: braking and cornering share one budget.}
The banked-curve result and the braking estimate each spent the tyre's
whole grip on a single job, cornering in one, braking in the other. A
tyre in fact supplies a horizontal force in any direction, and the total
it can supply is bounded by one number, $\mu N$, whatever the direction.
The accessible forces therefore fill a disc of radius $\mu N$, the
\engterm{friction circle} (Figure~\ref{fig:vol03ch04:friction-circle}).
Longitudinal force, the braking or driving effort, and lateral force,
the cornering effort, are two components of one vector that must stay
inside the circle:
\[
\sqrt{F_{\text{long}}^{2}+F_{\text{lat}}^{2}}\le\mu N.
\]
The consequence is that braking and cornering compete for the same
budget. A tyre spending all of $\mu N$ on cornering has none left for
braking; a driver who brakes hard into a bend and then asks the tyre to
turn as well drives the demand vector outside the circle and the tyre
lets go. The racing craft of trail braking is exactly the management of
this constraint: the brake is eased as the steering builds, trading
longitudinal grip for lateral grip so the sum stays on the rim. The
same circle governs a braked car on a curved off-ramp, a motorcycle
tipping into a corner under trailing throttle, and an aircraft tyre on a
crosswind landing rollout. The flat-corner limit $v_{\max}=\sqrt{\mu g
r}$ and the straight-line braking limit $a=\mu g$ are the two axis
intercepts of the same circle; every real manoeuvre that mixes the two
lives on a chord between them, and the design rule is that the vector
sum, not either component alone, is what the grip must cover.
\end{mastery}

% Figure: the friction circle. The tyre can supply a horizontal force
% in any direction up to a magnitude mu*N; the accessible forces fill a
% disc. Braking, cornering, and combined manoeuvres sit inside it.
\begin{figure}[htbp]
\centering
\begin{tikzpicture}[
  vec/.style={draw=engblue, -{Stealth}, very thick},
  comp/.style={draw=enggray, thin, dashed},
  lbl/.style={font=\scriptsize},
]
% the friction limit disc, radius 3 = mu*N
\draw[draw=engred, very thick, fill=engred!6] (0,0) circle (3);
% axes: forward (up) is braking/traction, side (right) is cornering
\draw[-{Stealth}, engblack, thick] (0,-3.4) -- (0,3.6)
  node[above, lbl] {longitudinal (brake / drive)};
\draw[-{Stealth}, engblack, thick] (-3.4,0) -- (3.6,0)
  node[right, lbl] {lateral (cornering)};
% a combined-load demand vector: braking + cornering at 40 deg from vertical
\def\ang{40}
\draw[vec] (0,0) -- ({3*sin(\ang)}, {-3*cos(\ang)})
  node[below right, lbl, engblue] {demand $\vec{F}$};
% its components
\draw[comp] ({3*sin(\ang)},{-3*cos(\ang)}) -- (0,{-3*cos(\ang)});
\draw[comp] ({3*sin(\ang)},{-3*cos(\ang)}) -- ({3*sin(\ang)},0);
\node[lbl, enggray, anchor=east] at (0,{-1.5*cos(\ang)}) {$F_{\text{brake}}$};
\node[lbl, enggray, anchor=south] at ({1.5*sin(\ang)},0) {$F_{\text{corner}}$};
% radius label
\draw[-{Stealth}, engamber, thick] (0,0) -- ({3*cos(30)},{3*sin(30)});
\node[lbl, engamber, anchor=south west] at ({1.6*cos(30)},{1.6*sin(30)})
  {$\mu N$};
% a point outside the circle (loss of grip)
\fill[engblack] ({3.4*sin(20)},{3.4*cos(20)}) circle (1.5pt);
\node[lbl, anchor=west] at ({3.4*sin(20)},{3.4*cos(20)})
  {outside: tyre slides};
\end{tikzpicture}
\caption[The friction circle for combined tyre loading]{The friction
circle. A tyre can return a horizontal force in any direction up to a
magnitude $\mu N$, so the forces it can supply fill a disc of radius
$\mu N$ (red). Pure braking spends the whole budget along the vertical
axis; pure cornering spends it along the horizontal axis; a driver who
brakes and turns at once spends it on a slant, and the two demands add
as vectors. A demand vector that reaches the rim uses all the available
grip; one that lies inside has grip to spare; one that would fall
outside cannot be met and the tyre slides. The circle is why trail
braking into a corner must ease the brake as steering builds: the sum
of the two, not either alone, must stay within the rim.}
\label{fig:vol03ch04:friction-circle}
\end{figure}


\subsection{A loop-the-loop by the work-energy theorem}

A small cart is released from rest at height $h$ on a frictionless
ramp and runs into a vertical circular loop of radius $R$
(Figure~\ref{fig:vol03ch04:loop}). The question that decides whether
the ride works is the minimum release height for which the cart keeps
contact with the track all the way around.

\textbf{The binding constraint is at the top.} At the top of the loop
the cart moves horizontally and the centre of the circle is directly
below it, so the net downward force supplies the centripetal
acceleration:
\[
mg + N = \frac{m v_{\text{top}}^{2}}{R},
\]
with $N$ the track's push (inward, downward at the top). Contact holds
only while $N\ge 0$. The slowest survivable top speed is the one that
just zeroes the normal force,
\[
N=0 \ \Rightarrow\ v_{\text{top}}^{2}=gR,
\]
at which gravity alone turns the cart. A cart slower than this at the
top needs an inward force larger than gravity can supply; the track
cannot pull, so the cart leaves the rail and falls inward before
reaching the top.

\textbf{Energy fixes the height.} The ramp and loop are frictionless,
so mechanical energy is conserved between release and the top. Measure
height from the bottom of the loop; the top sits at $2R$:
\[
mgh = \tfrac{1}{2}m v_{\text{top}}^{2}+mg(2R)
= \tfrac{1}{2}m(gR)+2mgR = \tfrac{5}{2}mgR.
\]
The mass cancels and the minimum release height is $h_{\min}=2.5R$,
half a radius above the top of the loop. A loop of radius
$R=5\,\si{\meter}$ needs a release height of at least $12.5\,\si{
\meter}$. Real rides clear the constraint with a margin and use a
clothoid (a curve of continuously varying radius) rather than a true
circle, so that the rider meets the tightest radius, and the highest
acceleration, only briefly near the top; the constant-radius loop of
this idealisation would subject the rider to its peak load over a long
arc at the bottom. The bottom-of-loop normal force, where the speed is
highest, is the load the structure must carry: at $h=2.5R$ the speed at
the bottom is $v_{\text{bot}}^{2}=2gh=5gR$, and $N_{\text{bot}}=mg+
mv_{\text{bot}}^{2}/R=6mg$, six times the cart's weight.

\subsection{A constrained pulley-cart system}

A cart of mass $m_{1}$ sits on a frictionless horizontal table.
A light, inextensible cord runs from the cart, over a frictionless
pulley at the table's edge, and down to a hanging block of mass
$m_{2}$. Released from rest, the block falls and pulls the cart. The
problem shows how a kinematic constraint closes a system that has more
unknowns than free bodies.

\textbf{Constraint first.} The cord is inextensible, so the cart's
horizontal displacement equals the block's vertical drop at every
instant; differentiating twice, the two bodies share one acceleration
magnitude $a$. The cord tension $T$ is the same throughout an ideal
massless cord over a frictionless pulley.

\textbf{Two free bodies.} For the cart (horizontal), the only
horizontal force is the tension:
\[
T = m_{1} a.
\]
For the block (vertical, taking down as positive), gravity pulls down
and the tension up:
\[
m_{2} g - T = m_{2} a.
\]
Add the two equations to cancel $T$:
\[
m_{2} g = (m_{1}+m_{2}) a \ \Rightarrow\
a = \frac{m_{2}}{m_{1}+m_{2}}\,g, \qquad
T = \frac{m_{1} m_{2}}{m_{1}+m_{2}}\,g.
\]
The acceleration is the hanging weight divided by the total inertia,
and the tension carries the reduced-mass combination $m_{1}m_{2}/
(m_{1}+m_{2})$ times $g$. Two limits check the result. A vanishing cart
($m_{1}\to 0$) gives $a\to g$ and $T\to 0$: the block is in free fall
and the cord goes slack. A massive cart ($m_{1}\gg m_{2}$) gives $a\to
0$ and $T\to m_{2}g$: the block hangs nearly static and the cord
carries its full weight. With $m_{1}=3\,\si{\kilogram}$ and
$m_{2}=2\,\si{\kilogram}$, $a=2(9.81)/5=3.92\,\si{\meter\per\second
\squared}$ and $T=6(9.81)/5=11.8\,\si{\newton}$, a tension between the
block's weight $19.6\,\si{\newton}$ and zero, as both limits require.

\subsection{Relative motion: crossing a current}

A boat moves at $3\,\si{\meter\per\second}$ relative to the water in a
river $200\,\si{\meter}$ wide; the current flows at $1.5\,\si{\meter\per
\second}$. Two questions separate the two natural steering choices: how
far downstream does the boat land if it points straight across, and at
what heading must it aim to land directly opposite?
Figure~\ref{fig:vol03ch04:river-crossing} draws both as vector
triangles.

\textbf{The relative-velocity law.} Velocities add as vectors:
$\vec{v}_{b/g}=\vec{v}_{b/w}+\vec{v}_{w/g}$, the boat's velocity
relative to the ground equals its velocity relative to the water plus
the water's velocity relative to the ground. The whole problem is the
reading of this one triangle.

\textbf{Pointing straight across.} The across-component of the ground
velocity is the full boat speed $3\,\si{\meter\per\second}$, since the
current is purely downstream and adds nothing across. The crossing time
is $t=200/3=66.7\,\si{\second}$. During that time the current carries
the boat downstream by $1.5\times 66.7=100\,\si{\meter}$. The boat
lands $100\,\si{\meter}$ downstream, on a track angled $\arctan(1.5/3)=
26.6^{\circ}$ from the perpendicular, and its speed over the ground is
$\sqrt{3^{2}+1.5^{2}}=3.35\,\si{\meter\per\second}$. This heading gives
the shortest crossing time, because the entire boat speed is spent on
the width.

\textbf{Landing directly opposite.} To cancel the downstream drift the
boat must aim upstream so that the across-water velocity has an
upstream component equal to the current. The required heading angle
$\phi$ upstream of the perpendicular satisfies $v_{b/w}\sin\phi=
v_{w/g}$, so $\sin\phi=1.5/3=0.5$ and $\phi=30^{\circ}$. The
across-component that remains is $v_{b/w}\cos\phi=3\cos 30^{\circ}=
2.60\,\si{\meter\per\second}$, so the crossing now takes $t=200/2.60=
77.0\,\si{\second}$, longer than the straight-across crossing by the
factor $1/\cos 30^{\circ}=1.15$. The boat trades $15\,\%$ more time for
a landing with no downstream offset. A heading any steeper than
$30^{\circ}$ would push the boat upstream; if the current exceeded the
boat speed, no heading could hold station and a direct-opposite landing
would be impossible, the same condition that decides whether a swimmer
can ferry across a fast river.

% Figure: relative-motion / river-crossing vector triangle. The boat's
% velocity relative to the water adds to the water's velocity relative
% to the ground to give the boat's velocity relative to the ground.
% Two steering choices are shown: heading straight across (swept
% downstream) and heading upstream to land directly opposite.
\begin{figure}[htbp]
\centering
\begin{tikzpicture}[
  vb/.style={draw=engblue, -{Stealth}, very thick},
  vw/.style={draw=engamber, -{Stealth}, very thick},
  vg/.style={draw=englime!60!black, -{Stealth}, very thick},
  lbl/.style={font=\scriptsize},
]
% banks
\draw[draw=enggray, thick] (-0.4,0) -- (8.5,0);
\draw[draw=enggray, thick] (-0.4,4) -- (8.5,4);
\fill[engblue!6] (-0.4,0) rectangle (8.5,4);
\node[lbl, enggray, anchor=west] at (-0.4,-0.3) {near bank};
\node[lbl, enggray, anchor=west] at (-0.4,4.3) {far bank};
% current direction arrow (downstream, to the right)
\node[lbl, engamber, anchor=west] at (5.6,2) {current $\vec{v}_{w}$};
\draw[vw] (5.4,2) -- ++(1.6,0);

% Case A: head straight across, swept downstream.
\coordinate (A0) at (0.6,0);
% boat rel water: straight up, magnitude 3 (scaled: 3 -> 3 units)
\draw[vb] (A0) -- ++(0,3) node[midway, left, lbl, engblue] {$\vec{v}_{b/w}$};
% water rel ground: to the right, magnitude 1.5
\draw[vw] (A0) ++(0,3) -- ++(1.5,0);
% boat rel ground: resultant from A0 to the tip
\draw[vg] (A0) -- ++(1.5,3) node[midway, right, lbl, englime!50!black] {$\vec{v}_{b/g}$};
\node[lbl, anchor=north] at (0.6,-0.1) {(a) point straight across};

% Case B: head upstream so resultant is straight across.
\coordinate (B0) at (4.3,0);
% boat rel water aimed up-left so that across-component lands opposite:
% magnitude 3, with x-component -1.5 to cancel current; y = sqrt(3^2-1.5^2)=2.598
\draw[vb] (B0) -- ++(-1.5,2.598) node[midway, left, lbl, engblue] {$\vec{v}_{b/w}$};
% water rel ground to the right, magnitude 1.5
\draw[vw] (B0) ++(-1.5,2.598) -- ++(1.5,0);
% boat rel ground: straight up, magnitude 2.598
\draw[vg] (B0) -- ++(0,2.598) node[midway, right, lbl, englime!50!black] {$\vec{v}_{b/g}$};
\node[lbl, anchor=north] at (4.3,-0.1) {(b) angle upstream};
\end{tikzpicture}
\caption[Relative-motion vector triangle for a river crossing]{The
two steering choices for crossing a current, each a closed vector
triangle of the relative-motion law $\vec{v}_{b/g}=\vec{v}_{b/w}+
\vec{v}_{w/g}$. In (a) the boat points straight across, so its velocity
relative to the water (blue) is perpendicular to the banks; the current
(amber) adds a downstream component and the velocity relative to the
ground (green) carries the boat across on a slanted track, landing
downstream of the start. The crossing time is shortest here because the
full speed counts toward the width. In (b) the boat angles upstream by
just enough that the current cancels the sideways drift, so the
ground velocity points straight across and the landing is directly
opposite; the price is a smaller across-component and a longer crossing
time. The triangle is the same law that governs an aircraft crabbing
into a crosswind.}
\label{fig:vol03ch04:river-crossing}
\end{figure}


\subsection{Straight-line car performance: traction and power limits}

A rear-driven car of mass $m=1{,}400\,\si{\kilogram}$ has an engine
that delivers $P=150\,\si{\kilo\watt}$ at the wheels and tyres with a
drive-axle friction coefficient $\mu=0.9$. Two questions decide its
straight-line performance: what limits the launch off the line, and
what sets the top speed. Figure~\ref{fig:vol03ch04:traction-power}
draws the two ceilings on the tractive force.

\textbf{The launch is traction-limited.} At low speed the engine could
deliver an enormous force, $F=P/v$, which grows without bound as
$v\to 0$, but the tyres cannot put more than $\mu m g$ to the ground
before they spin. The launch force is therefore the friction limit
\[
F_{\text{trac}}=\mu m g=0.9\times 1400\times 9.81=12.4\,\si{\kilo
\newton},
\]
giving a launch acceleration $a=\mu g=0.9\times 9.81=8.83\,\si{\meter
\per\second\squared}$, about $0.9g$. The car is grip-bound, not
power-bound, off the line: a more powerful engine would only spin the
tyres faster, not launch harder.

\textbf{The crossover speed.} The traction ceiling and the power
ceiling cross where $P/v=\mu m g$:
\[
v_{\times}=\frac{P}{\mu m g}=\frac{150{,}000}{12{,}400}=12.1\,\si{
\meter\per\second},
\]
about $44\,\si{\kilo\meter\per\hour}$. Below this speed the tyres are
the limit and the force is the flat $12.4\,\si{\kilo\newton}$; above it
the engine is the limit and the force falls along the hyperbola $P/v$.
The first stretch of a hard launch holds constant acceleration; past
the crossover the acceleration tapers as the power-limited force
declines.

\textbf{The top speed is power-limited.} Top speed is reached where the
power-limited tractive force has fallen to meet the road load, the sum
of aerodynamic drag and rolling resistance. Modelling the road load as
$R(v)=R_{0}+\tfrac{1}{2}\rho C_{d}A\,v^{2}$ with $R_{0}\approx
200\,\si{\newton}$ rolling and $\tfrac{1}{2}\rho C_{d}A\approx
0.60\,\si{\newton\second\squared\per\meter\squared}$ (a $C_{d}A$ near
$1\,\si{\meter\squared}$ in air), the balance $P/v=R(v)$ gives
$150{,}000/v=200+0.60\,v^{2}$. Solving numerically, $v_{\max}\approx
61\,\si{\meter\per\second}$, about $220\,\si{\kilo\meter\per\hour}$. The
cubic growth of the drag-times-speed power demand is why doubling top
speed needs roughly eight times the power: $P=R v\approx
\tfrac{1}{2}\rho C_{d}A\,v^{3}$ at high speed.

\textbf{The launch time to $100\,\si{\kilo\meter\per\hour}$.} A rough
$0$ to $100\,\si{\kilo\meter\per\hour}$ time splits at the crossover.
From rest to $v_{\times}=12.1\,\si{\meter\per\second}$ the acceleration
is the constant $8.83\,\si{\meter\per\second\squared}$, taking
$t_{1}=12.1/8.83=1.37\,\si{\second}$. From $v_{\times}$ to $27.8\,\si{
\meter\per\second}$ the force is power-limited; using the work-energy
form $P\,t_{2}\approx\tfrac{1}{2}m(v_{2}^{2}-v_{\times}^{2})$ as an
estimate ignoring road load, $t_{2}\approx\tfrac{1}{2}(1400)(27.8^{2}-
12.1^{2})/150{,}000=2.92\,\si{\second}$. The total is about
$4.3\,\si{\second}$, in the range a $150\,\si{\kilo\watt}$ car of this
mass actually records, and the estimate makes plain that the launch
phase is grip-bound while the rest is power-bound.

% Figure: straight-line vehicle performance. The tractive force a car
% can apply is the lesser of the traction limit (a flat ceiling set by
% tyre grip) and the power limit (a hyperbola F = P/v that falls with
% speed). The crossover speed separates the traction-limited launch
% from the power-limited top end.
\begin{figure}[htbp]
\centering
\begin{tikzpicture}
\begin{axis}[
  width=12cm, height=7cm,
  xlabel={speed $v$ (\si{\meter\per\second})},
  ylabel={available tractive force (\si{\kilo\newton})},
  xmin=0, xmax=70, ymin=0, ymax=14,
  axis lines=left,
  domain=5:70, samples=140,
  legend pos=north east,
  legend cell align=left,
  grid=major, grid style={enggray!20},
  tick label style={font=\scriptsize},
  label style={font=\small},
  legend style={font=\scriptsize},
]
% Traction limit: mu m g = 0.9 * 1400 * 9.81 = 12.36 kN, flat.
\addplot[engred, very thick, domain=0:70] {12.36};
\addlegendentry{traction limit $\mu m g$}
% Power limit: P/v, P = 150 kW = 150000 W, in kN -> 150/v.
\addplot[engblue, very thick] {150/x};
\addlegendentry{power limit $P/v$}
% Resistance (drag + rolling): R = 0.2 + 0.0006*v^2 kN (rough).
\addplot[enggray, thick, dashed, domain=0:70] {0.2 + 0.0006*x*x};
\addlegendentry{road load (drag + rolling)}
% crossover of traction and power limits: 150/v = 12.36 -> v = 12.1 m/s
\addplot[only marks, mark=*, engblack, mark size=2pt] coordinates {(12.1,12.36)};
\node[font=\scriptsize, anchor=south west] at (axis cs:12.1,12.36)
  {crossover $v\approx 12\,\si{\meter\per\second}$};
% top speed: power limit meets road load near v = 61 m/s
\addplot[only marks, mark=square*, engblack, mark size=2pt] coordinates {(61,2.46)};
\node[font=\scriptsize, anchor=south east] at (axis cs:61,2.46)
  {top speed};
\end{axis}
\end{tikzpicture}
\caption[Traction-limited and power-limited vehicle performance]{The
straight-line performance of a car is governed by two ceilings on the
tractive force. Below the crossover speed the tyres set the limit: the
ground can return at most $\mu m g$ regardless of engine output, so the
launch is traction-limited and the acceleration is the friction
maximum. Above the crossover the engine sets the limit: a fixed power
$P$ delivers force $P/v$, a hyperbola that falls as the car speeds up,
so the high-speed pull is power-limited. The two limits cross near
$12\,\si{\meter\per\second}$ for the numbers in the worked example.
Top speed is reached where the power-limited pull (blue) drops to meet
the road load (grey dashed), the sum of aerodynamic drag and rolling
resistance; no surplus force remains to accelerate.}
\label{fig:vol03ch04:traction-power}
\end{figure}


\subsection{Two blocks with friction by Newton's second law}

A block of mass $m_{1}=4\,\si{\kilogram}$ rests on a rough horizontal
table and is connected by a light cord over a frictionless pulley at
the table edge to a hanging block of mass $m_{2}=3\,\si{\kilogram}$.
The coefficient of kinetic friction between the table block and the
table is $\mu_{k}=0.25$. Find the acceleration of the system and the
cord tension once the blocks are moving.

\textbf{Free bodies and the constraint.} The inextensible cord ties the
two accelerations to a common magnitude $a$, and an ideal cord over a
frictionless pulley carries a single tension $T$. The table block feels
the tension forward, friction $f=\mu_{k}N$ backward, and a normal force
$N=m_{1}g$ from the table (no vertical acceleration). The hanging block
feels gravity down and the tension up.

\textbf{The two equations.} For the table block (horizontal),
\[
T-\mu_{k}m_{1}g=m_{1}a,
\]
and for the hanging block (down positive),
\[
m_{2}g-T=m_{2}a.
\]
Add the two to cancel $T$:
\[
m_{2}g-\mu_{k}m_{1}g=(m_{1}+m_{2})a
\ \Rightarrow\
a=\frac{(m_{2}-\mu_{k}m_{1})\,g}{m_{1}+m_{2}}.
\]
With the numbers, $a=(3-0.25\times 4)(9.81)/(7)=(3-1)(9.81)/7=
2.80\,\si{\meter\per\second\squared}$. The tension follows from the
hanging-block equation, $T=m_{2}(g-a)=3(9.81-2.80)=21.0\,\si{\newton}$,
a value between the hanging weight $29.4\,\si{\newton}$ and zero, as it
must be for a block that is accelerating downward but not in free fall.

\textbf{The motion-onset check.} The blocks move only if the driving
weight exceeds the static friction the table can muster. With a static
coefficient $\mu_{s}\approx 0.30$, the holding friction is at most
$\mu_{s}m_{1}g=0.30\times 4\times 9.81=11.8\,\si{\newton}$, while the
hanging weight is $m_{2}g=29.4\,\si{\newton}$. The drive beats the
hold, so the system accelerates; had the hanging block been light
enough that $m_{2}g<\mu_{s}m_{1}g$, the blocks would have stayed at rest
and the tension would have equalled the hanging weight with no motion.

\subsection{Simple-harmonic motion from a restoring force law}

A particle of mass $m$ moves under a force $F(x)=-kx$ that always points
back toward the origin in proportion to displacement. The shape of the
motion follows from the force law alone, without solving for the
trajectory by integration.

\textbf{The equation of motion.} Newton's second law gives $m\ddot{x}=
-kx$, which rearranges to $\ddot{x}+\omega_{n}^{2}x=0$ with
$\omega_{n}=\sqrt{k/m}$. The defining feature of simple-harmonic motion
is that the acceleration is proportional to the displacement and
opposite in sign; any force law linear in displacement and restoring in
direction produces it, which is why the pendulum at small angle, the
mass on a spring, the floating hydrometer bobbing in water, and the
charge in an inductor-capacitor loop all obey the same equation.

\textbf{The solution and its constants.} The general solution is
$x(t)=A\cos(\omega_{n}t)+B\sin(\omega_{n}t)$, or equivalently
$x(t)=C\cos(\omega_{n}t-\varphi)$ with amplitude $C=\sqrt{A^{2}+B^{2}}$
and phase $\varphi=\arctan(B/A)$. The angular frequency $\omega_{n}=
\sqrt{k/m}$ depends only on the system, not on the amplitude: a
stiffer spring or a lighter mass oscillates faster, and the period
$T=2\pi\sqrt{m/k}$ is the same whether the swing is large or small.
That amplitude-independence is the property that made the pendulum a
clock.

\textbf{A worked instance.} A $0.5\,\si{\kilogram}$ mass on a spring of
stiffness $k=200\,\si{\newton\per\meter}$ has $\omega_{n}=\sqrt{200/
0.5}=20\,\si{\radian\per\second}$, period $T=2\pi/20=0.314\,\si{
\second}$, and frequency $1/T=3.18\,\si{\hertz}$. Released from a
displacement of $0.04\,\si{\meter}$ at rest, the motion is $x(t)=0.04
\cos(20t)$, with peak speed $v_{\max}=\omega_{n}C=20\times 0.04=
0.80\,\si{\meter\per\second}$ at the centre and peak acceleration
$a_{\max}=\omega_{n}^{2}C=400\times 0.04=16\,\si{\meter\per\second
\squared}$ at the extremes. The energy
$\tfrac{1}{2}kC^{2}=\tfrac{1}{2}(200)(0.04)^{2}=0.16\,\si{\joule}$ is
constant and trades between kinetic and potential twice per cycle.

\subsection{Normal force on a vehicle over a crest and through a dip}

A car of mass $m$ travels at speed $v$ over a road whose profile curves.
At a crest of radius $R$ the road falls away beneath the car; in a dip
of radius $R$ it curves upward. The question that matters for tyre grip
and passenger comfort is the normal force the road exerts, which is what
a bathroom scale under the wheels would read.

\textbf{At the crest.} The centre of the circular arc is below the road,
so the centripetal acceleration points downward with magnitude
$v^{2}/R$. The vertical balance has gravity down and the normal force
$N$ up, with the net pointing down toward the centre:
\[
mg-N=\frac{mv^{2}}{R}
\ \Rightarrow\
N=m\!\left(g-\frac{v^{2}}{R}\right).
\]
The normal force is less than the weight: the car feels light over a
crest, and the passengers feel the float in the stomach. If the speed
reaches $v=\sqrt{gR}$, the normal force drops to zero and the car
leaves the road, the same $v^{2}=gR$ condition that set the minimum
speed at the top of the loop. For a crest of radius $R=80\,\si{\meter}$
the lift-off speed is $\sqrt{9.81\times 80}=28.0\,\si{\meter\per\second}$,
about $101\,\si{\kilo\meter\per\hour}$, which is why a humped bridge
taken too fast launches a car.

\textbf{In the dip.} The centre of the arc is now above the road, so
the centripetal acceleration points upward and the balance reverses:
\[
N-mg=\frac{mv^{2}}{R}
\ \Rightarrow\
N=m\!\left(g+\frac{v^{2}}{R}\right).
\]
The normal force exceeds the weight: the car and its passengers feel
heavy at the bottom of a dip, the same press felt at the bottom of the
loop. At $v=20\,\si{\meter\per\second}$ through a dip of $R=80\,\si{
\meter}$, the load factor is $N/(mg)=1+v^{2}/(gR)=1+400/(9.81\times 80)
=1.51$, so the suspension, the tyres, and the occupants carry half
again their static weight. The crest and the dip are one formula read
with the centre of curvature on opposite sides, and the magnitude
$v^{2}/R$ is the cornering acceleration of the path-coordinate
decomposition met at the start of the chapter.

\subsection{Variable acceleration: motion from a time-varying force}

The constant-acceleration formulas fail the moment the force changes
with time, and the recovery is to return to the definitions and
integrate. A test sled of mass $m=200\,\si{\kilogram}$ on a level track
feels a thrust that builds and then fades, $F(t)=F_{0}(1-t/\tau)$ for
$0\le t\le\tau$ and zero thereafter, with $F_{0}=1{,}600\,\si{\newton}$
and $\tau=4\,\si{\second}$. The sled starts from rest. The task is the
velocity and position at the end of the thrust, found by integrating
the acceleration rather than by reaching for $v=v_{0}+at$.

\textbf{Acceleration, then two integrals.} Newton's second law gives
the acceleration as a function of time,
\[
a(t)=\frac{F(t)}{m}=\frac{F_{0}}{m}\Bigl(1-\frac{t}{\tau}\Bigr)
=8.0\Bigl(1-\frac{t}{4}\Bigr)\,\si{\meter\per\second\squared},
\]
which starts at $8.0\,\si{\meter\per\second\squared}$ and falls
linearly to zero at $t=\tau$. Velocity is the running integral of
acceleration from rest,
\[
v(t)=\int_{0}^{t}a(t')\,dt'=\frac{F_{0}}{m}\Bigl(t-\frac{t^{2}}{2\tau}
\Bigr),
\]
and at $t=\tau$ this is $v(\tau)=\frac{F_{0}}{m}\cdot\frac{\tau}{2}=
8.0\times 2=16\,\si{\meter\per\second}$. The factor $\tau/2$ is the
area under the triangular acceleration curve, the same graphical
reading the chapter opened with. Position is the integral of velocity,
\[
x(t)=\int_{0}^{t}v(t')\,dt'=\frac{F_{0}}{m}\Bigl(\frac{t^{2}}{2}
-\frac{t^{3}}{6\tau}\Bigr),
\]
and at $t=\tau$, $x(\tau)=\frac{F_{0}}{m}\bigl(\frac{\tau^{2}}{2}-
\frac{\tau^{2}}{6}\bigr)=\frac{F_{0}}{m}\cdot\frac{\tau^{2}}{3}=8.0
\times\frac{16}{3}=42.7\,\si{\meter}$. A constant-acceleration estimate
using the average acceleration $\bar{a}=4.0\,\si{\meter\per\second
\squared}$ over $4\,\si{\second}$ would give the same final velocity
$\bar{a}\tau=16\,\si{\meter\per\second}$, because the mean of a linear
ramp is its midpoint value, but it would miss the position by giving
$\tfrac{1}{2}\bar{a}\tau^{2}=32\,\si{\meter}$ against the true
$42.7\,\si{\meter}$. The discrepancy is the warning: averaging the
acceleration recovers the impulse and so the final velocity, but it
discards the time history that the position integral needs, because the
sled spends its early seconds at the high end of the acceleration ramp
where it covers the most ground.

\textbf{The cross-check by impulse.} The impulse-momentum theorem gives
the final velocity without integrating twice. The impulse is the area
under $F(t)$, a triangle of height $F_{0}$ and base $\tau$, so
$J=\tfrac{1}{2}F_{0}\tau=\tfrac{1}{2}(1600)(4)=3{,}200\,\si{\newton
\second}$, and $v(\tau)=J/m=3200/200=16\,\si{\meter\per\second}$,
matching the integral. When only the final velocity is wanted, the
impulse is the shorter road; the position still requires the second
integral.

\subsection{A two-dimensional collision with restitution}

A real collision rarely lands the two bodies on one line, and the
two-dimensional case shows how the momentum vector and the restitution
coefficient together close a problem that momentum alone underdetermines.
A puck of mass $m_{1}=0.50\,\si{\kilogram}$ slides at $v_{1}=4.0\,\si{
\meter\per\second}$ along the $x$-axis and strikes a stationary puck of
mass $m_{2}=0.30\,\si{\kilogram}$. The contact is off-centre, so the
line joining the centres at impact (the line of impact) makes an angle
$\theta=30^{\circ}$ with the incoming velocity. The coefficient of
restitution along the line of impact is $e=0.80$. Find both pucks'
velocities after the collision.

\textbf{Resolve onto the line of impact and the tangent.} Frictionless
pucks exert force only along the line of impact, so the tangential
velocity components are unchanged and all the collision physics lives
on the line of impact, here called the $n$-axis. Decompose the incoming
velocity of puck 1:
\[
v_{1n}=v_{1}\cos\theta=4.0\cos 30^{\circ}=3.46\,\si{\meter\per\second},
\qquad
v_{1t}=v_{1}\sin\theta=4.0\sin 30^{\circ}=2.00\,\si{\meter\per\second}.
\]
Puck 2 starts at rest, so $v_{2n}=v_{2t}=0$.

\textbf{Two equations on the line of impact.} Along the $n$-axis,
momentum is conserved and restitution relates the separation speed to
the approach speed:
\[
m_{1}v_{1n}=m_{1}v_{1n}'+m_{2}v_{2n}', \qquad
v_{2n}'-v_{1n}'=e\,(v_{1n}-v_{2n})=e\,v_{1n}.
\]
Solving the pair,
\[
v_{1n}'=\frac{m_{1}-e\,m_{2}}{m_{1}+m_{2}}\,v_{1n}
=\frac{0.50-0.80(0.30)}{0.80}(3.46)=1.13\,\si{\meter\per\second},
\]
\[
v_{2n}'=\frac{m_{1}(1+e)}{m_{1}+m_{2}}\,v_{1n}
=\frac{0.50(1.80)}{0.80}(3.46)=3.90\,\si{\meter\per\second}.
\]

\textbf{Reassemble each velocity.} The tangential components ride
through untouched, so puck 1 leaves with $(v_{1n}',v_{1t})=(1.13,2.00)
\,\si{\meter\per\second}$, a speed of $\sqrt{1.13^{2}+2.00^{2}}=2.30
\,\si{\meter\per\second}$ at $\arctan(2.00/1.13)=60.5^{\circ}$ off the
line of impact, and puck 2 leaves with $(3.90,0)\,\si{\meter\per\second}$
straight along the line of impact at $3.90\,\si{\meter\per\second}$. As
a check, total $x$-momentum before is $m_{1}v_{1}=2.00\,\si{\kilogram
\meter\per\second}$; after, projecting both outgoing velocities back
onto the original $x$-direction recovers the same total to within
rounding. The kinetic energy falls from $\tfrac{1}{2}(0.50)(4.0)^{2}=
4.00\,\si{\joule}$ to $\tfrac{1}{2}(0.50)(2.30)^{2}+\tfrac{1}{2}(0.30)
(3.90)^{2}=1.32+2.28=3.60\,\si{\joule}$, a loss of $0.40\,\si{\joule}$,
which is the energy a restitution below one removes from the line of
impact only; the tangential motion carries no loss because no force
acted along it.

\subsection{A conical pendulum}

A bob on a string swept into a steady horizontal circle is the cleanest
problem in which a single tension does two jobs at once, holding the bob
up and turning it. A bob of mass $m=0.40\,\si{\kilogram}$ hangs on a
string of length $L=1.2\,\si{\meter}$ and is set swinging in a
horizontal circle so that the string makes a half-angle $\theta=
35^{\circ}$ with the vertical (Figure~\ref{fig:vol03ch04:conical-pendulum}).
Find the bob's speed, the period of its circular motion, and the string
tension.

\textbf{Two balances, not three.} The bob has no vertical acceleration,
so the vertical component of the tension carries the weight; the
horizontal acceleration is the centripetal $v^{2}/r$ toward the axis,
supplied by the horizontal component of the tension. With $r=L\sin
\theta$ the circle radius,
\[
T\cos\theta=mg, \qquad T\sin\theta=\frac{m v^{2}}{r}=\frac{m v^{2}}
{L\sin\theta}.
\]
Dividing the second by the first removes the tension and gives the
working relation $\tan\theta=v^{2}/(gr)$, the same expression that set
the design speed of the banked curve, with the bank angle replaced by
the cone angle.

\textbf{Carry the numbers.} The radius is $r=1.2\sin 35^{\circ}=0.688
\,\si{\meter}$. The speed follows from $v=\sqrt{gr\tan\theta}=\sqrt{9.81
(0.688)(0.700)}=2.17\,\si{\meter\per\second}$. The period is the
circumference over the speed,
\[
T_{\text{period}}=\frac{2\pi r}{v}=\frac{2\pi(0.688)}{2.17}=1.99\,\si{
\second},
\]
which the compact formula $T_{\text{period}}=2\pi\sqrt{L\cos\theta/g}=
2\pi\sqrt{1.2(0.819)/9.81}=1.99\,\si{\second}$ reproduces. The string
tension is $T=mg/\cos\theta=0.40(9.81)/0.819=4.79\,\si{\newton}$, half
again the bob's $3.92\,\si{\newton}$ weight, the surcharge for turning
the bob as well as holding it. The period depends on the cone angle
only through $\cos\theta$, so it falls as the cone opens and the bob
swings faster, the property that makes the conical pendulum a speed
governor: a steam engine's flyball governor opens its cone as the
engine speeds up and uses the rising bob to throttle the steam.

\subsection{A particle on a rotating turntable: the Coriolis term}

The polar-coordinate acceleration that opened the chapter carried a
Coriolis term whose meaning is sharpest when a particle moves radially
on a spinning platform. A turntable rotates at constant
$\Omega=4.0\,\si{\radian\per\second}$. A bead slides outward along a
frictionless radial groove at a constant radial speed $\dot{r}=0.50\,
\si{\meter\per\second}$ relative to the turntable. At the instant the
bead is at radius $r=0.30\,\si{\meter}$, find the force the groove wall
exerts on the bead and the bead's acceleration in the ground frame.

\textbf{Read the acceleration off the polar formula.} The groove holds
$\theta$ tied to the turntable, so $\dot{\theta}=\Omega=4.0\,\si{\radian
\per\second}$, $\ddot{\theta}=0$, $\dot{r}=0.50\,\si{\meter\per\second}$
constant so $\ddot{r}=0$. The polar acceleration components are
\[
a_{r}=\ddot{r}-r\dot{\theta}^{2}=0-(0.30)(4.0)^{2}=-4.8\,\si{\meter\per
\second\squared},
\]
\[
a_{\theta}=r\ddot{\theta}+2\dot{r}\dot{\theta}=0+2(0.50)(4.0)=4.0\,\si{
\meter\per\second\squared}.
\]
The radial component is the centripetal $-r\Omega^{2}$ pulling the bead
toward the axis; the transverse component is the Coriolis $2\dot{r}
\Omega$, present only because the bead moves radially while the platform
turns.

\textbf{The forces that produce them.} For a bead of mass $m=0.20\,\si{
\kilogram}$, the transverse force is $F_{\theta}=ma_{\theta}=0.20(4.0)=
0.80\,\si{\newton}$, and it can come only from the groove wall, because
the groove is the sole thing that can push sideways on the bead. The
wall therefore presses on the bead with $0.80\,\si{\newton}$ in the
direction of rotation, and by the third law the bead presses back on
the wall with the same force opposing the rotation, the load that a
real radial slot must carry. The radial force is $F_{r}=ma_{r}=0.20(
-4.8)=-0.96\,\si{\newton}$, an inward pull; in the frictionless groove
nothing supplies an inward radial force, so the bead is not in radial
equilibrium and $\ddot{r}$ would in fact grow once the constraint of
constant $\dot{r}$ is released. Holding $\dot{r}$ constant, as the
problem stipulates, requires an external inward agent supplying that
$0.96\,\si{\newton}$. The lesson the turntable teaches is the one the
rotating-frame mastery box stated abstractly: a body moving in a
rotating frame feels a Coriolis push perpendicular to its relative
velocity, and an engineer designing any radial-flow rotating machine,
a centrifugal pump impeller or a lawn-sprinkler arm, must carry that
sideways load in the vane (Figure~\ref{fig:vol03ch04:coriolis}).

\subsection{A uniform rod released from horizontal}

A rigid rod pivoted at one end and let fall is the cleanest rigid-body
problem in which translation and rotation are locked together by a
single constraint, and it pays the analyst back with a result that
surprises every first-time reader. A uniform rod of mass $m$ and length
$L$ is hinged at one end on a frictionless pin and released from rest in
the horizontal position. Find the angular acceleration at the instant of
release, the angular speed when the rod reaches vertical, and the linear
acceleration of the free tip at release.

\textbf{Angular acceleration at release, by the moment equation.} About
the fixed pin, the only moment is the weight acting at the rod's centre
of mass, a distance $L/2$ from the pin. With the rod horizontal the
moment arm is the full $L/2$, so the moment is $M=mg(L/2)$. The moment of
inertia of a uniform rod about one end is $I_{\text{end}}=\tfrac{1}{3}m
L^{2}$. The moment equation about the fixed axis gives
\[
\alpha=\frac{M}{I_{\text{end}}}=\frac{mg(L/2)}{\tfrac{1}{3}mL^{2}}
=\frac{3g}{2L}.
\]
The tip, at radius $L$, has tangential acceleration $a_{\text{tip}}=
\alpha L=\tfrac{3}{2}g$, which exceeds $g$. A coin balanced on the very
tip of a falling rod is left behind: the tip accelerates downward faster
than free fall, so the coin, which can only manage $g$, separates from
the rod and lands behind it. The point on the rod that falls at exactly
$g$ is the one at radius $r$ where $\alpha r=g$, namely $r=2L/3$;
everything beyond it outruns gravity.

\textbf{Angular speed at vertical, by energy.} The pin does no work, so
mechanical energy is conserved from release to the vertical position. The
centre of mass drops by $L/2$ as the rod swings down, releasing
potential energy $mg(L/2)$, which becomes rotational kinetic energy
$\tfrac{1}{2}I_{\text{end}}\omega^{2}$:
\[
mg\frac{L}{2}=\tfrac{1}{2}\Bigl(\tfrac{1}{3}mL^{2}\Bigr)\omega^{2}
\ \Rightarrow\
\omega=\sqrt{\frac{3g}{L}}.
\]
For a rod of length $L=1.0\,\si{\meter}$, $\omega=\sqrt{3(9.81)/1.0}=
5.42\,\si{\radian\per\second}$ at the bottom, and the tip there moves at
$\omega L=5.42\,\si{\meter\per\second}$. Using the end moment of inertia,
rather than the centre value $\tfrac{1}{12}mL^{2}$, is the step a careless
solution gets wrong: the rotation is about the pin, not the centre, and
the parallel-axis theorem $I_{\text{end}}=\tfrac{1}{12}mL^{2}+m(L/2)^{2}=
\tfrac{1}{3}mL^{2}$ reconciles the two.

\begin{mastery}
\textbf{The centre of percussion: where to strike so the pivot feels
nothing.} The falling rod pivots at a pin that must supply whatever
horizontal reaction the motion demands, and a companion question decides
where a blow can be struck so that the pin feels no jolt at all. Strike a
pivoted body a sharp transverse blow of impulse $J$ at a distance $b$
from the pivot. The blow does two things at once: it drives the centre of
mass, changing the linear momentum by $J$, and it torques the body about
the pivot, changing the angular momentum. Whether the pivot feels a
reaction impulse depends on whether these two demands are consistent with
a body that rotates about the pivot alone. Write the two impulse
balances. The linear impulse-momentum theorem about the centre of mass at
distance $d$ from the pivot gives $J+J_{\text{pin}}=m\,d\,\omega$, where
$\omega$ is the angular speed the blow imparts and $J_{\text{pin}}$ is the
reaction impulse at the pin. The angular impulse-momentum theorem about
the pivot gives $J\,b=I_{\text{pin}}\omega$, with $I_{\text{pin}}$ the
moment of inertia about the pivot. Eliminating $\omega$ between the two
and setting the pin reaction $J_{\text{pin}}$ to zero gives the striking
distance that leaves the pivot unloaded,
\[
b=\frac{I_{\text{pin}}}{m\,d}=\frac{k_{\text{pin}}^{2}}{d},
\]
the \engterm{centre of percussion} conjugate to the pivot, with
$k_{\text{pin}}$ the radius of gyration about the pivot. A blow at that
point drives the body's rotation about the pivot exactly, so the pivot
neither lurches forward nor snaps back. For a uniform rod pivoted at one
end, $I_{\text{pin}}=\tfrac{1}{3}mL^{2}$ and $d=L/2$, so $b=\tfrac{2}{3}L$:
striking two-thirds of the way down leaves the pivot quiet. This is the
sweet spot of a bat, a racket, and a hammer, the point a batter learns to
meet the ball on so the handle does not sting the hands, and it is why a
door stop is set near two-thirds of the door's width from the hinge so
the hinge pins are spared the slam. The same reasoning, run the other
way, tells a designer where a latching impact will punish a bearing: a
strike anywhere but the centre of percussion loads the pivot, and the
farther from that point, the harder the jolt the mounting must carry.
\end{mastery}

\subsection{An Atwood machine with a massive pulley}

The ideal pulley of the earlier worked example carried no inertia of its
own; a real pulley stores rotational energy and changes the answer, and
working the correction shows how a rotating constraint element enters the
force balance. Two blocks of masses $m_{1}=5\,\si{\kilogram}$ and
$m_{2}=3\,\si{\kilogram}$ hang from a cord over a pulley of mass $M=2\,
\si{\kilogram}$ and radius $R$, modelled as a uniform disc with
$I=\tfrac{1}{2}MR^{2}$. The cord does not slip on the pulley. Find the
acceleration of the blocks and the two cord tensions.

\textbf{Three free bodies, one constraint.} The cord no longer carries a
single tension: the pulley's rotational inertia means the tension on the
heavy side exceeds the tension on the light side, and that difference is
what angularly accelerates the pulley. The no-slip condition ties the
linear acceleration of the blocks to the pulley's angular acceleration,
$a=\alpha R$. Write the three equations of motion, taking the heavier
block to descend:
\[
m_{1}g-T_{1}=m_{1}a, \qquad
T_{2}-m_{2}g=m_{2}a, \qquad
(T_{1}-T_{2})R=I\alpha=\tfrac{1}{2}MR^{2}\frac{a}{R}.
\]
The third equation simplifies to $T_{1}-T_{2}=\tfrac{1}{2}Ma$. Add the
first two to get $(m_{1}-m_{2})g-(T_{1}-T_{2})=(m_{1}+m_{2})a$, then
substitute the pulley relation:
\[
a=\frac{(m_{1}-m_{2})g}{m_{1}+m_{2}+\tfrac{1}{2}M}.
\]
The pulley adds half its mass to the system inertia, the signature of a
uniform disc whose rim moves at the cord speed. With the numbers,
$a=(5-3)(9.81)/(5+3+1)=2(9.81)/9=2.18\,\si{\meter\per\second\squared}$,
against the $2.45\,\si{\meter\per\second\squared}$ an ideal massless
pulley would give; the heavier pulley slows the system by about
$11\,\%$. The tensions follow from the block equations, $T_{1}=m_{1}(g-a)
=5(9.81-2.18)=38.1\,\si{\newton}$ and $T_{2}=m_{2}(g+a)=3(9.81+2.18)=
36.0\,\si{\newton}$, and their difference $2.1\,\si{\newton}$ times $R$
supplies exactly the pulley's $\tfrac{1}{2}MaR$, the consistency check.

\subsection{An exploding shell and the centre of mass}

A projectile that bursts in flight tests the centre-of-mass theorem on a
problem where tracking every fragment would be hopeless but the centre
of mass is trivial. A shell of mass $m=12\,\si{\kilogram}$ is fired and,
at the top of its trajectory, is moving horizontally at $v=80\,\si{\meter
\per\second}$ when it bursts into two fragments. One fragment, of mass
$m_{1}=8\,\si{\kilogram}$, is thrown straight down with a speed of
$30\,\si{\meter\per\second}$ relative to the ground at the instant of the
burst. Find the velocity of the second fragment immediately after the
burst, and state where the centre of mass lands relative to the
unexploded trajectory.

\textbf{Momentum across the burst.} The explosion is internal, so the
total momentum is conserved through the instant of the burst; gravity
acts during the burst too, but over the vanishingly short burst time its
impulse is negligible against the explosive impulse. Before the burst the
momentum is all horizontal, $\vec{p}=m v\,\hat{\imath}=12(80)\hat{\imath}
=960\,\hat{\imath}\,\si{\kilogram\meter\per\second}$. The first fragment
carries $\vec{p}_{1}=m_{1}(0,-30)=(0,-240)\,\si{\kilogram\meter\per
\second}$. The second fragment ($m_{2}=4\,\si{\kilogram}$) carries the
balance:
\[
\vec{p}_{2}=\vec{p}-\vec{p}_{1}=(960,0)-(0,-240)=(960,240)\,\si{\kilogram
\meter\per\second}.
\]
Its velocity is $\vec{v}_{2}=\vec{p}_{2}/m_{2}=(240,60)\,\si{\meter\per
\second}$, a speed of $\sqrt{240^{2}+60^{2}}=247\,\si{\meter\per\second}$
directed $\arctan(60/240)=14.0^{\circ}$ above horizontal and forward. The
light fragment is flung forward and up at three times the original speed,
the recoil of throwing the heavy fragment down.

\textbf{Where the centre of mass goes.} The burst changes no external
force, so the centre of mass continues on the original parabola as
though nothing happened, until the first fragment strikes the ground and
breaks the system. Up to that moment the two fragments straddle the
undisturbed trajectory: the horizontal centre-of-mass speed stays
$80\,\si{\meter\per\second}$, since $(m_{1}\cdot 0+m_{2}\cdot 240)/12=
960/12=80\,\si{\meter\per\second}$ recovers it, and the vertical
centre-of-mass velocity stays zero at the burst instant, since $(8(-30)+
4(60))/12=0$. The arithmetic is the centre-of-mass theorem made
concrete: the fragments' momenta are whatever the explosion dictates,
but their mass-weighted average is fixed by the external force alone.

\subsection{Apparent weight: a load cell in a lift}

\engterm{Vignette.} A load cell is a force sensor: it reads the contact
force pressing on it, not the mass resting on it, and the distinction
becomes a measurement when the sensor accelerates. A
$70\,\si{\kilogram}$ technician stands on a load-cell platform bolted to
the floor of a service lift, watching the readout as the lift starts,
runs, and stops. The reading is the normal force $N$ the platform pushes
up with, and Newton's second law applied to the technician along the
vertical sets it (Figure~\ref{fig:vol03ch04:elevator-scale}).

\textbf{The one equation, read three ways.} Taking up as positive, the
forces on the technician are the platform's push $N$ up and the weight
$mg$ down, and the net equals the mass times the vertical acceleration:
\[
N-mg=ma \ \Rightarrow\ N=m(g+a).
\]
At rest or at constant speed, $a=0$ and the cell reads the true weight
$N=mg=70(9.81)=687\,\si{\newton}$. As the lift accelerates upward at, say,
$a=1.5\,\si{\meter\per\second\squared}$ to get moving, the reading rises
to $N=70(9.81+1.5)=792\,\si{\newton}$, a $15\,\%$ overshoot, and the
technician feels pressed into the floor. As the lift brakes to a stop at
the top of its run, the acceleration reverses to $a=-1.5\,\si{\meter\per
\second\squared}$ and the cell reads $N=70(9.81-1.5)=582\,\si{\newton}$,
the light-footed feeling of a lift slowing. The cabin's velocity does not
enter; only its acceleration does, which is why a smooth high-speed lift
feels weightless-normal at cruise and odd only at the start and the
stop.

\textbf{The measurement consequence.} A scale that must report mass, not
contact force, has to know it is not accelerating, or correct for the
acceleration it senses. A bathroom scale trusts the user to stand still;
a laboratory balance is levelled and isolated from floor vibration; a
truck weighbridge takes its reading only when the vehicle is stationary.
The free-fall limit makes the point sharpest: at $a=-g$ the reading is
$N=0$, the cell registering nothing because in free fall the platform and
the technician fall together and press on each other not at all. An
astronaut in orbit has not lost weight in the sense of $mg$; the
spacecraft and everything in it are in continuous free fall, so every
contact force, and every scale reading, is zero. The load cell measures
the contact force faithfully; reading mass off it is the engineer's
inference, valid only when the acceleration is known.

% Figure: apparent weight in an accelerating elevator. Three cases,
% drawn as a person on a scale: at rest (or constant speed), accelerating
% up, and accelerating down. The scale reads the normal force N, which
% equals m(g+a) up and m(g-a) down.
\begin{figure}[htbp]
\centering
\begin{tikzpicture}[
  force/.style={-{Stealth[length=2.6mm]}, very thick},
  lbl/.style={font=\small},
  cab/.style={draw=enggray, thick},
]
\foreach \dx/\cap/\acc/\Nlab/\col in {
  0/{rest}/{}/{$N=mg$}/engblue,
  4.4/{up}/{$a$}/{$N=m(g+a)$}/engred,
  8.8/{down}/{$a$}/{$N=m(g-a)$}/englime!65!black}
{
  % cab outline
  \draw[cab] (\dx,0) rectangle ++(2.6,3.6);
  % scale platform
  \draw[engblack, very thick] (\dx+0.5,0.5) -- (\dx+2.1,0.5);
  % body (block standing on scale)
  \fill[\col!25] (\dx+0.95,0.5) rectangle ++(0.7,1.6);
  \draw[\col, thick] (\dx+0.95,0.5) rectangle ++(0.7,1.6);
  % weight down (from body centre)
  \draw[force, draw=engblack] (\dx+1.3,1.3) -- ++(0,-1.0);
  \node[lbl, right] at (\dx+1.35,0.85) {$mg$};
  % normal force up (from scale into feet)
  \draw[force, draw=\col] (\dx+1.3,0.55) -- ++(0,1.0);
  \node[\col, lbl, left] at (\dx+0.6,1.05) {$N$};
  % heading
  \node[lbl, above] at (\dx+1.3,3.6) {\cap};
  % scale reading
  \node[\col, lbl, below] at (\dx+1.3,-0.05) {\Nlab};
}
% acceleration arrow for case 2 (up)
\draw[force, draw=engred, dashed] (4.4-0.55,1.4) -- ++(0,1.2);
\node[engred, lbl] at (4.4-0.85,2.0) {$a$};
% acceleration arrow for case 3 (down)
\draw[force, draw=englime!65!black, dashed] (8.8-0.55,2.6) -- ++(0,-1.2);
\node[englime!65!black, lbl] at (8.8-0.85,2.0) {$a$};
\end{tikzpicture}
\caption[Apparent weight in an accelerating elevator]{A person of mass
$m$ on a scale in an elevator. The scale reads the normal force $N$ it
pushes up with, which by Newton's second law applied to the person along
the vertical equals the static weight only when the cab is at rest or
moving at constant speed. Accelerating upward at $a$ raises the reading
to $N=m(g+a)$, the floating-heavy sensation of a lift starting up;
accelerating downward lowers it to $N=m(g-a)$, the light feeling of a
lift starting down, which reaches free fall and a zero reading at $a=g$.
The weight $mg$ never changes; the apparent weight the scale reports is
the contact force, and the contact force tracks the acceleration.}
\label{fig:vol03ch04:elevator-scale}
\end{figure}


% Figure: a conical pendulum. A bob on a string of length L sweeps a
% horizontal circle of radius r at constant speed; the string traces a
% cone of half-angle theta. The free-body diagram resolves the string
% tension into a vertical component balancing gravity and a horizontal
% component supplying the centripetal force.
\begin{figure}[htbp]
\centering
\begin{tikzpicture}[
  force/.style={-{Stealth[length=3mm]}, very thick},
  lbl/.style={font=\small},
]
% suspension point
\coordinate (O) at (0,3.4);
\fill[engblack] (O) circle (1.6pt);
\node[lbl, above] at (O) {pivot};
% vertical axis (dashed) down from pivot
\draw[enggray, dashed] (O) -- (0,0);
% bob position: string makes angle theta from vertical; take theta with
% sin=0.5547, cos=0.8321 (3-4-5-like: r=1.8, vertical drop=2.7, L=3.245)
\coordinate (B) at (1.8,0.7);
% string
\draw[engblue, very thick] (O) -- (B);
\node[engblue, lbl] at (0.75,2.2) {$L$};
% the cone half-angle theta at the pivot
\draw[enggray] (0,2.65) arc (-90:-56:0.75);
\node[lbl] at (0.35,2.55) {$\theta$};
% bob
\fill[engred] (B) circle (3.5pt);
\node[engred, lbl, right] at (B) {bob $m$};
% horizontal circle (perspective ellipse) the bob sweeps, radius r=1.8
\draw[englime!70!black, thick] (0,0.7) ellipse (1.8 and 0.42);
\node[englime!60!black, lbl] at (-1.3,0.25) {radius $r$};
% radius line from axis to bob at the bob's height
\draw[enggray, dashed] (0,0.7) -- (B);
\node[lbl, below] at (0.9,0.7) {$r$};
% forces at the bob:
% weight down
\draw[force, draw=engblack] (B) -- ++(0,-1.1);
\node[lbl, below] at (2.0,-0.5) {$mg$};
% tension along string toward pivot
\draw[force, draw=engamber] (B) -- ($(B)!0.42!(O)$);
\node[engamber, lbl, right] at (2.05,1.4) {$T$};
\end{tikzpicture}
\caption[Conical pendulum]{A conical pendulum: a bob on a string of
length $L$ sweeps a horizontal circle of radius $r=L\sin\theta$ at
constant speed, so the string traces a cone of half-angle $\theta$.
The bob has no vertical acceleration, so the vertical component of the
string tension balances gravity, $T\cos\theta=mg$; the horizontal
component supplies the centripetal force, $T\sin\theta=m v^2/r$.
Dividing the two gives $\tan\theta=v^2/(gr)$, and the period works out
to $T_{\text{period}}=2\pi\sqrt{L\cos\theta/g}$, which depends on the
cone angle only weakly at small angles and shrinks as the cone opens.}
\label{fig:vol03ch04:conical-pendulum}
\end{figure}


% Figure: the Coriolis deflection seen on a rotating turntable. A puck
% slid radially outward along a straight line in the inertial frame
% appears, to an observer co-rotating with the turntable, to curve to
% one side. The two paths are drawn: the straight inertial path and the
% curved path painted on the turntable.
\begin{figure}[htbp]
\centering
\begin{tikzpicture}[lbl/.style={font=\small}]
% turntable disc
\draw[enggray, very thick] (0,0) circle (3.2);
\fill[enggray!12] (0,0) circle (3.2);
\fill[engblack] (0,0) circle (2pt);
\node[lbl, below left] at (0,0) {centre};
% rotation arrow (counter-clockwise)
\draw[-{Stealth[length=2.5mm]}, engblue, thick] (135:3.5) arc (135:95:3.5);
\node[engblue, lbl] at (118:3.85) {$\Omega$};
% inertial straight path: launched at center toward angle 0 (to the right)
\draw[engred, very thick, -{Stealth[length=3mm]}] (0,0) -- (3.0,0);
\node[engred, lbl, above] at (2.2,0.05) {inertial path (straight)};
% curved path as seen on the rotating disc: deflects to the right of
% travel for counter-clockwise rotation. Plot a curve bending downward.
\draw[englime!60!black, very thick, -{Stealth[length=3mm]}]
  (0,0) .. controls (1.4,-0.35) and (2.4,-1.1) .. (2.7,-1.55);
\node[englime!50!black, lbl, below] at (2.0,-1.4) {path on the disc (curved)};
% small deflection bracket
\draw[enggray, dashed] (3.0,0) -- (2.7,-1.55);
\node[enggray, lbl, right] at (3.0,-0.8) {deflection};
\end{tikzpicture}
\caption[Coriolis deflection on a rotating turntable]{A puck launched
from the centre of a counter-clockwise turntable travels in a straight
line in the inertial frame (the red arrow), because once it leaves the
centre no horizontal force acts on it. An observer painting the puck's
track on the rotating disc records a curved path (the green arrow) that
bends to the right of the direction of travel. The curvature is the
Coriolis effect: in the rotating frame the puck obeys
$m\vec{a}'=-2m\vec{\Omega}\times\vec{v}'$, a deflecting pseudo-force
perpendicular to the puck's velocity. The same deflection sets the
rotation sense of large-scale winds and biases any projectile that
flies far enough for the planet to turn beneath it.}
\label{fig:vol03ch04:coriolis}
\end{figure}


\subsection{A block sliding on a moving block}

Two bodies in contact that may or may not slide relative to each other
force the analyst to test the no-slip assumption rather than assume it,
and a block riding on a pulled block is the cleanest setting for that
test. A block of mass $m_{2}=2.0\,\si{\kilogram}$ rests on top of a
block of mass $m_{1}=6.0\,\si{\kilogram}$ that sits on a frictionless
floor (Figure~\ref{fig:vol03ch04:block-on-block}). The contact between
the two blocks has static coefficient $\mu_{s}=0.30$ and kinetic
$\mu_{k}=0.25$. A horizontal force $P$ is applied to the lower block.
Find the largest $P$ for which the two blocks move together, and the
acceleration of each block when $P$ exceeds that value.

\textbf{Move-together phase.} While the blocks do not slip they share a
common acceleration $a=P/(m_{1}+m_{2})$, because the applied force
accelerates the whole $8.0\,\si{\kilogram}$ stack. The only horizontal
force on the upper block is the friction $f$ from the lower block's top
face, so that friction alone must produce the upper block's share of the
acceleration:
\[
f=m_{2}a=m_{2}\frac{P}{m_{1}+m_{2}}.
\]
The friction can rise no higher than $\mu_{s}m_{2}g$, so the blocks move
together only while
\[
m_{2}\frac{P}{m_{1}+m_{2}}\le\mu_{s}m_{2}g
\ \Rightarrow\
P\le\mu_{s}(m_{1}+m_{2})g=0.30(8.0)(9.81)=23.5\,\si{\newton}.
\]
The upper block's mass drops out of the inequality: the slip threshold
depends on the total mass and the friction coefficient, not on how the
mass is split between the two blocks. At the threshold the shared
acceleration is $a=\mu_{s}g=2.94\,\si{\meter\per\second\squared}$.

\textbf{Slipping phase.} For $P>23.5\,\si{\newton}$ the top block cannot
keep up; the contact slides and the friction saturates at its kinetic
value $f=\mu_{k}m_{2}g=0.25(2.0)(9.81)=4.91\,\si{\newton}$, forward on
the upper block and backward on the lower. The two blocks now accelerate
separately:
\[
a_{2}=\frac{\mu_{k}m_{2}g}{m_{2}}=\mu_{k}g=2.45\,\si{\meter\per\second
\squared},\qquad
a_{1}=\frac{P-\mu_{k}m_{2}g}{m_{1}}.
\]
The upper block's acceleration is frozen at $\mu_{k}g$, whatever $P$
does, because the only forward force it feels is the saturated friction.
The lower block accelerates faster and slides out from under it. For a
pull of $P=40\,\si{\newton}$, $a_{1}=(40-4.91)/6.0=5.85\,\si{\meter\per
\second\squared}$ against the upper block's $2.45\,\si{\meter\per\second
\squared}$: the lower block runs ahead, and the top block, lagging,
slides toward the rear. The discipline the problem teaches is to solve
the coupled case first, check the friction it demands against the limit,
and switch to the decoupled case only when the check fails.

% Figure: a block resting on a block, the lower one pulled. Free-body
% diagrams of both blocks side by side, showing the friction that must
% accelerate the upper block and the limit at which it slips.
\begin{figure}[htbp]
\centering
\begin{tikzpicture}[
  force/.style={draw=engblue, -{Stealth}, thick},
  wt/.style={draw=engblack, -{Stealth}, thick},
  fr/.style={draw=engamber, -{Stealth}, thick},
  nrm/.style={draw=engred, -{Stealth}, thick},
  lbl/.style={font=\scriptsize},
]
% left: the physical picture
\begin{scope}
  \draw[draw=enggray, thick] (-0.3,0) -- (5,0);
  \fill[enggray!12] (-0.3,0) rectangle (5,-0.35);
  % lower block m1
  \draw[draw=engblack, thick, fill=engblue!8] (0.4,0) rectangle (3.4,1.0);
  \node[lbl] at (1.9,0.5) {$m_1$};
  % upper block m2
  \draw[draw=engblack, thick, fill=englime!18] (1.2,1.0) rectangle (2.9,1.8);
  \node[lbl] at (2.05,1.4) {$m_2$};
  % applied pull on lower block
  \draw[force] (3.4,0.5) -- ++(1.3,0) node[right, lbl, engblue] {$P$};
  \node[lbl, anchor=north] at (2.3,-0.15) {applied to lower block};
\end{scope}
% right: stacked free-body diagrams
\begin{scope}[shift={(7.6,0)}]
  % upper block FBD
  \draw[draw=engblack, thick, fill=englime!18] (-0.7,1.6) rectangle (0.7,2.4);
  \node[lbl] at (0,2.0) {$m_2$};
  \draw[wt] (0,1.6) -- ++(0,-0.9) node[below, lbl] {$m_2 g$};
  \draw[nrm] (0,2.4) -- ++(0,0.9) node[above, lbl, engred] {$N_2$};
  \draw[fr] (0,2.0) -- ++(1.1,0) node[right, lbl, engamber] {$f$ (drives $m_2$)};
  % lower block FBD
  \draw[draw=engblack, thick, fill=engblue!8] (-1.0,-0.7) rectangle (1.0,0.3);
  \node[lbl] at (-0.4,-0.2) {$m_1$};
  \draw[wt] (-0.4,-0.7) -- ++(0,-0.8) node[below, lbl] {$m_1 g$};
  \draw[nrm] (-0.4,0.3) -- ++(0,0.8) node[above, lbl, engred] {$N_1$};
  \draw[force] (1.0,-0.2) -- ++(1.1,0) node[right, lbl, engblue] {$P$};
  \draw[fr] (0.4,0.3) -- ++(-1.1,0) node[left, lbl, engamber] {$f$ (reaction)};
  \draw[nrm] (0.4,0.3) -- ++(0,-0.7) node[pos=1.0, below, lbl, engred] {$N_2$};
\end{scope}
\end{tikzpicture}
\caption[Two stacked blocks and their free-body diagrams]{A block
$m_2$ rides on a block $m_1$ that is pulled by a force $P$ along a
frictionless floor. The only horizontal force on the upper block is the
friction $f$ from the lower block's top face, so that friction is what
must accelerate $m_2$; it can grow no larger than $\mu_s m_2 g$. While
$P$ is small the two blocks share one acceleration and the friction
stays below its limit; when $P$ is large enough that the shared
acceleration would need more friction than $\mu_s m_2 g$, the upper
block slips and the two accelerations part. The paired diagrams keep
the third-law friction couple straight: the lower block feels the equal
and opposite reaction $f$ backward on its top face.}
\label{fig:vol03ch04:block-on-block}
\end{figure}


\subsection{Terminal velocity under quadratic drag}

A body falling through air soon meets a drag that grows with speed until
it balances the weight, and the motion settles to a constant terminal
speed the constant-acceleration formulas never reach. A skydiver of
mass $m=80\,\si{\kilogram}$ falls through air of density $\rho=1.2\,\si{
\kilogram\per\meter\cubed}$ with a drag area $C_{d}A=0.70\,\si{\meter
\squared}$ in the belly-to-earth posture. Find the terminal speed, and
the speed reached after a fall of $100\,\si{\meter}$ from rest.

\textbf{The terminal speed is a force balance.} At terminal speed the
acceleration is zero, so the quadratic drag equals the weight:
\[
\tfrac{1}{2}\rho C_{d}A\,v_{t}^{2}=mg
\ \Rightarrow\
v_{t}=\sqrt{\frac{2mg}{\rho C_{d}A}}=\sqrt{\frac{2(80)(9.81)}{1.2(0.70)}}
=43.2\,\si{\meter\per\second},
\]
about $156\,\si{\kilo\meter\per\hour}$, the familiar figure for a
belly-down freefall. A diver who pulls into a head-down dart cuts
$C_{d}A$ by roughly a factor of four and raises $v_{t}$ by $\sqrt{4}=2$,
to over $300\,\si{\kilo\meter\per\hour}$, the physics behind
speed-skydiving.

\textbf{The approach to terminal speed.} The equation of motion is
$m\dot{v}=mg-\tfrac{1}{2}\rho C_{d}A\,v^{2}$, which in terms of $v_{t}$
and the drag parameter $k=\tfrac{1}{2}\rho C_{d}A/m=g/v_{t}^{2}$ reads
$\dot{v}=g(1-v^{2}/v_{t}^{2})$. Its closed solution from rest is
$v(t)=v_{t}\tanh(gt/v_{t})$, and the speed as a function of distance
fallen, obtained by writing $\dot{v}=v\,dv/dx$ and integrating, is
\[
v(x)=v_{t}\sqrt{1-\exp\!\left(-\frac{2gx}{v_{t}^{2}}\right)}.
\]
At $x=100\,\si{\meter}$ the exponent is $-2(9.81)(100)/43.2^{2}=-1.052$,
so $v=43.2\sqrt{1-e^{-1.052}}=43.2\sqrt{1-0.349}=34.8\,\si{\meter\per
\second}$, already $81\,\%$ of terminal. A drag-free estimate would have
given $v=\sqrt{2gx}=\sqrt{2(9.81)(100)}=44.3\,\si{\meter\per\second}$,
larger than the terminal speed itself, which is the tell that the
drag-free model is nonsense here: it lets the speed grow without bound
while the real diver is already nearly capped. The characteristic length
over which drag asserts itself is $v_{t}^{2}/(2g)=95\,\si{\meter}$, so a
fall of $100\,\si{\meter}$ is about one such length and the two models
have already parted by a quarter.

\subsection{A rocket's vertical ascent against gravity}

The rocket derivation of the exercises stripped away gravity to isolate
the mass-ratio result; a real launch must lift against gravity for the
whole burn, and carrying that term shows how thrust-to-weight and burn
time set the burnout speed. A sounding rocket of initial mass
$m_{0}=1{,}200\,\si{\kilogram}$ burns propellant at $\dot{m}=40\,\si{
\kilogram\per\second}$ with exhaust speed $v_{e}=2{,}400\,\si{\meter\per
\second}$, rising vertically from rest. The burn lasts until half the
initial mass is gone. Find the thrust, the initial and final
thrust-to-weight ratios, and the burnout speed, treating $g$ as constant
and neglecting air drag.

\textbf{Thrust and thrust-to-weight.} The thrust is the momentum the
exhaust carries away per unit time, $F=\dot{m}v_{e}=40(2{,}400)=96\,\si{
\kilo\newton}$, constant through the burn. The initial weight is
$m_{0}g=1{,}200(9.81)=11.8\,\si{\kilo\newton}$, so the lift-off
thrust-to-weight ratio is $F/(m_{0}g)=96/11.8=8.1$, a brisk launch. As
propellant burns the vehicle lightens, so the ratio climbs; at burnout
the mass is $m_{f}=600\,\si{\kilogram}$, the weight $5.89\,\si{\kilo
\newton}$, and the ratio $96/5.89=16.3$. The rising ratio is why a
launch vehicle's acceleration grows through a stage even at constant
thrust, and why the crew feels the heaviest press just before staging.

\textbf{Burnout speed with the gravity term.} The vertical equation of
motion is the rocket equation with a gravity loss:
\[
m\frac{dv}{dt}=\dot{m}v_{e}-mg.
\]
Integrating over the burn adds a gravity-loss term $-g\,t_{b}$ to the
ideal velocity, where $t_{b}$ is the burn time. The burn time is the
propellant burned over the burn rate, $t_{b}=(m_{0}-m_{f})/\dot{m}=
600/40=15\,\si{\second}$. The ideal velocity is the mass-ratio result
$v_{e}\ln(m_{0}/m_{f})=2{,}400\ln 2=1{,}663\,\si{\meter\per\second}$, and
the gravity loss is $g\,t_{b}=9.81(15)=147\,\si{\meter\per\second}$, so
the burnout speed is
\[
v_{b}=v_{e}\ln\frac{m_{0}}{m_{f}}-g\,t_{b}=1{,}663-147=1{,}516\,\si{
\meter\per\second}.
\]
The gravity loss costs about $9\,\%$ of the ideal velocity here, and it
would cost more for a sluggish launch: a vehicle with a lift-off
thrust-to-weight nearer one would spend far longer fighting gravity for
the same $\Delta v$, which is the quantitative reason a launcher runs a
high thrust-to-weight and climbs quickly through the vertical part of
its ascent, the same loss the two-stage launch case study budgeted at
the vehicle scale.

\subsection{A falling spool: rolling on its own string}

A spool that unwinds as it falls couples a translation and a rotation
through the string rather than through a road, and it pays back a clean
lesson in how a constraint at a radius other than the outer edge changes
the answer. A uniform spool of mass $m$, outer radius $R$, and radius of
gyration $k$ (so $I_{G}=mk^{2}$) hangs from a string wrapped around an
inner hub of radius $r$ and is released from rest. The string is
vertical and fixed at the top. Find the downward acceleration of the
centre and the tension in the string.

\textbf{Two equations and the string constraint.} The forces on the
spool are its weight $mg$ down and the string tension $T$ up, applied at
the hub radius $r$. Newton's second law for the centre and the moment
equation about the centre are
\[
mg-T=ma_{G}, \qquad Tr=I_{G}\alpha=mk^{2}\alpha.
\]
The string is fixed and unwinds without slipping, so the centre drops at
the rate the string pays out from the hub, $a_{G}=\alpha r$. Substitute
$\alpha=a_{G}/r$ into the moment equation to get $T=mk^{2}a_{G}/r^{2}$,
and put that into the force equation:
\[
mg-\frac{mk^{2}a_{G}}{r^{2}}=ma_{G}
\ \Rightarrow\
a_{G}=\frac{g}{1+k^{2}/r^{2}}.
\]
The acceleration is the free-fall value divided by one plus the ratio of
the gyration radius squared to the hub radius squared. A spool wound on a
small hub ($r\ll k$) falls slowly, most of the released potential energy
going into spin; a spool with the string on the outer rim ($r=R$)
behaves like the rolling cylinder of the incline example with the slope
turned vertical. For a uniform disc, $k^{2}=R^{2}/2$; if the string is
on the rim ($r=R$), $a_{G}=g/(1+1/2)=\tfrac{2}{3}g$, matching the
free-rolling cylinder, and if the string is on a hub of half the radius
($r=R/2$), $a_{G}=g/(1+2)=\tfrac{1}{3}g$, a markedly slower fall. The
tension is $T=m(g-a_{G})=mg\,(k^{2}/r^{2})/(1+k^{2}/r^{2})$, which for
the half-radius hub is $\tfrac{2}{3}mg$, a string carrying two-thirds of
the weight while the spool descends at a third of $g$.

\subsection{Two spinning discs coupled by a clutch}

When two rotating bodies are suddenly locked together, angular momentum
survives the engagement while rotational energy does not, the exact
rotational parallel of the inelastic collision, and a clutch is the
device that makes the parallel concrete. A disc of moment of inertia
$I_{1}=0.40\,\si{\kilogram\meter\squared}$ spins at $\omega_{1}=120\,\si{
\radian\per\second}$; a second, initially stationary disc of
$I_{2}=0.60\,\si{\kilogram\meter\squared}$ is pressed against it by a
clutch until the two turn as one. Find the common final angular speed
and the fraction of the rotational kinetic energy lost in the
engagement.

\textbf{Angular momentum is conserved across the engagement.} The clutch
forces are internal to the two-disc system and exert equal and opposite
torques on the two discs, so the total angular momentum about the common
axis is the same before and after:
\[
I_{1}\omega_{1}=(I_{1}+I_{2})\omega_{f}
\ \Rightarrow\
\omega_{f}=\frac{I_{1}\omega_{1}}{I_{1}+I_{2}}=\frac{0.40(120)}{1.00}=
48\,\si{\radian\per\second}.
\]
The result is the rotational image of the perfectly inelastic collision,
with moment of inertia in the place of mass and angular speed in the
place of velocity.

\textbf{The energy lost to the slipping clutch.} The kinetic energy
before is $\tfrac{1}{2}I_{1}\omega_{1}^{2}=\tfrac{1}{2}(0.40)(120)^{2}=
2{,}880\,\si{\joule}$; after, $\tfrac{1}{2}(I_{1}+I_{2})\omega_{f}^{2}=
\tfrac{1}{2}(1.00)(48)^{2}=1{,}152\,\si{\joule}$. The loss is
$1{,}728\,\si{\joule}$, a fraction $1{,}728/2{,}880=0.60$ of the
starting energy, dissipated as heat in the slipping clutch faces. The
lost fraction equals $I_{2}/(I_{1}+I_{2})=0.60$, the same form the
inelastic collision gave with the struck mass in the numerator, and it
is why a clutch engaged across a large speed difference runs hot and
wears: the energy the friction faces must absorb is fixed by the inertia
ratio and the speed gap, not by how gently the clutch is let out. Letting
the clutch out slowly spreads the same heat over more time, easing the
peak temperature, but it does not reduce the total energy the engagement
must throw away.

\subsection{A car braking and turning at once: the friction circle in
numbers}

The friction-circle mastery box argued that braking and cornering share
one grip budget; a worked case puts numbers to the trade and shows how
much braking a driver must surrender to hold a bend. A car of mass
$m=1{,}300\,\si{\kilogram}$ rounds a flat curve of radius $r=90\,\si{
\meter}$ at $v=25\,\si{\meter\per\second}$ on a road with tyre-grip
coefficient $\mu=0.85$. Find the lateral grip the corner already spends
and the largest braking deceleration still available without sliding.

\textbf{The lateral demand.} Holding the curve needs a centripetal
acceleration $a_{\text{lat}}=v^{2}/r=625/90=6.94\,\si{\meter\per\second
\squared}$, supplied entirely by tyre friction on the flat road. The
total grip the tyres can muster is $\mu g=0.85(9.81)=8.34\,\si{\meter\per
\second\squared}$. The cornering alone therefore spends a fraction
$a_{\text{lat}}/(\mu g)=6.94/8.34=0.83$ of the available grip, leaving
little in reserve.

\textbf{The braking left over.} The friction circle bounds the vector
sum of the longitudinal and lateral accelerations by $\mu g$:
\[
\sqrt{a_{\text{brake}}^{2}+a_{\text{lat}}^{2}}\le\mu g
\ \Rightarrow\
a_{\text{brake}}\le\sqrt{(\mu g)^{2}-a_{\text{lat}}^{2}}
=\sqrt{8.34^{2}-6.94^{2}}=4.63\,\si{\meter\per\second\squared}.
\]
The car cornering this hard can brake at only $4.63\,\si{\meter\per
\second\squared}$, about $0.47g$, against the $8.34\,\si{\meter\per
\second\squared}$ ($0.85g$) it could manage braking in a straight line.
Asking for the full straight-line braking while holding the corner would
demand $\sqrt{8.34^{2}+6.94^{2}}=10.9\,\si{\meter\per\second\squared}$
from the tyres, well outside the circle, so the tyres would break away
and the car would run wide. The lesson the numbers carry is that grip
spent turning is grip unavailable for stopping, and the driver who
brakes late into a fast bend finds the pedal does far less than it does
on the straight, the everyday face of the friction circle.

\subsection{A particle leaving a smooth sphere}

A body sliding over a curved surface stays on it only while the surface
can push, and finding where it leaves is a clean use of the normal-force
condition together with energy conservation. A small bead released from
rest at the top of a frictionless sphere of radius $R$ slides down the
outside (Figure~\ref{fig:vol03ch04:bead-sphere}). Find the angle from
the vertical at which it leaves the surface.

\textbf{Two relations at angle $\theta$.} While the bead is in contact
at an angle $\theta$ from the top, the inward net of the radial
component of gravity and the normal force supplies the centripetal
acceleration:
\[
mg\cos\theta-N=\frac{mv^{2}}{R}.
\]
Energy conservation from the top, where the bead started at rest and has
dropped a height $R(1-\cos\theta)$, fixes the speed:
\[
\tfrac{1}{2}mv^{2}=mgR(1-\cos\theta)
\ \Rightarrow\
v^{2}=2gR(1-\cos\theta).
\]

\textbf{The leaving condition.} The bead leaves where the surface can no
longer push, that is where $N$ falls to zero. Setting $N=0$ in the
radial equation gives $mg\cos\theta=mv^{2}/R$, and substituting the
energy result for $v^{2}$:
\[
g\cos\theta=\frac{2gR(1-\cos\theta)}{R}=2g(1-\cos\theta)
\ \Rightarrow\
\cos\theta=\frac{2}{3}.
\]
The bead leaves at $\theta=\arccos(2/3)=48.2^{\circ}$ from the vertical,
at a height $\tfrac{2}{3}R$ above the centre, with a speed $v=\sqrt{2gR/
3}$. The result depends on neither the sphere's radius nor the bead's
mass; it is fixed by the balance of gravity's radial pull against the
centripetal demand. The same condition sets where a skier lifts off a
convex crest, where water leaves a rounded weir, and where a fairground
rider on a domed slide parts from the surface, one geometric answer
serving every smooth convex descent.

% Figure: a bead (or particle) sliding from the top of a frictionless
% sphere leaves the surface where the normal force vanishes, at the
% angle whose cosine is 2/3. The geometry and the leaving point.
\begin{figure}[htbp]
\centering
\begin{tikzpicture}[lbl/.style={font=\scriptsize}]
% sphere
\draw[draw=enggray, thick, fill=engblue!5] (0,0) circle (3);
% vertical through centre
\draw[enggray, dashed] (0,-3.3) -- (0,3.6);
\node[lbl, enggray, anchor=south] at (0,3.6) {vertical};
% centre
\fill[engblack] (0,0) circle (1.2pt);
\node[lbl, anchor=north east] at (0,0) {$O$};
% leaving angle: cos(theta)=2/3 -> theta = 48.19 deg
\def\th{48.19}
% radius to the leaving point
\draw[engblack, thick] (0,0) -- (\th:3);
\node[lbl, anchor=south west] at (\th:1.5) {$R$};
% the bead at the leaving point
\fill[engred] (\th:3) circle (2.5pt);
\node[lbl, engred, anchor=west] at (\th:3.15) {bead leaves here};
% the top starting point
\fill[englime] (90:3) circle (2.5pt);
\node[lbl, anchor=south] at (90:3.15) {start (rest)};
% angle arc at centre
\draw[engamber] (90:1.0) arc (90:\th:1.0);
\node[lbl, engamber] at (68:1.35) {$\theta$};
% weight component drawn at the bead
\draw[-{Stealth}, engblack, thick] (\th:3) -- ++(0,-1.4)
  node[below, lbl] {$mg$};
% free-fall parabola sketch after leaving
\draw[engred, dashed, thick]
  ({3*cos(\th)},{3*sin(\th)})
  .. controls ({3*cos(\th)+0.9},{3*sin(\th)-0.7}) ..
  ({3*cos(\th)+1.6},{3*sin(\th)-2.3});
\end{tikzpicture}
\caption[A bead leaving a frictionless sphere]{A bead released from
rest at the top of a smooth sphere slides down and leaves the surface
at the angle $\theta$ measured from the vertical where the normal force
falls to zero. While in contact the inward net of gravity's radial
component and the normal force supplies the centripetal
$mv^{2}/R$; energy conservation fixes the speed at each angle, and the
normal force reaches zero at $\cos\theta=2/3$, that is $\theta\approx
48.2^{\circ}$. Beyond that point the sphere can only push, not pull, so
the bead separates and falls on a free parabola (red dashed). The
leaving angle is independent of the sphere's size and of the bead's
mass; it is set entirely by the balance of gravity and the centripetal
demand.}
\label{fig:vol03ch04:bead-sphere}
\end{figure}


\subsection{A car coasting to rest: rolling resistance and drag}

A vehicle taken out of gear on the level rolls to a stop under two
resistances a driver rarely separates, and the coast-down is the
standard test that pulls them apart. A car of mass $m=1{,}400\,\si{
\kilogram}$ coasts from $v_{0}=30\,\si{\meter\per\second}$ on a flat
road with the clutch disengaged. The resisting force is the sum of a
speed-independent rolling resistance $R_{0}=c_{r}mg$ with $c_{r}=0.012$
and an aerodynamic drag $\tfrac{1}{2}\rho C_{d}A\,v^{2}$ with
$\tfrac{1}{2}\rho C_{d}A=0.35\,\si{\newton\second\squared\per\meter
\squared}$. The task is the speed at which the two resistances are
equal, and the distance to slow from $30$ to $20\,\si{\meter\per
\second}$, worked by energy rather than by solving the equation of
motion in time.

\textbf{Where the two resistances cross.} Rolling resistance is
$R_{0}=0.012(1400)(9.81)=165\,\si{\newton}$, flat with speed. The drag
grows as $v^{2}$ and matches the rolling term where $\tfrac{1}{2}\rho
C_{d}A\,v^{2}=R_{0}$, that is $v=\sqrt{R_{0}/(0.35)}=\sqrt{165/0.35}=
21.7\,\si{\meter\per\second}$, about $78\,\si{\kilo\meter\per\hour}$.
Below this speed the roll dominates the coast and above it the drag
does, which is why highway coasting bleeds speed fast at first and then
drags on at low speed: the resisting force falls as the car slows.

\textbf{The distance by the work-energy theorem.} The resisting force
varies with speed, so the distance is not $v_{0}^{2}/(2a)$ with a single
$a$. Write the work-energy theorem in the form $m\,v\,dv/dx=-R(v)$ and
separate,
\[
dx=-\frac{m\,v\,dv}{R_{0}+bv^{2}}, \qquad b=\tfrac{1}{2}\rho C_{d}A,
\]
which integrates to
\[
x=\frac{m}{2b}\ln\!\frac{R_{0}+bv_{0}^{2}}{R_{0}+bv_{1}^{2}}.
\]
The integrand is a standard $v\,dv/(R_{0}+bv^{2})$ whose antiderivative
is a logarithm, the same structure the quadratic-drag terminal-velocity
example met. With $v_{0}=30$ and $v_{1}=20\,\si{\meter\per\second}$,
$R_{0}+bv_{0}^{2}=165+0.35(900)=480\,\si{\newton}$ and $R_{0}+bv_{1}^{2}
=165+0.35(400)=305\,\si{\newton}$, so
\[
x=\frac{1400}{2(0.35)}\ln\!\frac{480}{305}=2000\ln(1.574)=907\,\si{
\meter}.
\]
The car coasts nearly a kilometre shedding a third of its speed, and a
naive constant-force estimate using the initial resistance
$480\,\si{\newton}$ would give $x=m(v_{0}^{2}-v_{1}^{2})/(2\times 480)=
1400(500)/960=729\,\si{\meter}$, short by a fifth, because it ignores
the resistance falling as the car slows. The coast-down distance is the
measurement a manufacturer uses to separate $c_{r}$ from $C_{d}A$: two
coast intervals at different speeds give two equations in the two
unknowns, the road-test analogue of the terminal-velocity balance.

\subsection{A bouncing ball: restitution as a geometric series}

A ball dropped onto a hard floor bounces to a fraction of its release
height each time, and the sequence of bounces is the cleanest place to
watch a coefficient of restitution act over many impacts rather than
one (Figure~\ref{fig:vol03ch04:bounce-series}). A ball is released from
rest at height $h_{0}=2.0\,\si{\meter}$ onto a floor with coefficient of
restitution $e=0.80$. Find the height of the $n$th bounce, the total
distance the ball travels before coming to rest, and the total time it
spends bouncing.

\textbf{Each bounce scales the height by $e^{2}$.} The ball reaches the
floor with speed $v=\sqrt{2gh_{0}}$; restitution sends it back up at
$ev$, which rises to $h_{1}=(ev)^{2}/(2g)=e^{2}h_{0}$. The rule
compounds: the $n$th peak is
\[
h_{n}=e^{2n}h_{0},
\]
so the peaks form a geometric sequence with ratio $e^{2}=0.64$. Here
$h_{1}=0.64(2.0)=1.28\,\si{\meter}$, $h_{2}=0.819\,\si{\meter}$, and the
heights fall off geometrically toward zero.

\textbf{Total path length by a geometric sum.} After the first drop of
$h_{0}$, the ball climbs and falls $h_{n}$ on each subsequent bounce,
contributing $2h_{n}$ to the path. The total distance is
\[
D=h_{0}+2\sum_{n=1}^{\infty}e^{2n}h_{0}=h_{0}\Bigl(1+\frac{2e^{2}}{1-
e^{2}}\Bigr)=h_{0}\,\frac{1+e^{2}}{1-e^{2}}.
\]
With $e^{2}=0.64$, $D=2.0(1.64/0.36)=9.11\,\si{\meter}$: the ball
travels more than four times its release height before stopping,
because the early bounces are tall and many.

\textbf{Total bounce time, also a geometric sum.} The time for a free
flight to height $h$ and back is $2\sqrt{2h/g}$; the first drop alone is
$\sqrt{2h_{0}/g}$. Each bounce height scales by $e^{2}$, so each flight
time scales by $e$, and the times form a geometric series with ratio
$e$:
\[
t_{\infty}=\sqrt{\frac{2h_{0}}{g}}+2\sqrt{\frac{2h_{0}}{g}}
\sum_{n=1}^{\infty}e^{n}=\sqrt{\frac{2h_{0}}{g}}\Bigl(1+\frac{2e}{1-e}
\Bigr)=\sqrt{\frac{2h_{0}}{g}}\,\frac{1+e}{1-e}.
\]
With $\sqrt{2(2.0)/9.81}=0.639\,\si{\second}$ and $(1+e)/(1-e)=1.8/0.2=
9.0$, the ball stops bouncing after $t_{\infty}=0.639(9.0)=5.75\,\si{
\second}$. The result is the physical face of a convergent series: the
ball completes infinitely many ever-shorter bounces in a finite time,
and the run of clicks a dropped ball makes crowds into silence at
exactly $t_{\infty}$. A restitution nearer one stretches the tail
($e=0.95$ gives $(1+e)/(1-e)=39$), which is why a superball rattles on
far longer than a tennis ball from the same height.

\subsection{A rotating habitat: artificial gravity and its Coriolis
price}

A spun-up space habitat replaces gravity with the centripetal
requirement of circular motion, and pricing the trade shows why the
Coriolis term the chapter opened with becomes a design constraint rather
than a curiosity (Figure~\ref{fig:vol03ch04:spin-gravity}). A ring
habitat is to give its crew an apparent weight equal to Earth gravity at
the rim. Find the spin rate for a rim radius $R=90\,\si{\meter}$, the
apparent-gravity gradient a standing crew member feels head to foot, and
the sideways Coriolis acceleration on a tool dropped from the rim.

\textbf{The spin rate for one gravity.} A body at the rim of a ring
turning at rate $\Omega$ needs a centripetal acceleration $\Omega^{2}R$,
supplied by the floor pushing inward; the crew feel the reaction as an
outward weight $m\Omega^{2}R$. Setting this equal to $mg$ gives
\[
\Omega=\sqrt{\frac{g}{R}}=\sqrt{\frac{9.81}{90}}=0.330\,\si{\radian\per
\second},
\]
a rotation period $T=2\pi/\Omega=19.0\,\si{\second}$, about three
revolutions per minute. The rim speed is $\Omega R=29.8\,\si{\meter\per
\second}$.

\textbf{The gradient a standing body feels.} Apparent gravity falls with
radius as $\Omega^{2}r$, so a crew member of height $1.8\,\si{\meter}$
standing with feet at $R$ and head at $R-1.8$ feels less weight at the
head. The fractional change over the height is $\Delta g/g=(R-(R-1.8))/R
=1.8/90=0.02$, a two per cent head-to-foot gradient. A small habitat
worsens the gradient: halving the radius and doubling the spin rate to
hold one gravity doubles the gradient to four per cent, the kind of
mismatch between the pull on the head and the feet that unsettles the
inner ear.

\textbf{The Coriolis push on a dropped tool.} A tool released from rest
at the rim falls inward, and in the rotating frame it picks up a
Coriolis acceleration $2\Omega v_{r}$ perpendicular to its inward
velocity $v_{r}$. Early in the fall, when $v_{r}$ has grown to, say,
$3\,\si{\meter\per\second}$, the sideways acceleration is $2(0.330)(3)=
1.98\,\si{\meter\per\second\squared}$, a fifth of the apparent gravity,
enough that the tool visibly lands ahead of the spot straight below the
release point rather than at it. The deflection is why proposed spin
habitats are drawn large and slow: a big radius buys a low spin rate,
which shrinks both the head-to-foot gradient and the Coriolis push that
make a fast small ring disorienting to move in.

\subsection{A tow-rope snatch: impulsive tension with slack}

A vehicle recovered by a rope that starts slack loads the rope not by
the steady pull a driver imagines but by an impulsive jerk when the
slack runs out, and the impulse-momentum theorem sizes the peak the rope
must survive. A stalled car of mass $m=1{,}500\,\si{\kilogram}$ is towed
by a recovery vehicle through a rope that is momentarily slack; the
recovery vehicle is already moving so that the relative speed when the
rope snaps taut is $v_{\text{rel}}=2.5\,\si{\meter\per\second}$. The rope
stretches and arrests the relative motion over a snatch time
$\Delta t=0.20\,\si{\second}$. Find the impulse the rope carries, the
mean tension, and the peak tension for a triangular force pulse.

\textbf{The impulse follows from the velocity change.} Treat the stalled
car alone. The rope must change its speed by the relative closing speed
$v_{\text{rel}}$ over the snatch, so the impulse the rope delivers is
\[
J=m\,\Delta v=1500(2.5)=3{,}750\,\si{\newton\second}.
\]
The impulse is fixed by the mass and the closing speed, whatever the
detailed shape of the tension pulse; it is the integrated area under the
tension-versus-time curve.

\textbf{Mean and peak tension.} The mean tension over the snatch is the
impulse over the time,
\[
\bar{T}=\frac{J}{\Delta t}=\frac{3{,}750}{0.20}=18.8\,\si{\kilo\newton},
\]
already above the stalled car's weight of $14.7\,\si{\kilo\newton}$. The
tension does not hold at the mean; it rises from zero as the rope
stretches to a peak and falls again. For a triangular pulse the peak is
twice the mean, $T_{\text{peak}}=2\bar{T}=37.5\,\si{\kilo\newton}$, and a
stiffer rope that arrests the motion in half the time doubles both the
mean and the peak. The lesson the snatch teaches is the one the
impulse-momentum section stated: a fixed velocity change delivered over
a short time demands a large force, and the way to soften a recovery is
to lengthen $\Delta t$ with a stretchy kinetic rope that spreads the same
$3{,}750\,\si{\newton\second}$ over a longer snatch, which is why a
proper recovery strap is elastic and a chain, which arrests the motion in
a few milliseconds, can shatter a tow eye.

\subsection{A spool pulled by its string: the direction of rolling}

A spool wound with string on an inner hub rolls in a direction most
first-time solvers guess wrong, and the moment equation about the
contact point settles it without argument
(Figure~\ref{fig:vol03ch04:pulled-spool}). A spool of outer radius $R$,
inner hub radius $r$, mass $m$, and centre moment of inertia
$I_{G}=mk^{2}$ rests on a rough floor. A horizontal string wound off the
bottom of the hub is pulled with a steady force $P$. The spool rolls
without slipping. Find the acceleration of the centre and the direction
it rolls.

\textbf{Take moments about the contact point.} The friction force acts
at the contact $C$, so taking moments about $C$ removes it from the
equation and leaves only the applied pull. The string leaves the bottom
of the hub, a height $R-r$ above the contact, so its moment about $C$ is
$P(R-r)$, and it drives the spool in the direction of the pull. The
rolling body turns about $C$ as an instantaneous centre with moment of
inertia $I_{C}=I_{G}+mR^{2}=m(k^{2}+R^{2})$, so the moment equation about
$C$ is
\[
P(R-r)=I_{C}\alpha=m(k^{2}+R^{2})\alpha,
\]
and with $a_{G}=\alpha R$ the centre accelerates at
\[
a_{G}=\frac{P(R-r)R}{m(k^{2}+R^{2})},
\]
positive, so the spool rolls \emph{toward} the pull and winds the
remaining string up as it comes. The result surprises because intuition
expects the string to unwind and drive the spool away; the moment about
the contact shows the pull instead rotates the whole spool the other way.

\textbf{The critical case.} Tilt the pull upward until its line of action
passes through the contact point $C$; then the moment arm is zero, the
moment vanishes, and the spool neither rolls forward nor back. Past that
angle the moment reverses sign and the spool rolls away from the pull.
The critical string angle above the horizontal satisfies $\cos\theta_{
\text{crit}}=r/R$, the geometry that aims the string line through $C$,
and it is the boundary an operator winding a cable spool must respect:
pull below the critical angle and the spool reels the cable in, pull
above it and the spool pays the cable out. Numbers make it concrete: a
spool with $R=0.30\,\si{\meter}$ and $r=0.12\,\si{\meter}$ has
$\cos\theta_{\text{crit}}=0.40$, so $\theta_{\text{crit}}=66^{\circ}$;
any pull flatter than $66^{\circ}$ above horizontal rolls the spool
toward the operator.

\subsection{A block on an accelerating wedge: working in the
non-inertial frame}

A block riding on a wedge that is itself accelerating is the cleanest
problem in which the pseudo-force method of the non-inertial mastery box
pays for itself, turning a two-body kinematics tangle into a single
static balance in the wedge's frame. A wedge with a frictionless incline
at angle $\beta=30^{\circ}$ is pushed along a level floor with a
horizontal acceleration $A$. A block of mass $m$ rests on the incline.
Find the horizontal acceleration $A$ that holds the block stationary
relative to the wedge, requiring no normal-direction sliding.

\textbf{Move into the wedge frame and add the pseudo-force.} In the
frame of the accelerating wedge the block is at rest, so the applied
forces plus the inertial pseudo-force $-mA$ (pointing backward, opposite
the wedge's acceleration) must balance. The forces on the block are
gravity $mg$ down, the normal force $N$ perpendicular to the incline,
and the horizontal pseudo-force $mA$ pointing backward. Resolve along
the incline, where the frictionless surface can supply nothing: the
block stays put only if the along-incline components of gravity and the
pseudo-force cancel. Taking the incline rising in the direction the
wedge is pushed, gravity's along-incline component is $mg\sin\beta$ down
the slope and the pseudo-force's is $mA\cos\beta$ up the slope, so
balance requires
\[
mA\cos\beta=mg\sin\beta \ \Rightarrow\ A=g\tan\beta.
\]
With $\beta=30^{\circ}$, $A=9.81\tan 30^{\circ}=5.66\,\si{\meter\per
\second\squared}$. Pushed at exactly this acceleration the block hangs on
the frictionless slope without sliding, the incline analogue of the
banked-curve design speed at which the bank alone turns the car.

\textbf{The normal force it takes.} Resolving perpendicular to the
incline gives the normal force the surface must carry,
$N=mg\cos\beta+mA\sin\beta=m g(\cos\beta+\tan\beta\sin\beta)=mg/\cos
\beta$, which for $\beta=30^{\circ}$ is $mg/0.866=1.15\,mg$, the same
$1/\cos\beta$ surcharge the conical pendulum's string tension carried.
Pushed harder than $g\tan\beta$ the block would climb the slope; pushed
softer it would slide down. The inertial-frame route to the same answer
tracks both the block and the wedge in ground coordinates and eliminates
the shared horizontal acceleration by hand; the pseudo-force route
collapses that bookkeeping into one static triangle, the payoff
d'Alembert's principle promised.

% Figure: a ball bouncing with coefficient of restitution e < 1. Each
% bounce returns to e^2 times the previous height, so the peaks form a
% geometric sequence and the whole motion ends in a finite time.
\begin{figure}[htbp]
\centering
\begin{tikzpicture}[scale=1.0]
% ground
\draw[engblack, thick] (0,0) -- (11.2,0);
% first drop from h0 = 3, then peaks at e^2 h with e^2 = 0.64
% peak heights: 3, 1.92, 1.229, 0.786, 0.503
% We draw parabolic arcs between contact points spaced by shrinking widths.
% contact x-positions chosen to shrink like the time between bounces (factor e)
\def\hzero{3.0}
% arc 1: fall from (0,hzero) to (1.6,0)
\draw[engblue, very thick] (0,\hzero) .. controls (0.8,1.3) .. (1.6,0);
% arc 2: up to peak 1.92 then down; width scales by e=0.8 -> 2.56 wide
\draw[engblue, very thick] (1.6,0) .. controls (2.88,2.56) .. (4.16,0);
% arc 3: peak 1.229, width 2.048
\draw[engblue, very thick] (4.16,0) .. controls (5.184,1.64) .. (6.208,0);
% arc 4: peak 0.786, width 1.638
\draw[engblue, very thick] (6.208,0) .. controls (7.027,1.05) .. (7.846,0);
% arc 5: peak 0.503, width 1.311
\draw[engblue, very thick] (7.846,0) .. controls (8.501,0.67) .. (9.157,0);
% remaining bounces collapse toward the accumulation point
\draw[engblue, very thick] (9.157,0) .. controls (9.68,0.43) .. (10.2,0);
\draw[engblue, very thick] (10.2,0) .. controls (10.53,0.27) .. (10.86,0);
% accumulation point marker
\fill[engred] (11.0,0) circle (2pt);
\node[font=\scriptsize, engred, anchor=south west] at (11.02,0.02)
  {rest at finite $t_\infty$};
% height markers
\draw[enggray, dashed] (0,\hzero) -- (-0.15,\hzero);
\node[font=\scriptsize, anchor=east] at (-0.15,\hzero) {$h_0$};
\draw[enggray, dashed] (2.88,1.92) -- (-0.15,1.92);
\node[font=\scriptsize, anchor=east] at (-0.15,1.92) {$e^{2}h_0$};
\draw[enggray, dashed] (5.184,1.229) -- (-0.15,1.229);
\node[font=\scriptsize, anchor=east] at (-0.15,1.229) {$e^{4}h_0$};
% axis label
\node[font=\scriptsize, anchor=north] at (5.6,-0.15) {time / horizontal offset (drawn shrinking by $e$ each bounce)};
\end{tikzpicture}
\caption[Bounce heights of a ball with restitution below one]{A ball
dropped from height $h_0$ onto a floor of restitution $e<1$ returns to
$e^{2}h_0$ after the first bounce, $e^{4}h_0$ after the second, and
$e^{2n}h_0$ after the $n$th, a geometric sequence of peaks. The time
between successive bounces shrinks by the factor $e$, so the arcs
crowd together and the infinite sequence of bounces completes in a
finite total time $t_\infty$. The ball comes to rest after infinitely
many ever-shorter bounces, the same accumulation a mathematician meets
in a convergent geometric series and an engineer hears as the run of
clicks that ends in silence.}
\label{fig:vol03ch04:bounce-series}
\end{figure}


% Figure: a rotating habitat producing artificial gravity at the rim.
% A tool dropped from the rim deflects sideways relative to the floor,
% the Coriolis signature seen inside the rotating frame.
\begin{figure}[htbp]
\centering
\begin{tikzpicture}[scale=1.0,
  lbl/.style={font=\scriptsize},
]
% station rim
\draw[engblue, very thick] (0,0) circle (3.0);
\draw[engblue, thick, dashed] (0,0) circle (2.55); % floor thickness hint
\fill[engblack] (0,0) circle (1.5pt);
\node[lbl, anchor=north east] at (-0.05,-0.05) {axis};
% spin arrow
\draw[-{Stealth}, engamber, very thick] (2.1,2.1) arc (45:15:2.97);
\node[lbl, engamber, anchor=south west] at (2.15,2.15) {$\Omega$};
% a person standing at the rim (bottom), feet outward
\draw[engblack, very thick] (0,-3.0) -- (0,-2.35);
\fill[engblack] (0,-2.25) circle (0.13);
\node[lbl, anchor=west] at (0.18,-2.7) {crew};
% apparent weight arrow (centripetal, points to axis; felt as outward "down")
\draw[-{Stealth}, engred, very thick] (0,-2.9) -- (0,-2.0)
  node[right, lbl, engred] {felt weight $m\Omega^{2}R$};
% dropped tool at right rim, released; radially inward "fall" plus Coriolis
% release point at (3,0); ideal inward path is along -x; actual path curves
\fill[engblue] (3.0,0) circle (2pt);
\node[lbl, engblue, anchor=south west] at (2.75,0.15) {release};
% straight radial (naive) path
\draw[enggray, dashed] (3.0,0) -- (0.6,0);
\node[lbl, enggray, anchor=south] at (1.7,0.02) {naive radial};
% actual curved path (Coriolis deflects against the spin sense)
\draw[engblue, very thick, -{Stealth}]
  (3.0,0) .. controls (2.2,-0.35) and (1.3,-0.7) .. (0.55,-0.85);
\node[lbl, engblue, anchor=north] at (1.4,-0.75) {actual (Coriolis)};
\end{tikzpicture}
\caption[Artificial gravity and the Coriolis deflection in a spinning
habitat]{A rotating habitat spun at rate $\Omega$ presses its crew
against the rim: a person standing at radius $R$ feels an apparent
weight $m\Omega^{2}R$ directed outward, the reaction to the centripetal
force the floor supplies, and choosing $\Omega=\sqrt{g/R}$ sets that
apparent weight equal to Earth gravity. Motion inside the habitat
carries the Coriolis signature. A tool released from the rim would, to
a naive eye, fall straight in along a radius (dashed), but in the
rotating frame it deflects sideways (solid) because the Coriolis
acceleration $2\Omega v_{r}$ acts perpendicular to its inward motion.
The deflection is the reason a small habitat, which needs a high spin
rate for a given apparent gravity, feels disorienting, and the reason
proposed spin habitats are drawn large and slow.}
\label{fig:vol03ch04:spin-gravity}
\end{figure}


% Figure: a spool pulled by a string wound under its inner hub. The
% direction the spool rolls follows from the moment about the contact
% point; a pull whose line of action passes through the contact gives
% no rolling at all, the critical-angle case.
\begin{figure}[htbp]
\centering
\begin{tikzpicture}[scale=1.0,
  frc/.style={draw=engblue, -{Stealth}, very thick},
  lbl/.style={font=\scriptsize},
]
% left panel: horizontal pull under the hub, spool rolls toward the pull
\begin{scope}[shift={(1.2,1.3)}]
  \draw[engblack, thick] (-1.8,-1.3) -- (2.6,-1.3);
  \draw[engblue, thick, fill=engblue!5] (0,0) circle (1.3);   % outer radius R
  \draw[enggray, thick, fill=enggray!12] (0,0) circle (0.6);  % inner hub r
  \fill[engblack] (0,0) circle (1.2pt);
  \fill[engred] (0,-1.3) circle (1.6pt);
  \node[lbl, engred, anchor=north] at (0,-1.35) {$C$};
  % string leaves the bottom of the hub, pulled horizontally right
  \draw[frc] (0,-0.6) -- (2.3,-0.6) node[right, lbl, engblue] {$P$};
  % moment arm from C up to the line of action
  \draw[enggray, dashed] (0,-1.3) -- (0,-0.6);
  \node[lbl, enggray, anchor=east] at (-0.05,-0.95) {$R{-}r$};
  \node[lbl, engamber, anchor=south] at (0,1.45) {horizontal pull};
  \draw[-{Stealth}, engamber, thick] (-0.75,-1.62) arc (150:30:0.5);
  \node[lbl, engamber, anchor=north] at (0,-1.75) {rolls toward $P$};
\end{scope}
% right panel: pull aimed through the contact point, no rolling
\begin{scope}[shift={(6.2,1.3)}]
  \draw[engblack, thick] (-1.8,-1.3) -- (2.6,-1.3);
  \draw[engblue, thick, fill=engblue!5] (0,0) circle (1.3);
  \draw[enggray, thick, fill=enggray!12] (0,0) circle (0.6);
  \fill[engblack] (0,0) circle (1.2pt);
  \fill[engred] (0,-1.3) circle (1.6pt);
  \node[lbl, engred, anchor=north] at (0,-1.35) {$C$};
  % string leaves the bottom of the hub aimed so its line passes through C
  % hub-bottom at (0,-0.6); aim the line through C at (0,-1.3): direction
  % points down-and-right toward a far point; draw pull along that line
  \draw[frc] (0,-0.6) -- (2.03,-1.90) node[right, lbl, engblue] {$P$};
  % show the line continued to C
  \draw[enggray, dashed] (0,-0.6) -- (0,-1.3);
  \node[lbl, engamber, anchor=south] at (0,1.45) {pull aimed through $C$};
  \node[lbl, engamber, anchor=north] at (0,-1.75) {no rolling};
\end{scope}
\end{tikzpicture}
\caption[A spool pulled by a string wound on its inner hub]{A spool of
outer radius $R$ pulled by a string wound on an inner hub of radius $r$.
The direction of rolling follows from the moment of the pull about the
contact point $C$, taken as the instantaneous centre. A horizontal pull
off the bottom of the hub (left) has a line of action a height $R-r$
above the contact, so its moment $P(R-r)$ rolls the spool toward the
pull, winding the string up as it comes, the result that surprises most
first-time solvers. Tilt the pull until its line passes through the
contact (right) and the moment arm vanishes: the spool neither rolls
forward nor back but is on the verge of sliding. The lesson is to take
moments about $C$ rather than about the centre, because the friction
force at $C$ then drops out of the moment equation and the sense of
rolling reads straight off the geometry of the pull.}
\label{fig:vol03ch04:pulled-spool}
\end{figure}


\subsection{Vignette: a tuned mass damper on a tall structure}

\engterm{Vignette.} A slender tower, a long footbridge, or a tall
chimney has a low natural frequency and little inherent damping, so a
wind gust or a crowd's footfall that matches that frequency can build a
sway to alarming amplitude. The cure is not to stiffen the structure,
which is expensive and often impossible, but to hang a second mass on a
spring and damper inside it, tuned so that its own natural frequency
matches the structure's, and let the two exchange energy
(Figure~\ref{fig:vol03ch04:tmd-response}).

\textbf{Why the absorber works.} Model the structure near its resonance
as a single mass $M$ on a spring $K$ with natural frequency
$\omega_{n}=\sqrt{K/M}$. Attach a smaller mass $m$ on a spring $k$ and
damper $c$, chosen so the absorber's natural frequency $\sqrt{k/m}$
equals $\omega_{n}$. When the structure tries to sway at $\omega_{n}$,
the absorber swings in opposition, a quarter-cycle behind, so the spring
force it feeds back to the structure pushes against the motion. The
single tall resonance peak of the bare structure splits into two lower
peaks, and the damper $c$ on the absorber bleeds the resonant energy
away as heat. The absorber does not carry any static load and adds no
stiffness; it works purely by moving out of phase and dissipating in its
own damped mass. The mass ratio $m/M$ sets how far apart the two peaks
push and how much the sway drops; a few per cent of the structure's mass
is typical, so the device is small against the thing it protects.

\textbf{What the numbers look like.} A footbridge of modal mass
$M=50\,\si{\tonne}$ swaying at $f_{n}=1.0\,\si{\hertz}$ under a marching
crowd might carry an absorber of $m=1{,}000\,\si{\kilogram}$, two per
cent of the modal mass, on a spring tuned to $1.0\,\si{\hertz}$
($k=m\omega_{n}^{2}=1{,}000(2\pi)^{2}\approx 39.5\,\si{\kilo\newton\per
\meter}$) with a damper set near the optimum $\zeta\approx 0.09$ for that
mass ratio. Such an absorber can cut the peak sway by more than half, the
difference between a bridge that feels lively and one that feels solid.
Tall buildings carry the same idea at large scale: a heavy block or a
suspended pendulum near the top, riding on bearings or oil dampers,
tuned to the building's sway period of several seconds. The mechanism is
the one this chapter has built from the ground up: a second mass whose
inertia and damping are placed to fight a resonance the primary mass
cannot damp on its own. Den Hartog's classical treatment gives the
optimum tuning and damping in closed form
\cite{text:den-hartog-vibrations}.

% Figure: the effect of a tuned mass damper. The steady-state response
% amplitude of a structure against forcing frequency, with and without
% an attached, tuned, damped absorber. The single resonance peak is
% split into two smaller peaks.
\begin{figure}[htbp]
\centering
\begin{tikzpicture}
\begin{axis}[
  width=12cm, height=7cm,
  xlabel={forcing frequency ratio $\omega/\omega_n$},
  ylabel={response amplitude (normalised)},
  xmin=0.5, xmax=1.5, ymin=0, ymax=10,
  axis lines=left,
  samples=300,
  legend pos=north east,
  legend cell align=left,
  grid=major, grid style={enggray!20},
  tick label style={font=\scriptsize},
  label style={font=\small},
  legend style={font=\scriptsize},
]
% bare structure, light damping zeta = 0.02: 1/sqrt((1-r^2)^2 + (2 zeta r)^2)
\addplot[engred, very thick, domain=0.5:1.5]
  {1/sqrt((1-x^2)^2 + (2*0.02*x)^2)};
\addlegendentry{bare structure}
% with tuned mass damper: a two-peaked curve, sketched with a smooth
% double-lobe. Use a fitted rational-like proxy that stays bounded.
\addplot[engblue, very thick, domain=0.5:1.5, samples=400]
  {2.6/sqrt((1 - (x/0.90)^2)^2 + (2*0.10*(x/0.90))^2)
   + 2.6/sqrt((1 - (x/1.11)^2)^2 + (2*0.10*(x/1.11))^2) - 1.4};
\addlegendentry{with tuned mass damper}
% mark the tuning frequency
\addplot[enggray, dashed, thick] coordinates {(1,0) (1,10)};
\node[font=\scriptsize, enggray, anchor=south west] at (axis cs:1.01,8.6)
  {tuned at $\omega_n$};
\end{axis}
\end{tikzpicture}
\caption[Response of a structure with a tuned mass damper]{The
steady-state response amplitude of a lightly damped structure driven
near its natural frequency $\omega_n$. Alone (red) the structure shows
one tall, sharp resonance peak that a matched forcing, wind gust or
footfall, can drive to damaging amplitude. Adding a smaller mass on a
spring and damper, tuned so its own natural frequency matches
$\omega_n$, splits the single peak into two lower ones (blue): the
absorber and the structure exchange energy out of phase, and the damper
on the absorber bleeds the vibration away. The absorber does not stiffen
the structure; it moves in opposition to it and carries the resonant
motion into its own damped mass. The device is why tall towers hang a
pendulum or a sliding block near the top, and why a footbridge that
sways underfoot is cured with a tuned absorber rather than a rebuild.}
\label{fig:vol03ch04:tmd-response}
\end{figure}


\subsection{Vignette: a jet turned by a vane and the impulse turbine}

\engterm{Vignette.} A stream of water driven from a nozzle carries
momentum, and a vane placed in its path that turns the stream must
supply the force to change that momentum, feeling an equal and opposite
push in return (Figure~\ref{fig:vol03ch04:jet-vane}). The steady
momentum-flux argument of the impulse-momentum section sizes the force
and, when the vane is allowed to move, sets the speed at which it
extracts the most power, the working principle of the impulse water
turbine.

\textbf{Force on a stationary vane.} A jet of cross-section $A$ leaving a
nozzle at speed $v$ carries a mass flow $\dot{m}=\rho A v$. Brought to
rest against a flat plate, its momentum flux $\dot{m}v=\rho A v^{2}$ is
destroyed and the plate feels a force of that magnitude. A cupped vane
that turns the jet through $180^{\circ}$ does better: the stream leaves
with its momentum reversed, a change of $2\dot{m}v$, so the force is
$2\rho A v^{2}$, twice the flat-plate value. A jet $30\,\si{\milli
\meter}$ across ($A=7.07\times 10^{-4}\,\si{\meter\squared}$) at
$40\,\si{\meter\per\second}$ carries $\dot{m}=1{,}000(7.07\times
10^{-4})(40)=28.3\,\si{\kilogram\per\second}$ and pushes a reversing
vane with $2\dot{m}v=2(28.3)(40)=2.26\,\si{\kilo\newton}$, a substantial
load from a modest hose.

\textbf{A moving vane and the best speed.} If the vane recedes at speed
$u$ in the jet's direction, the jet overtakes it at the relative speed
$v-u$, and the mass actually caught per unit time on a single moving
vane, or the steady mass on a wheel of vanes fed continuously, sets the
force. For a wheel of buckets fed by the full jet, the force is
$\dot{m}(v-u)(1+\cos\phi)$ with $\phi$ the turning angle, and the power
delivered is that force times $u$. Maximising $u(v-u)$ over $u$ gives
$u=v/2$: the bucket speed that extracts the most power is half the jet
speed, at which point an ideal $180^{\circ}$ bucket recovers nearly all
the jet's kinetic energy, leaving the water to fall away almost at rest.
Real impulse turbines, the Pelton wheels of high-head hydro plants, run
their buckets near this half-jet-speed optimum and reach efficiencies
above ninety per cent. The chain from the impulse-momentum theorem to
the turbine is direct: momentum carried by the stream, changed by the
vane, delivered as force, converted to power at the speed that best
matches the vane to the jet.

% Figure: a fluid jet striking a curved vane and being turned. The
% momentum flux of the incoming and outgoing stream sets the force on
% the vane. Shown for a vane that reverses the jet (turning angle 180),
% the impulse-turbine geometry.
\begin{figure}[htbp]
\centering
\begin{tikzpicture}[
  jet/.style={draw=engblue, -{Stealth}, very thick},
  frc/.style={draw=engred, -{Stealth}, very thick},
  lbl/.style={font=\scriptsize},
]
% the curved vane: a half-cup opening to the left
\draw[draw=engblack, very thick]
  (1.6,1.5) .. controls (-0.4,1.5) and (-0.4,-1.5) .. (1.6,-1.5);
\node[lbl, anchor=west] at (1.7,0) {vane};
% incoming jet from the left, hitting the cup centre
\draw[jet] (-3.4,0) -- (-0.35,0) node[midway, above, lbl, engblue] {$\dot{m},\,v$};
\node[lbl, anchor=south] at (-3.0,0.1) {incoming jet};
% outgoing streams, turned back (split up and down)
\draw[jet] (0.2,0.9) -- (-2.2,1.5) node[above, lbl, engblue] {$v'$};
\draw[jet] (0.2,-0.9) -- (-2.2,-1.5) node[below, lbl, engblue] {$v'$};
% reaction force on the vane, to the right (jet pushes vane in flow dir)
\draw[frc] (1.0,0) -- (3.0,0) node[right, lbl, engred] {$F$ on vane};
% if the vane moves, its velocity u
\draw[-{Stealth}, engamber, thick] (1.6,-2.0) -- (3.0,-2.0)
  node[right, lbl, engamber] {vane speed $u$};
\end{tikzpicture}
\caption[A jet turned by a curved vane]{A steady jet of mass flow
$\dot{m}$ and speed $v$ strikes a curved vane and is turned back on
itself. The vane changes the jet's momentum, and by the
impulse-momentum theorem it must exert the force that produces that
change; the jet pushes back on the vane with the equal and opposite
reaction (red), in the direction of the original flow. For a jet turned
through $180^{\circ}$ against a stationary vane the force is $2\dot{m}v$,
twice the value for a jet merely stopped, because the stream leaves with
its momentum reversed rather than only removed. When the vane itself
moves away at speed $u$, the relevant flow speed is the relative speed
$v-u$, and the power delivered to the vane, $F u$, is largest near
$u=v/2$, the result that sets the bucket speed of an impulse water
turbine.}
\label{fig:vol03ch04:jet-vane}
\end{figure}


\subsection{Vignette: a centrifugal clutch and its engagement speed}

\engterm{Vignette.} A chainsaw, a go-kart, and a mobility scooter all
start their engines with the drive disconnected and let the drive engage
by itself as the engine speeds up, so that the machine idles without
creeping and pulls away smoothly once the throttle is opened. The device
that does this without a pedal or a lever is the centrifugal clutch, and
it works by pitting a spring against the $\omega^{2}$ growth of the
centrifugal requirement that the rotating-machinery failure section met
as a hazard and this vignette meets as a tool.

\textbf{The engagement threshold is a force balance.} A centrifugal
clutch carries shoes that ride on the driving shaft and are held inward
by a spring until the shaft spins fast enough to fling them out against a
surrounding drum. A shoe of mass $m$ whose centre of mass sits at radius
$r$ from the shaft axis needs a centripetal force $m r\omega^{2}$ to hold
its circular path; the spring supplies a preload $F_{s}$ pulling it
inward. The shoe stays retracted while the spring wins and begins to
press on the drum only when the centrifugal demand overtakes the
preload, at the engagement speed
\[
m r\omega_{\text{eng}}^{2}=F_{s} \ \Rightarrow\ \omega_{\text{eng}}=
\sqrt{\frac{F_{s}}{m r}}.
\]
Below $\omega_{\text{eng}}$ the drive is free and the engine idles
disconnected; above it the shoes press on the drum with a net outward
force $m r\omega^{2}-F_{s}$ that grows as the square of the speed, and
the friction between shoe and drum, $\mu(m r\omega^{2}-F_{s})$ per shoe,
transmits torque to the output. The clutch is thus a speed-triggered
switch built from nothing but a mass, a spring, and a friction face.

\textbf{What the numbers look like.} A small clutch might carry two shoes
of $m=40\,\si{\gram}$ each with a centre-of-mass radius $r=25\,\si{
\milli\meter}$ and a spring preload $F_{s}=45\,\si{\newton}$ per shoe.
The engagement speed is
\[
\omega_{\text{eng}}=\sqrt{\frac{45}{0.040(0.025)}}=\sqrt{45{,}000}=212\,
\si{\radian\per\second},
\]
about $2{,}025$ revolutions per minute, comfortably above a typical idle
near $1{,}500$ and below the running speed of a few thousand, so the tool
idles free and drives under load. At a running speed of $6{,}000$
revolutions per minute ($\omega=628\,\si{\radian\per\second}$) the
outward force per shoe is $m r\omega^{2}-F_{s}=0.040(0.025)(628)^{2}-45=
394-45=349\,\si{\newton}$, and with two shoes at a friction coefficient
$\mu=0.35$ pressing at drum radius $R_{d}=40\,\si{\milli\meter}$ the
transmissible torque is $2\mu(m r\omega^{2}-F_{s})R_{d}=2(0.35)(349)(
0.040)=9.8\,\si{\newton\meter}$. The torque climbs steeply with speed
because the pressing force follows $\omega^{2}$, so a centrifugal clutch
grips harder exactly when the engine is making the most power, the
property that makes it self-scheduling. The cost is the same heat the
slipping-clutch worked example counted: every engagement slips while the
output speeds up to the input, and the energy lost to that slip heats the
shoes, which is why a clutch made to engage often is built with a heat
path to the drum and why riding a chainsaw at part-throttle, holding the
clutch on the edge of engagement, burns the shoes.

\section{End-to-end case study: a two-vehicle accident
reconstruction}\pathtag{standard}

An intersection collision puts the chapter's tools to work on one
event from the physical evidence to a defensible verdict. A car
travelling east on a through road struck a car travelling north that
had entered from a side road; the two vehicles locked and slid to rest
together. One driver claims the eastbound car was within the posted
$50\,\si{\kilo\meter\per\hour}$ limit. The evidence at the scene is a
pre-impact skid mark left by the eastbound car, the masses of the two
vehicles, the post-impact slide, and the geometry of the rest
position. The reconstruction proceeds in four steps, carrying the
numbers through.

\textbf{Evidence.} The eastbound car (car A) has mass $m_{A}=1{,}400
\,\si{\kilogram}$ and left a straight pre-impact skid of length
$d_{\text{skid}}=24\,\si{\meter}$ on dry asphalt before the point of
impact. The northbound car (car B) has mass $m_{B}=1{,}000\,\si{
\kilogram}$ and left no pre-impact skid; a witness puts its speed at
roughly $18\,\si{\meter\per\second}$ ($65\,\si{\kilo\meter\per\hour}$),
which the reconstruction will test against the physics. After impact
the locked wreck slid $d_{\text{post}}=11\,\si{\meter}$ along a
straight mark before stopping. The road is dry asphalt, drag
coefficient of friction $\mu=0.75$ for the locked-wheel skid and
$\mu_{\text{slide}}=0.65$ for the post-impact slide of the deformed,
partly-airborne wreck.

\textbf{Step 1: speed at the start of the skid, from the skid mark.}
A car in a full locked-wheel skid is decelerated by friction alone.
The work-energy theorem equates the kinetic energy at the start of the
skid to the friction work over the skid length:
\[
\tfrac{1}{2}m_{A}v_{A,\text{skid}}^{2}=\mu m_{A} g\,d_{\text{skid}}.
\]
The mass cancels, the signature of skid-based speed estimation:
\[
v_{A,\text{skid}}=\sqrt{2\mu g\,d_{\text{skid}}}=\sqrt{2(0.75)(9.81)
(24)}=18.8\,\si{\meter\per\second}.
\]
That is the speed at the instant the wheels locked, about
$67.7\,\si{\kilo\meter\per\hour}$. Figure~\ref{fig:vol03ch04:skid-speed}
reads this point off the dry-asphalt curve and shows how the square-root
relation makes the surface coefficient the dominant uncertainty: the
same $24\,\si{\meter}$ skid would read only $49\,\si{\kilo\meter\per
\hour}$ if the road had in fact been wet.

\textbf{Step 2: speed at impact.} The skid mark ends at the point of
impact, so $v_{A,\text{skid}}=18.8\,\si{\meter\per\second}$ is the
speed at the \emph{start} of the recorded skid, when the wheels first
locked, and the car kept braking over the $24\,\si{\meter}$ until it
struck. The collision needs car A's speed at the moment of impact,
$v_{A}$, which is lower than the start-of-skid speed by whatever
braking happened over the mark. The reconstruction takes the
conservative reading: it treats the impact as occurring at the end of
the skid, so $v_{A}\approx v_{A,\text{skid}}=18.8\,\si{\meter\per
\second}$, and records that the true speed when the driver first
reacted was higher still, because $18.8\,\si{\meter\per\second}$ is
already a value the car had braked down to. This impact speed is what
enters the collision below; any refinement that splits the skid into
pre- and post-lock segments only raises the inferred pre-brake speed,
so it would strengthen, not weaken, the verdict that follows.

\textbf{Step 3: the collision, by conservation of momentum.} The two
cars lock together, a perfectly inelastic collision. Momentum is
conserved across the brief contact in both directions independently,
because the impulsive contact forces are internal to the two-car
system. Car A carries momentum east; car B carries momentum north
(Figure~\ref{fig:vol03ch04:oblique-collision}). With car B's witnessed
$v_{B}=18\,\si{\meter\per\second}$:
\[
p_{x}=m_{A}v_{A}=1400(18.8)=26{,}300\,\si{\kilogram\meter\per\second},
\]
\[
p_{y}=m_{B}v_{B}=1000(18.0)=18{,}000\,\si{\kilogram\meter\per\second}.
\]
The locked wreck (mass $m_{A}+m_{B}=2{,}400\,\si{\kilogram}$) departs
with these components:
\[
v_{\text{wreck}}=\frac{\sqrt{p_{x}^{2}+p_{y}^{2}}}{m_{A}+m_{B}}
=\frac{\sqrt{26{,}300^{2}+18{,}000^{2}}}{2400}
=\frac{31{,}870}{2400}=13.3\,\si{\meter\per\second},
\]
heading $\phi=\arctan(p_{y}/p_{x})=\arctan(18{,}000/26{,}300)=
34.4^{\circ}$ north of east. The kinetic energy lost in the crush is
\[
\Delta T=\Bigl(\tfrac{1}{2}m_{A}v_{A}^{2}+\tfrac{1}{2}m_{B}v_{B}^{2}
\Bigr)-\tfrac{1}{2}(m_{A}+m_{B})v_{\text{wreck}}^{2}
=(247+162)-212=197\,\si{\kilo\joule},
\]
in kilojoules, about $48\,\%$ of the incident kinetic energy, absorbed
by the deforming structures, heat, and sound. The momentum survives the
contact intact; the energy does not, the contrast the collision figure
draws.

\textbf{Step 4: the independent cross-check from the post-impact
slide.} The reconstruction has two routes to the wreck's departure
speed, and a sound verdict needs them to agree. The post-impact slide
length gives a second, independent estimate of $v_{\text{wreck}}$
through the same friction-energy balance used for the skid:
\[
v_{\text{wreck,slide}}=\sqrt{2\mu_{\text{slide}}\,g\,d_{\text{post}}}
=\sqrt{2(0.65)(9.81)(11)}=11.8\,\si{\meter\per\second}.
\]
The momentum route gave $13.3\,\si{\meter\per\second}$ and the slide
route gives $11.8\,\si{\meter\per\second}$, a $12\,\%$ spread. The two
estimates rest on different evidence (the masses and incoming speeds
for the first, the post-impact mark and slide friction for the second),
so agreement to within the friction-coefficient uncertainty is the
reconstruction's internal consistency check. The spread is well inside
the $\pm 20\,\%$ band that the friction coefficient alone carries, so
the reconstruction holds together.

\textbf{Verdict.} Car A entered its skid at $18.8\,\si{\meter\per
\second}$, or $67.7\,\si{\kilo\meter\per\hour}$, and was still moving
near that speed at impact. The posted limit is $50\,\si{\kilo\meter\per
\hour}$ ($13.9\,\si{\meter\per\second}$). The skid-based speed exceeds
the limit by $36\,\%$, and it is a lower bound on the pre-brake speed
because the driver had already shed speed braking through the skid.
The claim that car A was within the limit is not consistent with the
physical evidence: a car obeying the limit, braking on dry asphalt at
$\mu=0.75$, would have stopped in $v^{2}/(2\mu g)=13.9^{2}/(2\times
0.75\times 9.81)=13.1\,\si{\meter}$, far short of the
$24\,\si{\meter}$ skid actually left. The two independent estimates of
the wreck speed agree, which raises confidence that the input masses
and the witnessed northbound speed are sound. The reconstruction's one
soft input is the surface coefficient; a forensic report would state
the speed as a band, $\mu=0.7$ to $0.8$ giving
$18.2$ to $19.4\,\si{\meter\per\second}$, and would record the road
condition, tyre state, and temperature on which the band depends.

\begin{principle}[Two estimates from independent evidence]
A reconstruction earns its verdict by computing the same unknown two
ways from evidence that does not share a common error, then checking
that the answers agree to within the inputs' stated uncertainty. Here
the wreck's departure speed comes once from the conservation of
momentum in the collision and once from the friction work of the
post-impact slide; the masses and incoming speeds feed only the first,
the slide mark and slide friction only the second. Agreement is
evidence; a wide disagreement would point to a wrong mass, a
mismeasured mark, or a misjudged surface, and would have to be resolved
before any speed is quoted. A single calculation, however careful, has
no such check on its inputs.
\end{principle}

% Figure: pre-braking speed inferred from skid-mark length for three road
% surfaces, v = sqrt(2 mu g d), plotted against skid length.
\begin{figure}[htbp]
\centering
\begin{tikzpicture}
\begin{axis}[
  width=12cm, height=7cm,
  xlabel={skid-mark length $d$ (\si{\meter})},
  ylabel={inferred speed $v$ (\si{\kilo\meter\per\hour})},
  xmin=0, xmax=60, ymin=0, ymax=160,
  axis lines=left,
  domain=0:60, samples=120,
  legend pos=north west,
  legend cell align=left,
  grid=major, grid style={enggray!20},
  tick label style={font=\scriptsize},
  label style={font=\small},
  legend style={font=\scriptsize},
]
% v = sqrt(2 mu g d) in m/s, then * 3.6 to km/h. g = 9.81.
% dry asphalt mu = 0.75: 3.6*sqrt(2*0.75*9.81*x) = 3.6*sqrt(14.715*x)
\addplot[engblue, very thick] {3.6*sqrt(14.715*x)};
\addlegendentry{dry asphalt, $\mu=0.75$}
% wet asphalt mu = 0.40: 3.6*sqrt(2*0.40*9.81*x) = 3.6*sqrt(7.848*x)
\addplot[engamber, very thick] {3.6*sqrt(7.848*x)};
\addlegendentry{wet asphalt, $\mu=0.40$}
% packed snow mu = 0.18: 3.6*sqrt(2*0.18*9.81*x) = 3.6*sqrt(3.532*x)
\addplot[engred, very thick] {3.6*sqrt(3.532*x)};
\addlegendentry{packed snow, $\mu=0.18$}
% mark the case-study point: dry asphalt, d = 24 m -> v ~= 67.5 km/h
\addplot[only marks, mark=*, engblack, mark size=2pt] coordinates {(24,67.4)};
\node[font=\scriptsize, anchor=south west] at (axis cs:24,67.4) {$d=24\,\si{\meter}\Rightarrow 67\,\si{\kilo\meter\per\hour}$};
\end{axis}
\end{tikzpicture}
\caption[Pre-braking speed from skid-mark length]{Pre-braking speed
inferred from a measured skid-mark length, $v=\sqrt{2\mu g d}$, for
three road surfaces. The relation comes from the work-energy theorem:
the friction force $\mu m g$ acting over the skid distance $d$ removes
all the kinetic energy $\tfrac{1}{2}mv^{2}$, and the mass cancels.
Because the speed grows as the square root of the distance, a long
skid implies a high speed only weakly, while the surface coefficient
$\mu$ enters under the same root and dominates the answer: the same
$24\,\si{\meter}$ skid reads $67\,\si{\kilo\meter\per\hour}$ on dry
asphalt (black marker) but only about $49\,\si{\kilo\meter\per\hour}$
on a wet road. The accident-reconstruction case study reads the dry
curve at $d=24\,\si{\meter}$.}
\label{fig:vol03ch04:skid-speed}
\end{figure}


% Figure: oblique two-vehicle collision momentum vector diagram. Car A
% travels east, car B travels north; they lock and move off at a common
% heading. The vector triangle shows total momentum conserved.
\begin{figure}[htbp]
\centering
\begin{tikzpicture}[
  pvec/.style={draw=engblue, -{Stealth}, very thick},
  rvec/.style={draw=engred, -{Stealth}, very thick},
  lbl/.style={font=\scriptsize},
]
% scale: 1 unit = 10000 kg m/s.
% Car A: m=1400, v=14 east  -> pA = 19600 -> 1.96 units east
% Car B: m=1000, v=18 north -> pB = 18000 -> 1.80 units north
% total p = (1.96, 1.80), magnitude 2.66 units = 26600 kg m/s
% origin
\coordinate (O) at (0,0);
% pA along x
\draw[pvec] (O) -- (3.92,0) node[midway, below, lbl] {$p_A=19\,600$};
% pB along y from tip of pA (vector triangle)
\draw[pvec] (3.92,0) -- (3.92,3.60) node[midway, right, lbl] {$p_B=18\,000$};
% resultant from O to tip
\draw[rvec] (O) -- (3.92,3.60) node[midway, above left, lbl, engred]
  {$P=26\,600$};
% angle annotation
\node[lbl, anchor=west] at (1.4,1.0) {$\phi=42.6^{\circ}$};
% small axis compass
\draw[->, enggray] (-0.6,-0.6) -- (-0.6,0.2) node[above, lbl, enggray] {N};
\draw[->, enggray] (-0.6,-0.6) -- (0.2,-0.6) node[right, lbl, enggray] {E};
% labels for the two incoming cars (schematic)
\node[lbl, engblue, anchor=north] at (1.96,-0.7) {car A: east};
\node[lbl, engblue, anchor=west] at (4.05,1.8) {car B: north};
% post-impact wreck heading note
\node[lbl, engred, anchor=west] at (2.3,3.0) {wreck heading};
\end{tikzpicture}
\caption[Oblique two-vehicle collision: momentum vector triangle]{The
momentum vector triangle for the oblique collision in the
accident-reconstruction case study. Car A ($1{,}400\,\si{\kilogram}$)
travels east at $14\,\si{\meter\per\second}$, contributing
$p_A=19{,}600\,\si{\kilogram\meter\per\second}$; car B
($1{,}000\,\si{\kilogram}$) travels north at $18\,\si{\meter\per
\second}$, contributing $p_B=18{,}000\,\si{\kilogram\meter\per
\second}$. The two perpendicular momenta add head to tail, so the
locked wreck departs along the resultant $P=\sqrt{p_A^{2}+p_B^{2}}=
26{,}600\,\si{\kilogram\meter\per\second}$ at heading
$\phi=\arctan(p_B/p_A)=42.6^{\circ}$ north of east. The post-impact
skid direction recorded at the scene is the physical check on this
heading, and a mismatch is the first sign that a stated incoming speed
is wrong.}
\label{fig:vol03ch04:oblique-collision}
\end{figure}


% Figure: free-body diagram of a car on a banked curve. Normal force,
% weight, friction, and the required centripetal direction.
\begin{figure}[htbp]
\centering
\begin{tikzpicture}[
  force/.style={draw=engblue, -{Stealth}, very thick},
  wt/.style={draw=engblack, -{Stealth}, very thick},
  fr/.style={draw=engamber, -{Stealth}, very thick},
  cen/.style={draw=engred, -{Stealth}, thick, dashed},
  lbl/.style={font=\scriptsize},
]
% bank angle 18 degrees. Road surface line.
\def\bank{18}
% ground / banked road
\draw[draw=enggray, thick] (-3,0) -- (3,1.949);  % tan(18)=0.3249, slope over 6 -> 1.949
\fill[enggray!12] (-3,0) -- (3,1.949) -- (3,-0.4) -- (-3,-0.4) -- cycle;
% car block sitting on the bank near centre, rotated
\begin{scope}[shift={(0,0.6494)}, rotate=\bank]
  \draw[draw=engblack, thick, fill=engblue!10, rounded corners=1.5pt]
    (-0.7,0) rectangle (0.7,0.45);
\end{scope}
% point of application = centre of car base on the road
\coordinate (P) at (0,0.6494);
% Normal force: perpendicular to road surface (road slopes up-right at 18 deg,
% so normal points up-left at 18 deg from vertical): direction (−sin18, cos18)
\draw[force] (P) -- ++({-1.6*sin(\bank)}, {1.6*cos(\bank)})
  node[above, lbl, engblue] {$N$};
% Weight straight down
\draw[wt] (P) -- ++(0,-1.6) node[below, lbl] {$mg$};
% Friction up the slope (toward the high side) along surface: (cos18, sin18)
\draw[fr] (P) -- ++({1.3*cos(\bank)}, {1.3*sin(\bank)})
  node[right, lbl, engamber] {$f$ (up-slope)};
% required centripetal direction: horizontal, toward centre of curve (left)
\draw[cen] (P) -- ++(-1.8,0) node[left, lbl, engred] {$a_c=v^2/r$ (toward centre)};
% bank angle arc
\draw[draw=enggray] (1.2,0.78) arc (\bank:0:1.2);
\node[lbl, enggray] at (1.5,0.45) {$\theta$};
\end{tikzpicture}
\caption[Free-body diagram of a car on a banked curve]{Free-body
diagram of a car cornering on a road banked at angle $\theta$. Three
forces act: the weight $mg$ straight down, the normal force $N$
perpendicular to the road surface, and the tyre friction $f$ in the
plane of the road. The acceleration is horizontal and points toward
the centre of the curve with magnitude $v^{2}/r$ (red dashed). Banking
tilts the normal force so that its horizontal component supplies part
of the centripetal requirement without help from friction; at the
design speed $v_{0}=\sqrt{g r\tan\theta}$ the friction is zero and the
normal force alone turns the car. Above that speed friction must point
down-slope to add inward force; below it friction must point up-slope
to keep the car from sliding down the bank. The friction-limited
maximum speed adds $\mu N$ on the inward side and is worked in the text.}
\label{fig:vol03ch04:banked-curve}
\end{figure}


% Figure: vertical loop-the-loop. A cart released from height h on a ramp
% enters a circular loop of radius R; minimum release height for the cart
% to maintain contact at the top is h = 2.5 R.
\begin{figure}[htbp]
\centering
\begin{tikzpicture}[
  track/.style={draw=engblack, very thick},
  hvec/.style={draw=engblue, -{Stealth}, thick},
  fvec/.style={draw=engred, -{Stealth}, thick},
  lbl/.style={font=\scriptsize},
]
% loop centre and radius
\def\R{1.3}
\coordinate (C) at (0,\R);
% circular loop
\draw[track] (C) circle (\R);
% entry ramp on the left, descending from release height h = 2.5 R
\def\hh{3.25} % 2.5 * 1.3
\draw[track] (-4,\hh) .. controls (-2.6,1.6) and (-2.0,0.05) .. (-\R,\R);
% ground line at base of loop
\draw[draw=enggray] (-4.2,0) -- (1.6,0);
% release height marker
\draw[draw=enggray, dashed] (-4,\hh) -- (1.4,\hh);
\draw[<->, enggray] (1.4,0) -- (1.4,\hh) node[midway, right, lbl, enggray] {$h=2.5R$};
% loop radius marker
\draw[<->, enggray] (C) -- ++(0,\R) node[midway, right, lbl, enggray] {$R$};
% cart at the top of the loop
\fill[engblue] (0,2*\R) circle (2.5pt);
\node[lbl, anchor=south] at (0,2*\R+0.08) {top: $v_{\text{top}}^2=gR$};
% forces at the top: weight down, normal down (toward centre)
\draw[fvec] (0,2*\R) -- ++(0,-0.7) node[right, lbl, engred] {$mg+N$};
% cart at release
\fill[engblue] (-4,\hh) circle (2.5pt);
\node[lbl, anchor=west] at (-3.95,\hh+0.1) {release, $v=0$};
\end{tikzpicture}
\caption[Vertical loop-the-loop and the minimum release height]{A cart
released from rest at height $h$ runs down a frictionless ramp and
into a vertical circular loop of radius $R$. At the top of the loop
the cart needs a downward centripetal force $mv_{\text{top}}^{2}/R$
supplied by gravity plus the track's normal force, $mg+N=mv_{\text{
top}}^{2}/R$. Contact is maintained only while $N\ge 0$, so the
slowest survivable top speed is set by $N=0$, giving
$v_{\text{top}}^{2}=gR$. The work-energy theorem then fixes the
minimum release height: equating $\tfrac{1}{2}mv_{\text{top}}^{2}+
mg(2R)$ at the top to $mgh$ at release gives $h=2.5R$, independent of
the cart's mass. A cart released any lower leaves the track before the
top and falls inward, the failure the worked example computes.}
\label{fig:vol03ch04:loop}
\end{figure}


\section{End-to-end case study: a two-stage launch
budget}\pathtag{standard}

A launch to low Earth orbit exercises the variable-mass form of
Newton's second law over the whole flight, and it is the cleanest
setting in which to watch the rocket equation, the staging argument,
and the loss budget interact. The task here is to size a two-stage
vehicle that delivers a $1{,}000\,\si{\kilogram}$ payload to a
low-orbit speed of $7.8\,\si{\kilo\meter\per\second}$, and to show why
a single stage cannot do the job with the same propellant technology.
The reasoning runs from the required velocity through the rocket
equation to the stage masses, then back through the gravity and drag
losses that the ideal velocity must also cover. Numbers are carried at
three figures; the exhaust velocities are representative of kerosene
and hydrogen upper stages rather than any one flown vehicle.

\textbf{Step 1: the velocity the vehicle must supply.} Orbital speed at
low altitude follows from setting gravity equal to the centripetal
requirement, $g=v^{2}/r$, with $r$ the orbital radius. For a low orbit
$v_{\text{orb}}\approx 7.8\,\si{\kilo\meter\per\second}$. The vehicle
must supply more than this, because some of its ideal velocity is spent
fighting gravity during the climb and air during the ascent. The
\engterm{gravity loss} is the time integral of $g\sin\gamma$ over the
boost, with $\gamma$ the flight-path angle above horizontal; a vehicle
that climbs steeply and slowly loses the most, and a typical figure for
a well-flown ascent is $1.2$ to $1.5\,\si{\kilo\meter\per\second}$. The
\engterm{drag and steering loss} adds another $0.3$ to
$0.5\,\si{\kilo\meter\per\second}$. Taking $1.3$ and
$0.5\,\si{\kilo\meter\per\second}$ for the two, the ideal velocity the
propulsion must deliver is
\[
\Delta v_{\text{ideal}}=v_{\text{orb}}+\Delta v_{\text{grav}}
+\Delta v_{\text{drag}}=7.8+1.3+0.5=9.6\,\si{\kilo\meter\per\second}.
\]
Figure~\ref{fig:vol03ch04:deltav-budget} draws this budget: the useful
orbital speed is the largest single block, and the $1.8\,\si{\kilo\meter
\per\second}$ of loss is the surcharge for launching from inside an
atmosphere and a gravity field.

\textbf{Step 2: why one stage fails.} The rocket equation $\Delta
v=v_{e}\ln(m_{0}/m_{f})$ sets the achievable velocity from the mass
ratio. A good kerosene first stage has exhaust velocity $v_{e}\approx
3.0\,\si{\kilo\meter\per\second}$ (specific impulse near
$300\,\si{\second}$). For a single stage to reach
$9.6\,\si{\kilo\meter\per\second}$ the required mass ratio is
\[
\frac{m_{0}}{m_{f}}=\exp\!\left(\frac{\Delta v_{\text{ideal}}}{v_{e}}
\right)=\exp\!\left(\frac{9.6}{3.0}\right)=\exp(3.20)=24.5.
\]
The final mass $m_{f}$ must hold the structure, the engines, the empty
tanks, and the payload. A mass ratio of $24.5$ leaves a final mass of
$1/24.5=4.1\,\%$ of the lift-off mass for all of that. Real tankage and
engines for a kerosene stage run a structural fraction near $6$ to
$10\,\%$ of the propellant loaded, so $4.1\,\%$ for structure plus
payload is not reachable: the tanks alone weigh more than the budget
allows, and there is nothing left for payload. The single stage is
defeated by its own dead weight, which it must carry to orbit along
with the payload.

\textbf{Step 3: staging splits the burden.} A two-stage vehicle drops
the first stage's empty mass before lighting the second, so the second
stage never has to accelerate the first stage's spent tanks. Split the
$9.6\,\si{\kilo\meter\per\second}$ as $\Delta v_{1}=4.2$ and $\Delta
v_{2}=5.4\,\si{\kilo\meter\per\second}$, giving the upper stage the
larger share because it carries the payload the farthest. Use $v_{e,1}=
3.0$ for the kerosene first stage and $v_{e,2}=4.4\,\si{\kilo\meter\per
\second}$ for a hydrogen upper stage (specific impulse near
$450\,\si{\second}$). The stage mass ratios are
\[
\frac{m_{0,1}}{m_{f,1}}=\exp\!\left(\frac{4.2}{3.0}\right)=\exp(1.40)=
4.06,\qquad
\frac{m_{0,2}}{m_{f,2}}=\exp\!\left(\frac{5.4}{4.4}\right)=\exp(1.23)=
3.41.
\]
Each ratio is now modest, well inside what real tankage can carry.
Work the upper stage first because it sets the payload it must deliver.
Take a structural fraction $\varepsilon=0.10$ for each stage (empty
mass is $10\,\%$ of the loaded propellant-plus-structure). For the
upper stage carrying the $1{,}000\,\si{\kilogram}$ payload, write
$m_{0,2}=m_{\text{prop,2}}+m_{\text{struct,2}}+m_{\text{pay}}$ and
$m_{f,2}=m_{\text{struct,2}}+m_{\text{pay}}$ with
$m_{\text{struct,2}}=\varepsilon\,(m_{\text{prop,2}}+m_{\text{struct,2}})$.
The mass-ratio condition $m_{0,2}/m_{f,2}=3.41$ with $\varepsilon=0.10$
solves to an upper-stage gross mass near $5.0\,\si{\tonne}$, of which
about $3.5\,\si{\tonne}$ is propellant, $0.4\,\si{\tonne}$ is structure,
and $1.0\,\si{\tonne}$ is payload. The first stage then treats the
fuelled upper stage as its payload ($5.0\,\si{\tonne}$) and the same
construction with $m_{0,1}/m_{f,1}=4.06$ gives a first-stage gross mass
near $26\,\si{\tonne}$. The full vehicle leaves the pad at about
$26\,\si{\tonne}$ to place $1\,\si{\tonne}$ in orbit, a payload
fraction near $3.8\,\%$.

\textbf{Step 4: the cross-check on the loss budget.} A second route to
the same answer checks the assumed losses. The gravity loss equals
$\int g\sin\gamma\,dt$, and for a vehicle whose thrust-to-weight ratio
at lift-off is $T/W=1.3$, the early vertical climb lasts until the
pitch-over, a few tens of seconds. With a boost phase of order
$150\,\si{\second}$ and an average $\sin\gamma$ near $0.6$ over the
climb, the gravity loss estimate is $g\,\overline{\sin\gamma}\,t\approx
9.81(0.6)(150)/1000\approx 0.88$ to $1.5\,\si{\kilo\meter\per\second}$
depending on how steeply the trajectory is flown. The figure used in
Step 1, $1.3\,\si{\kilo\meter\per\second}$, sits in this band, so the
budget is internally consistent: the loss assumed at the start is
recovered from the flight time and climb angle the staging implies. A
higher lift-off thrust-to-weight shortens the boost and shrinks the
gravity loss, which is why launch vehicles run $T/W>1.2$ rather than
lifting off sluggishly.

\textbf{What the case study shows.} Three results survive the
arithmetic. The first is the dominance of the exponential: each
kilometre per second of velocity multiplies the mass ratio, so the
loss budget is not a rounding term but the difference between a
reachable and an unreachable design. The second is the power of
staging: dropping dead mass partway up converts one impossible mass
ratio into two easy ones, which is why every orbital launcher stages.
The third is the smallness of the payload fraction: under four per cent
of the lift-off mass reaches orbit, and the rest is propellant and the
structure to hold it, a ratio that has shaped launch economics since
the first orbital flights.

\begin{principle}[The rocket equation is multiplicative in velocity and
exponential in mass]
A velocity requirement adds, but the mass it demands multiplies. Two
velocity increments $\Delta v_{1}$ and $\Delta v_{2}$ that sum to a
total cost a mass ratio $\exp((\Delta v_{1}+\Delta v_{2})/v_{e})=
\exp(\Delta v_{1}/v_{e})\exp(\Delta v_{2}/v_{e})$, the product of the
two separate ratios. Every loss term, every margin, every reserve adds
to the velocity and so multiplies the lift-off mass, which is why
launch-vehicle design fights for each metre per second and why the
useful payload is a thin slice of the whole. Staging breaks the single
exponential into a product of smaller ones by discarding spent mass, a
move the equation rewards because the discarded structure no longer
rides the later increments.
\end{principle}

% Figure: the delta-v budget for a two-stage launch to low Earth orbit,
% drawn as a stacked accounting of the ideal (Tsiolkovsky) delta-v
% against the losses (gravity, drag, steering) that the launch must
% also pay for, and the orbital velocity that is the useful product.
\begin{figure}[htbp]
\centering
\begin{tikzpicture}
\begin{axis}[
  width=12cm, height=7cm,
  xbar stacked,
  xmin=0, xmax=10,
  ymin=-0.6, ymax=1.6,
  xlabel={velocity budget (\si{\kilo\meter\per\second})},
  ytick={0,1},
  yticklabels={delivered, required},
  axis lines=left,
  bar width=20pt,
  tick label style={font=\scriptsize},
  label style={font=\small},
  legend style={font=\scriptsize, at={(0.5,-0.22)}, anchor=north, legend columns=2},
  legend cell align=left,
  enlarge y limits=0.4,
]
% "required" row (y=1): ideal delta-v split between the two stages.
% Stage 1 ideal 4.2, stage 2 ideal 5.4, total ideal 9.6 km/s.
\addplot[fill=engblue!70, draw=engblue] coordinates {(4.2,1)};
\addplot[fill=englime!70, draw=englime!60!black] coordinates {(5.4,1)};
% "delivered" row (y=0): the 7.8 km/s orbital speed plus the losses.
\addplot[fill=enggray!45, draw=enggray] coordinates {(7.8,0)};
\addplot[fill=engamber!75, draw=engamber] coordinates {(1.3,0)};
\addplot[fill=engred!65, draw=engred] coordinates {(0.5,0)};
\legend{stage 1 ideal $\Delta v$, stage 2 ideal $\Delta v$, orbital
  speed (useful), gravity loss, drag and steering loss}
\end{axis}
\end{tikzpicture}
\caption[Two-stage launch velocity budget]{The velocity budget for a
two-stage launch to low Earth orbit. The lower bar (delivered) shows
where the ideal velocity goes: about $7.8\,\si{\kilo\meter\per\second}$
becomes useful orbital speed, while roughly $1.3\,\si{\kilo\meter\per
\second}$ is spent against gravity during the climb and about
$0.5\,\si{\kilo\meter\per\second}$ against atmospheric drag and the
steering turn. The upper bar (required) shows the same total split
between the two stages by the Tsiolkovsky equation: the first stage
supplies about $4.2\,\si{\kilo\meter\per\second}$ and the second about
$5.4\,\si{\kilo\meter\per\second}$, summing to the $9.6\,\si{\kilo\meter
\per\second}$ ideal that the orbital speed plus the losses demand. The
gap between the $7.8\,\si{\kilo\meter\per\second}$ useful product and
the $9.6\,\si{\kilo\meter\per\second}$ paid is the cost of leaving the
ground inside an atmosphere and a gravity field.}
\label{fig:vol03ch04:deltav-budget}
\end{figure}


\begin{mastery}
\textbf{Impulse-momentum for variable-mass systems, beyond the
rocket.} The momentum form $\vec{F}=d\vec{p}/dt$ governs any system
that gains or sheds mass, and the rocket is one instance of a wider
pattern. The general statement, for a main body of mass $m$ exchanging
mass with its surroundings, is
\[
m\frac{d\vec{v}}{dt}=\vec{F}_{\text{ext}}+\dot{m}_{\text{in}}
(\vec{u}_{\text{in}}-\vec{v})-\dot{m}_{\text{out}}(\vec{u}_{\text{out}}
-\vec{v}),
\]
where $\vec{u}_{\text{in}}$ and $\vec{u}_{\text{out}}$ are the
velocities at which mass joins and leaves and $\vec{v}$ is the body's
velocity. The bracketed terms are the relative velocities of the
incoming and outgoing streams; each carries momentum across the body's
boundary and acts as a thrust or a drag. Three engineering systems
read straight off this equation. A rocket has $\dot{m}_{\text{in}}=0$
and an exhaust stream leaving at relative velocity $-v_{e}$, recovering
$m\dot{v}=\dot{m}_{\text{out}}v_{e}$. A conveyor belt loading sand
from a stationary hopper has $\vec{u}_{\text{in}}=\vec{0}$, so the belt
motor must supply an extra force $\dot{m}_{\text{in}}v$ to accelerate
the newly-arrived material to belt speed, a force present even at
constant belt speed and absent from a naive $\sum F=ma$ written for
constant mass. A raindrop growing as it falls through a stationary mist
has $\vec{u}_{\text{in}}=\vec{0}$ and a downward $\vec{v}$, and the
mass it sweeps up retards it, so the drop accelerates more slowly than
$g$. The error of writing $\sum F=ma$ for these systems is to forget
that the mass crossing the boundary arrives or departs with a momentum
of its own; the relative-velocity terms are the bookkeeping that
forgetting omits. Meriam and Kraige treat the variable-mass equation
and its three regimes at problem-set depth
\cite{text:meriam-kraige2018}.
\end{mastery}

\section{End-to-end case study: a Hohmann transfer and
rendezvous}\pathtag{standard}

A spacecraft that must move from one circular orbit to a higher one,
and meet a target already on the higher orbit, exercises the
conservation laws, the central-force results, and the velocity-budget
discipline of the launch case study, now applied to motion that is all
coasting between two short burns. The task is to raise a chaser from a
low parking orbit to a higher circular orbit, compute the two burns and
the transfer time, and then find the phasing the rendezvous demands.
The vehicle is in orbit about Earth; the gravitational parameter is
$\mu=GM_{\oplus}=3.986\times 10^{14}\,\si{\meter\cubed\per\second
\squared}$, and Earth's radius is $R_{\oplus}=6{,}371\,\si{\kilo\meter}$.
The chaser starts on a circular orbit at altitude $300\,\si{\kilo
\meter}$ ($r_{1}=6{,}671\,\si{\kilo\meter}$); the target circles at
altitude $1{,}000\,\si{\kilo\meter}$ ($r_{2}=7{,}371\,\si{\kilo\meter}$).
Numbers are carried at four figures.

\textbf{Step 1: the two circular speeds.} A circular orbit balances
gravity against the centripetal requirement, $\mu/r^{2}=v^{2}/r$, so
the circular speed is $v_{\text{circ}}=\sqrt{\mu/r}$, the central-force
result of the worked example specialised to a circle. At the two radii,
\[
v_{1}=\sqrt{\frac{3.986\times 10^{14}}{6.671\times 10^{6}}}=7{,}730\,\si{
\meter\per\second}, \quad
v_{2}=\sqrt{\frac{3.986\times 10^{14}}{7.371\times 10^{6}}}=7{,}353\,\si{
\meter\per\second}.
\]
The higher orbit is the slower one, the counterintuitive signature of
orbital motion: climbing to a higher orbit ends with the craft moving
more slowly, not faster, because the speed it gains in the burns is
spent climbing against gravity and then some.

\textbf{Step 2: the transfer ellipse.} The minimum-energy two-burn
path between the two circular orbits is a Hohmann transfer ellipse whose
periapsis touches the inner orbit and whose apoapsis touches the outer
(Figure~\ref{fig:vol03ch04:hohmann}). Its semi-major axis is the
average of the two radii,
\[
a_{t}=\frac{r_{1}+r_{2}}{2}=\frac{6{,}671+7{,}371}{2}=7{,}021\,\si{\kilo
\meter}.
\]
The speed at any point on the ellipse follows from the
\engterm{vis-viva} relation $v^{2}=\mu(2/r-1/a)$, itself the energy
conservation law $\tfrac{1}{2}v^{2}-\mu/r=-\mu/(2a)$ rearranged. At
periapsis ($r=r_{1}$) and apoapsis ($r=r_{2}$) of the transfer ellipse,
\[
v_{p}=\sqrt{\mu\Bigl(\frac{2}{r_{1}}-\frac{1}{a_{t}}\Bigr)}=\sqrt{3.986
\times 10^{14}\Bigl(\frac{2}{6.671\times 10^{6}}-\frac{1}{7.021\times
10^{6}}\Bigr)}=7{,}926\,\si{\meter\per\second},
\]
\[
v_{a}=\sqrt{\mu\Bigl(\frac{2}{r_{2}}-\frac{1}{a_{t}}\Bigr)}=\sqrt{3.986
\times 10^{14}\Bigl(\frac{2}{7.371\times 10^{6}}-\frac{1}{7.021\times
10^{6}}\Bigr)}=7{,}173\,\si{\meter\per\second}.
\]

\textbf{Step 3: the two burns.} Each burn is the speed change at a
tangent point, a scalar difference because the transfer ellipse is
tangent to both circles and the velocities are parallel there:
\[
\Delta v_{1}=v_{p}-v_{1}=7{,}926-7{,}730=196\,\si{\meter\per\second},
\]
\[
\Delta v_{2}=v_{2}-v_{a}=7{,}353-7{,}173=180\,\si{\meter\per\second}.
\]
The first burn speeds the chaser up at periapsis to stretch the orbit
out to the higher apoapsis; the second burn speeds it up again at
apoapsis to round the slow apoapsis speed up to the higher orbit's
circular speed. Both burns are prograde. The total is $\Delta v_{\text{
total}}=\Delta v_{1}+\Delta v_{2}=376\,\si{\meter\per\second}$, the
velocity budget for the whole manoeuvre, which the rocket equation of
the launch case study converts into a propellant mass: a chaser of dry
mass $500\,\si{\kilogram}$ burning a thruster of exhaust velocity
$v_{e}=3{,}000\,\si{\meter\per\second}$ needs a mass ratio $\exp(376/
3000)=1.134$, so about $67\,\si{\kilogram}$ of propellant, $13\,\%$ of
the dry mass, a modest budget that makes the Hohmann transfer the
default for orbit raising when time is not pressing.

\textbf{Step 4: the transfer time.} The coast is half the period of
the transfer ellipse, and the period follows from Kepler's third law in
the form $T=2\pi\sqrt{a^{3}/\mu}$:
\[
t_{\text{transfer}}=\frac{T_{t}}{2}=\pi\sqrt{\frac{a_{t}^{3}}{\mu}}
=\pi\sqrt{\frac{(7.021\times 10^{6})^{3}}{3.986\times 10^{14}}}
=2{,}924\,\si{\second},
\]
about $48.7\,\si{\minute}$. The chaser fires the first burn, coasts
unpowered for forty-nine minutes along the half-ellipse, and fires the
second burn at arrival. No force acts during the coast except gravity,
so the long quiet middle of the manoeuvre is pure central-force motion
of the kind the orbit worked example treated.

\textbf{Step 5: the phasing for rendezvous.} Reaching the higher orbit
is not the same as meeting the target on it. The chaser arrives at the
apoapsis point after the coast; the target must be at that same point at
that same moment (Figure~\ref{fig:vol03ch04:phasing}). The target's
angular rate on its circular orbit is
\[
\omega_{\text{target}}=\frac{v_{2}}{r_{2}}=\frac{7{,}353}{7.371\times
10^{6}}=9.976\times 10^{-4}\,\si{\radian\per\second}.
\]
During the transfer the target sweeps an angle $\omega_{\text{target}}
\,t_{\text{transfer}}=9.976\times 10^{-4}\times 2{,}924=2.917\,\si{
\radian}=167.1^{\circ}$. The chaser sweeps exactly $180^{\circ}$ along
the half-ellipse from periapsis to apoapsis. For both to reach the
apoapsis point together, the target must lead the chaser's arrival point
at departure by the angle the chaser's $180^{\circ}$ exceeds the
target's sweep,
\[
\phi_{\text{lead}}=180^{\circ}-167.1^{\circ}=12.9^{\circ},
\]
measured ahead of the chaser's arrival point in the direction of
motion. The chaser may burn only when the target sits $12.9^{\circ}$
ahead; at any other geometry the two miss, and the chaser must wait for
the phase to come around, the orbital analogue of waiting for a launch
window. The wait between successive windows is set by how fast the two
orbits drift in phase, the \engterm{synodic period}
$T_{\text{syn}}=2\pi/|\omega_{1}-\omega_{2}|$ built from the two orbital
rates.

\textbf{The cross-check on the energy.} A second route confirms the
burns through the orbital energy. The specific orbital energy is
$\varepsilon=-\mu/(2a)$, so the chaser's energy rises from $-\mu/(2r_{1})
=-2.988\times 10^{7}\,\si{\joule\per\kilogram}$ on the inner circle to
$-\mu/(2r_{2})=-2.704\times 10^{7}\,\si{\joule\per\kilogram}$ on the
outer, a gain of $2.84\times 10^{6}\,\si{\joule\per\kilogram}$. The two
burns must supply exactly this specific-energy gain through the work
their thrust does, and computing $\tfrac{1}{2}(v_{1}+\Delta v_{1})^{2}-
\tfrac{1}{2}v_{1}^{2}$ at periapsis plus the matching apoapsis term
returns the same energy gain to within rounding, the independent check
that the speeds and the budget hold together.

\begin{principle}[Orbit changes are bought in velocity, not distance]
The cost of moving between orbits is set by the velocity changes at the
burn points, not by the distance climbed. A Hohmann transfer charges
two scalar burns at the tangent points and asks nothing of the long
coast between them, because gravity does the work of the climb and the
craft pays only to enter and leave the transfer ellipse. The vis-viva
relation prices each burn from the orbital energies, and the phasing
condition turns the geometry into a schedule: reaching the orbit and
meeting a body on it are separate requirements, the first set by
velocity and the second by timing.
\end{principle}

% Figure: a Hohmann transfer between two coplanar circular orbits, with
% the two impulsive burns at the tangent points and the transfer ellipse
% touching the inner orbit at periapsis and the outer orbit at apoapsis.
\begin{figure}[htbp]
\centering
\begin{tikzpicture}[
  burn/.style={-{Stealth[length=3mm]}, very thick, draw=engred},
  lbl/.style={font=\small},
]
% central body
\fill[engamber] (0,0) circle (5pt);
\node[engamber, lbl, below left] at (0,0) {Earth};
% inner circular orbit, radius 1.7
\draw[engblue, very thick] (0,0) circle (1.7);
\node[engblue, lbl] at (-1.7,-1.35) {inner orbit $r_1$};
% outer circular orbit, radius 3.4
\draw[enggray, very thick] (0,0) circle (3.4);
\node[enggray, lbl] at (2.4,2.7) {outer orbit $r_2$};
% transfer ellipse: periapsis at inner orbit (left, x=-1.7),
% apoapsis at outer orbit (right, x=+3.4).
% semi-major a = (1.7+3.4)/2 = 2.55; centre offset along x:
% centre x = apoapsis - a = 3.4 - 2.55 = 0.85.
% b = sqrt(r1*r2) = sqrt(1.7*3.4) = sqrt(5.78) = 2.404.
\draw[englime!70!black, very thick, dashed] (0.85,0) ellipse (2.55 and 2.404);
\node[englime!60!black, lbl] at (0.85,-2.7) {transfer ellipse};
% burn 1 at periapsis (left tangent point on inner orbit), x=-1.7,y=0
\coordinate (b1) at (-1.7,0);
\fill[engblack] (b1) circle (1.6pt);
\draw[burn] (b1) ++(0,0) -- ++(0,0.9);
\node[engred, lbl, left] at (-1.7,0.5) {$\Delta v_1$};
% burn 2 at apoapsis (right tangent point on outer orbit), x=+3.4,y=0
\coordinate (b2) at (3.4,0);
\fill[engblack] (b2) circle (1.6pt);
\draw[burn] (b2) ++(0,0) -- ++(0,0.9);
\node[engred, lbl, right] at (3.4,0.5) {$\Delta v_2$};
% direction-of-motion arrows
\draw[-{Stealth[length=2.5mm]}, engblue, thick] (0,-1.7) ++(0.25,0) arc (0:-22:1.7);
\draw[-{Stealth[length=2.5mm]}, enggray, thick] (0,3.4) ++(-0.25,0) arc (90:112:3.4);
\end{tikzpicture}
\caption[Hohmann transfer between two circular orbits]{A Hohmann
transfer raises a spacecraft from an inner circular orbit of radius
$r_1$ to an outer circular orbit of radius $r_2$ with two prograde
burns. The first burn $\Delta v_1$ at periapsis injects the craft onto
an ellipse whose periapsis touches the inner orbit and whose apoapsis
touches the outer orbit; the craft coasts half an ellipse to apoapsis,
where the second burn $\Delta v_2$ circularises onto the outer orbit.
The transfer ellipse has semi-major axis $a=(r_1+r_2)/2$, and the
coast lasts half its period. The Hohmann transfer is the
minimum-energy two-impulse transfer between coplanar circular orbits,
which is why it sets the reference $\Delta v$ against which other
transfer schemes are measured.}
\label{fig:vol03ch04:hohmann}
\end{figure}


% Figure: the rendezvous phasing geometry for a Hohmann transfer. The
% chaser leaves the inner orbit when the target on the outer orbit leads
% by the lead angle, so that target and chaser arrive at the apoapsis
% point together after the half-ellipse coast.
\begin{figure}[htbp]
\centering
\begin{tikzpicture}[lbl/.style={font=\small}]
% central body
\fill[engamber] (0,0) circle (4pt);
% inner orbit r=1.6 (chaser start), outer orbit r=3.2 (target)
\draw[engblue, thick] (0,0) circle (1.6);
\draw[enggray, thick] (0,0) circle (3.2);
% chaser starts at periapsis point, left (angle 180), x=-1.6
\coordinate (chaser) at (180:1.6);
\fill[engblue] (chaser) circle (2.2pt);
\node[engblue, lbl, left] at (chaser) {chaser};
% arrival point is apoapsis, right (angle 0), x=+3.2
\coordinate (arrive) at (0:3.2);
\fill[engblack] (arrive) circle (1.6pt);
\node[lbl, right] at (arrive) {rendezvous};
% target leads: it must be at lead angle ahead of the arrival point so
% it reaches arrival as the chaser does. Draw target on outer orbit at
% angle 50 deg (measured from +x toward the leading direction).
\coordinate (target) at (50:3.2);
\fill[engred] (target) circle (2.2pt);
\node[engred, lbl, above right] at (target) {target};
% lead angle arc between arrival (0 deg) and target (50 deg)
\draw[-{Stealth[length=2.5mm]}, engred, thick] (0:1.05) arc (0:50:1.05);
\node[engred, lbl] at (25:1.35) {$\phi_{\text{lead}}$};
% transfer ellipse (dashed) periapsis at -1.6, apoapsis at +3.2
% a=(1.6+3.2)/2=2.4, centre x = 3.2-2.4=0.8, b=sqrt(1.6*3.2)=sqrt(5.12)=2.263
\draw[englime!70!black, very thick, dashed] (0.8,0) ellipse (2.4 and 2.263);
% direction of motion
\draw[-{Stealth[length=2.2mm]}, enggray] (90:3.2) ++(-0.2,0) arc (90:108:3.2);
\end{tikzpicture}
\caption[Rendezvous phasing for a Hohmann transfer]{Phasing for a
rendezvous by Hohmann transfer. The chaser leaves the inner orbit at
periapsis of the transfer ellipse and arrives at apoapsis after a
half-orbit coast. For the target to be at the apoapsis point at that
moment, it must lead the chaser's arrival point by the angle it sweeps
during the coast, $\phi_{\text{lead}}=\pi-\omega_{\text{target}}\,
t_{\text{transfer}}$, where $\omega_{\text{target}}$ is the target's
angular rate and $t_{\text{transfer}}$ is the half-period of the
transfer ellipse. The required lead angle sets the launch or departure
window: the chaser may burn only when the target sits at exactly this
angle ahead.}
\label{fig:vol03ch04:phasing}
\end{figure}


\begin{mastery}
\textbf{Kepler's laws from the inverse-square force.} The three laws
Kepler read off Tycho Brahe's data are consequences of the
inverse-square central force, and deriving them shows how two
conservation laws fix the shape of every bound orbit. Start from the
force $\vec{F}=-\mu m\hat{r}/r^{2}$ on a body of mass $m$ about a fixed
centre. Because the force is central it exerts no moment about the
centre, so the angular momentum $L=m r^{2}\dot{\theta}$ is conserved.
The areal sweep rate is $\tfrac{1}{2}r^{2}\dot{\theta}=L/(2m)$, constant
in time: the radius vector sweeps equal areas in equal times, which is
\engterm{Kepler's second law}, and it holds for any central force, not
only the inverse square.

The first and third laws need the inverse-square form. Write the radial
equation of motion $\ddot{r}-r\dot{\theta}^{2}=-\mu/r^{2}$ and change
the variable to $u=1/r$ with $\theta$ as the independent variable. Using
$\dot{\theta}=Lu^{2}/m$ to convert time derivatives, the radial equation
becomes the linear oscillator equation
\[
\frac{d^{2}u}{d\theta^{2}}+u=\frac{\mu m^{2}}{L^{2}},
\]
whose solution is $u=\frac{\mu m^{2}}{L^{2}}\bigl(1+e\cos\theta\bigr)$,
that is $r=\dfrac{p}{1+e\cos\theta}$ with semi-latus rectum $p=L^{2}/
(\mu m^{2})$ and eccentricity $e$ fixed by the initial energy. That
$r(\theta)$ is the polar equation of a conic with the centre at one
focus: an ellipse for $e<1$, the bound case, which is \engterm{Kepler's
first law}. \engterm{Kepler's third law} follows from the second:
integrating the constant areal rate over one period gives the ellipse
area $\pi a b=\frac{L}{2m}T$, and substituting $b=a\sqrt{1-e^{2}}$ and
$p=a(1-e^{2})=L^{2}/(\mu m^{2})$ yields $T^{2}=4\pi^{2}a^{3}/\mu$, the
period squared proportional to the semi-major axis cubed, with a
constant that depends only on the central mass. The Hohmann transfer
time computed above is one use of this last law; the orbit-period data
file of the central-force worked example is another. Landau and
Lifshitz derive the orbit equation and the conic solution from the same
two conserved quantities \cite[\S{}15]{text:landau-mechanics}.
\end{mastery}

\section{Failure: rotating machinery imbalance}\pathtag{core}

Rotating machinery whose mass distribution is not symmetric about
its rotation axis develops dynamic loads on its bearings and
foundations at the rotation frequency. The loads are unwanted (they
have no purpose in the machine's function) and large (they scale as
$\omega^2$, so doubling the rotation speed quadruples the imbalance
force).

For a thin rotor of mass $m$ with centre of mass displaced from the
geometric rotation axis by distance $e$ (the \engterm{eccentricity}),
the imbalance force at rotation rate $\omega$ is
\[
F_{\text{imb}} = m e \omega^2.
\]
The force rotates with the rotor at the rotation rate and produces
a sinusoidal load on the bearings at that frequency. Vibration
sensors on the bearing housings register the force directly; the
amplitude of the rotation-frequency component is the standard
imbalance measurement, and the phase identifies the angular location
of the heavy spot on the rotor.

% Figure: imbalance force F = m e omega^2 versus rotation speed, on log-log
% style linear axes, for several eccentricities, with the omega^2 growth.
\begin{figure}[htbp]
\centering
\begin{tikzpicture}
\begin{axis}[
  width=12cm, height=7cm,
  xlabel={rotation speed $N$ (revolutions per minute)},
  ylabel={imbalance force $F_{\text{imb}}$ (\si{\newton})},
  xmin=0, xmax=6000, ymin=0, ymax=2000,
  axis lines=left,
  domain=0:6000, samples=120,
  legend pos=north west,
  legend cell align=left,
  grid=major, grid style={enggray!20},
  tick label style={font=\scriptsize},
  label style={font=\small},
  legend style={font=\scriptsize},
]
% F = m e omega^2, omega = N 2pi/60. m = 20 kg rotor.
% F = m e (2 pi N /60)^2 = m e (0.10472 N)^2 = m e 0.010966 N^2.
% e = 0.05 mm = 5e-5 m: coeff = 20*5e-5*0.010966 = 1.0966e-5
\addplot[engblue, very thick] {1.0966e-5 * x^2};
\addlegendentry{$e=0.05\,\si{\milli\meter}$ (balanced)}
% e = 0.15 mm: coeff = 20*1.5e-4*0.010966 = 3.290e-5
\addplot[engamber, very thick] {3.290e-5 * x^2};
\addlegendentry{$e=0.15\,\si{\milli\meter}$}
% e = 0.50 mm: coeff = 20*5e-4*0.010966 = 1.0966e-4
\addplot[engred, very thick] {1.0966e-4 * x^2};
\addlegendentry{$e=0.50\,\si{\milli\meter}$ (out of spec)}
\end{axis}
\end{tikzpicture}
\caption[Imbalance force versus rotation speed]{Rotating-imbalance
force $F_{\text{imb}}=m e\,\omega^{2}$ for a $20\,\si{\kilogram}$
rotor at three residual eccentricities, against rotation speed in
revolutions per minute. The quadratic growth in speed is the load
designer's central fact: tripling the speed raises the force ninefold
at fixed eccentricity, and the residual eccentricity allowed by an
ISO~1940/1 balance grade therefore tightens as the rated speed rises.
A rotor balanced to $e=0.05\,\si{\milli\meter}$ (blue) generates a
modest force even at $6{,}000$ revolutions per minute; the same rotor
left at $e=0.50\,\si{\milli\meter}$ (red) drives nearly
$2{,}000\,\si{\newton}$ of rotating bearing load, cyclic at the
rotation frequency, into the bearings, shaft, and foundation.}
\label{fig:vol03ch04:imbalance-force}
\end{figure}


The $\omega^{2}$ scaling is the load designer's central fact.
Figure~\ref{fig:vol03ch04:imbalance-force} plots the force against
speed for a $20\,\si{\kilogram}$ rotor at three residual
eccentricities; tripling the speed raises the force ninefold at fixed
eccentricity, which is why the residual eccentricity that a balance
grade permits tightens as the rated speed rises. The signature in a
vibration spectrum is a dominant line at the rotation frequency, the
\engterm{one-times} line. The sample spectrum in
\texttt{imbalance\_spectrum.csv} shows that line standing an order of
magnitude above its neighbours; a strong two-times line would instead
point to misalignment, and a half-times line to mechanical looseness,
so the spectrum diagnoses the mechanism as well as its severity. Den
Hartog's treatment of machinery vibration and isolation remains the
working practical reference \cite{text:den-hartog-vibrations}; the
allowable residual imbalance for each class of machinery is set by the
balance grades of ISO~1940/1.

For a long rotor (length comparable to or larger than its diameter),
imbalance is not a single eccentric mass but a distribution of
eccentric masses along the rotor's length. The distribution is
characterised by two independent quantities: the static imbalance
(the net eccentric mass at the rotor's centroid plane) and the
couple imbalance (a pair of equal-and-opposite eccentric masses at
different axial positions, producing no static unbalance but a
rotating moment on the bearings). Both must be measured and
corrected. The standard two-plane balancing procedure adds correction
weights to the rotor at two axial planes and at angular positions
determined from the vibration measurements; the procedure is now
automated on dedicated balancing machines, and machine-specific
ISO standards govern the allowable residual imbalance for each
class of rotating machinery (turbines, motors, generators, pumps,
compressors) per ISO 1940/1.

The failure modes of an unbalanced rotor:

\begin{itemize}[leftmargin=*]
  \item \engterm{Bearing wear}. The cyclic radial load accelerates
    bearing fatigue; rolling-element bearings spall and rumble
    audibly; journal bearings develop scoring and increased
    clearance.
  \item \engterm{Shaft fatigue}. The cyclic bending stress in the
    shaft at the imbalance frequency causes fatigue cracks at any
    stress concentration along the shaft (keyways, fillets,
    splines, threaded ends). The shaft eventually fractures at the
    concentration, often during a routine start-up when the
    transient through critical speed amplifies the imbalance force
    momentarily.
  \item \engterm{Foundation resonance}. If the rotation frequency
    matches a natural frequency of the foundation or the machine's
    support structure, the imbalance force is amplified by the
    structural quality factor (often $10$ to $50$). The foundation
    may crack, the supports may loosen, or in severe cases the
    machine may walk on its mounts.
  \item \engterm{Bearing-cap failure}. The bolts holding the bearing
    cap to the machine housing experience the cyclic load; if the
    preload is insufficient or the bolt pattern is undersized, the
    cap loses preload, the bolts fatigue, and the cap separates.
  \item \engterm{Through-housing crack}. In an extreme case, the
    cumulative cyclic load opens fatigue cracks through the
    machine's main housing. The repair is replacement of the
    housing, often a several-month outage for large industrial
    machinery.
\end{itemize}

The mechanisms generalise across the rotating machinery in this book:
turbines (chapter 13 of Volume IV), centrifugal pumps and compressors
(chapter 9 of Volume IV), electric motors and generators (chapter 12
of Volume VIII), engines (chapter 12 of Volume VIII), and the rotors
of helicopters and turbofans (chapter 13 of Volume VIII). In each
domain the imbalance mechanism is the same; the magnitudes and the
acceptable residual imbalance per ISO 1940/1 are different.

Cross-references. Chapter 6 of this volume develops the underlying
vibration theory. Chapter 8 of Volume IV treats critical speed and
rotor-bearing-foundation dynamics. Chapter 2 of Volume X treats the
fatigue mechanism that imbalance forces drive. Chapter 14 of Volume X
treats maintenance failures, of which detected-but-untracked imbalance
trends are the canonical example.

\begin{principle}[Rotating mass distribution is engineering geometry]
For a rotating machine, the location of the centre of mass relative
to the rotation axis is a load-creating geometric property of the
same kind as a force is. The forces induced by imbalance scale as
$\omega^2$, are cyclic at the rotation frequency, and pass directly
into the bearings, shaft, foundation, and housing. Acceptable
residual imbalance is set by ISO 1940/1 for each class of machinery;
designed-in balance and post-build balancing are not optional
features but parts of the machine's design specification.
\end{principle}

\section{Project}\pathtag{core}

\begin{project}[Video-track a thrown object; fit a model]
The reader photographs (with a video camera or smartphone) a thrown
or rolled object in flight, extracts its trajectory frame by frame,
and fits a dynamics model to the data. The deliverable is a
twenty- to thirty-page engineering report.

\textbf{Track}. Physical observation plus analysis.
\textbf{Hazard class}: low. Standard projectiles (tennis ball, soft
ball, sock filled with rice). Throw in an open area; no risk of
injury to bystanders.

\textbf{Apparatus}. A smartphone or video camera capable of
$60\,\text{fps}$ or higher at $1080\text{p}$ resolution; a flat,
visually distinctive background; a measuring tape or known-length
reference object placed in the scene; a chosen projectile.

\textbf{Procedure}. The reader records video of three different
trajectories of the same projectile (varying initial launch angle
or speed). For each video, the reader extracts the projectile's
position from each frame, typically by hand or with a free video-
analysis tool such as Tracker (Brown), and produces a table of
$(x_i, y_i, t_i)$ values. The file
\texttt{data/projectile\_track.csv} gives a worked sample of such a
table (a tennis-ball toss at $60$ frames per second) that the reader
can use to test the analysis pipeline before recording original
video; the file \texttt{code/projectile\_drag.py} provides the
drag-model integrator for Model~B.

\textbf{Analysis}. The reader fits two models to the data:

\begin{itemize}[leftmargin=*]
  \item Model A: no air resistance. Predicts $y(x)$ as a parabola
    with $y$-intercept and apex determined by initial conditions;
    fit two free parameters (initial speed and launch angle).
  \item Model B: linear or quadratic drag. Predicts numerically-
    integrated trajectories with three free parameters (initial
    speed, launch angle, drag coefficient).
\end{itemize}

For each trajectory, the reader reports the best-fit parameters,
the residuals (measured minus predicted position), and the
goodness-of-fit ($R^2$ or chi-square).

\textbf{Reflection ($1{,}500$ to $3{,}000$ words)}. The reader
addresses:

\begin{itemize}[leftmargin=*]
  \item Which model better describes which projectile (a tennis
    ball is closer to drag-free at low speeds; a sock filled with
    rice may show more drag).
  \item How the inferred drag coefficient compares to published
    values for the projectile shape and Reynolds number.
  \item One source of measurement error in the position extraction
    that limits the fit's accuracy.
  \item One quantity the model does not capture (Magnus force on a
    spinning ball, wind, aerodynamic lift, ground rebound) and what
    extension would be needed.
\end{itemize}

\textbf{Deliverable}. The video files (or links to them), the
extracted position tables, the analysis code, the plots of measured
versus predicted trajectories, and the reflection.

\textbf{Time}. Two to four hours of recording and frame extraction,
four to eight hours of analysis and fitting, four to eight hours of
write-up.
\end{project}

\section{Exercises}

The exercises mix calculation, derivation, estimation, simulation,
design, diagnosis, failure analysis, and judgment. Solutions are
provided for calculation and derivation exercises; estimation
exercises receive published reference estimates with the editor's
working; design, diagnosis, failure-analysis, and judgment exercises
receive rubrics rather than single answers. Exercises are tagged
with their reader-path tier.

\subsection*{Calculation}\pathtag{core}

\begin{exercise}
A particle moves along the $x$-axis with position $x(t) = 3t^2 - 2t
+ 5$ (metres, with $t$ in seconds). Find the velocity and
acceleration at $t = 2\,\si{\second}$.
\paragraph{Solution sketch.} Differentiate: $v(t)=\dot{x}=6t-2$ and
$a(t)=\dot{v}=6$, constant. At $t=2\,\si{\second}$, $v=6(2)-2=
10\,\si{\meter\per\second}$ and $a=6\,\si{\meter\per\second\squared}$.
The acceleration is constant because the position is quadratic in
time; the motion is uniformly accelerated.
\end{exercise}

\begin{exercise}
A car accelerates from rest at $3\,\si{\meter\per\second\squared}$ for
$10\,\si{\second}$, then decelerates uniformly to rest over the next
$15\,\si{\second}$. Find the total distance travelled and the maximum
speed reached.
\paragraph{Solution sketch.} Phase 1: $v_{\max}=at=3\times 10=
30\,\si{\meter\per\second}$; distance $d_{1}=\tfrac{1}{2}at^{2}=
\tfrac{1}{2}(3)(10)^{2}=150\,\si{\meter}$. Phase 2 decelerates from
$30$ to rest in $15\,\si{\second}$ (deceleration $2\,\si{\meter\per
\second\squared}$); distance $d_{2}=\tfrac{1}{2}v_{\max}t_{2}=
\tfrac{1}{2}(30)(15)=225\,\si{\meter}$. Total $d=375\,\si{\meter}$,
$v_{\max}=30\,\si{\meter\per\second}$. The total is the area under the
velocity triangle of
Figure~\ref{fig:vol03ch04:kinematics-1d}.
\end{exercise}

\begin{exercise}
A wheel of radius $0.5\,\si{\meter}$ rolls without slipping along a
horizontal surface with linear velocity $4\,\si{\meter\per\second}$.
Find the angular velocity of the wheel and the velocity of the top
point of the wheel (the point diametrically opposite the
ground-contact point).
\paragraph{Solution sketch.} Rolling without slipping gives
$v_{G}=\omega R$, so $\omega=v_{G}/R=4/0.5=8\,\si{\radian\per
\second}$. The contact point is the instantaneous centre, so the top
point sits at twice the centre's distance from it and moves at
$2v_{G}=8\,\si{\meter\per\second}$, horizontal and forward
(Figure~\ref{fig:vol03ch04:rolling-velocity}).
\end{exercise}

\begin{exercise}
A solid cylinder of mass $5\,\si{\kilogram}$ and radius $0.1\,\si{
\meter}$ rolls without slipping down an inclined plane at
$15^{\circ}$. Find its linear acceleration.
\paragraph{Solution sketch.} For a solid cylinder, $I_{G}=\tfrac{1}{2}
mR^{2}$, so the worked result $a_{G}=\dfrac{g\sin\beta}{1+I_{G}/
(mR^{2})}=\tfrac{2}{3}g\sin\beta$ applies. With $\beta=15^{\circ}$,
$a_{G}=\tfrac{2}{3}(9.81)\sin 15^{\circ}=\tfrac{2}{3}(9.81)(0.259)=
1.69\,\si{\meter\per\second\squared}$. The mass and radius cancel; the
acceleration depends only on the angle and the shape factor.
\end{exercise}

\begin{exercise}
A projectile is launched with speed $30\,\si{\meter\per\second}$ at
$40^{\circ}$ above horizontal. Find its range, time of flight, and
maximum height (no air resistance).
\paragraph{Solution sketch.} With $v_{0}=30\,\si{\meter\per\second}$,
$\theta=40^{\circ}$, $g=9.81\,\si{\meter\per\second\squared}$:
time of flight $t=2v_{0}\sin\theta/g=2(30)(0.643)/9.81=3.93\,\si{
\second}$; range $R=v_{0}^{2}\sin 2\theta/g=900\sin 80^{\circ}/9.81=
90.4\,\si{\meter}$; maximum height $h=v_{0}^{2}\sin^{2}\theta/(2g)=
900(0.413)/19.62=18.9\,\si{\meter}$. All three follow from the
componentwise equations of motion.
\end{exercise}

\begin{exercise}
A $1{,}500\,\si{\kilogram}$ car moving at $20\,\si{\meter\per\second}$
collides with a stationary $1{,}000\,\si{\kilogram}$ car. The cars
stick together. Find the speed of the combined wreck immediately
after impact and the fraction of kinetic energy lost.
\paragraph{Solution sketch.} Momentum conservation:
$v_{f}=m_{1}v_{1}/(m_{1}+m_{2})=1500(20)/2500=12\,\si{\meter\per
\second}$. Kinetic energy before $=\tfrac{1}{2}(1500)(20)^{2}=
300\,\si{\kilo\joule}$; after $=\tfrac{1}{2}(2500)(12)^{2}=180\,\si{
\kilo\joule}$; fraction lost $=120/300=0.40$, that is $40\,\%$. The
lost energy is the mass ratio: the inelastic loss equals
$m_{2}/(m_{1}+m_{2})$ of the incident kinetic energy
(Figure~\ref{fig:vol03ch04:collision}).
\end{exercise}

\begin{exercise}
A bullet of mass $10\,\si{\gram}$ travelling at $400\,\si{\meter\per
\second}$ embeds in a $4\,\si{\kilogram}$ pendulum block. Find the
speed of the block immediately after the bullet embeds and the
height to which the block rises on its pendulum chain.
\paragraph{Solution sketch.} The ballistic pendulum has two stages.
Embedding is inelastic, so use momentum: $v=m_{b}v_{b}/(m_{b}+M)=
0.010(400)/4.010=0.998\,\si{\meter\per\second}$. The swing is then
energy-conserving: $h=v^{2}/(2g)=(0.998)^{2}/19.62=0.0508\,\si{\meter}
\approx 51\,\si{\milli\meter}$. The two stages use different
conservation laws by design: momentum across the fast embedding,
energy across the slow swing.
\end{exercise}

\begin{exercise}
A car of mass $1{,}500\,\si{\kilogram}$ on a curve of radius
$80\,\si{\meter}$ travels at $25\,\si{\meter\per\second}$. Find the
centripetal acceleration and the lateral friction force the road
must supply.
\paragraph{Solution sketch.} Centripetal acceleration $a_{c}=v^{2}/r=
625/80=7.81\,\si{\meter\per\second\squared}$, about $0.80g$. The road
supplies it through tyre friction: $F=ma_{c}=1500(7.81)=
11.7\,\si{\kilo\newton}$. The required friction coefficient is
$\mu=a_{c}/g=0.80$, near the dry-asphalt limit, so this corner is
close to the no-skid edge; on a wet road ($\mu\approx 0.4$) the car
could not hold it at this speed.
\end{exercise}

\subsection*{Derivation}\pathtag{standard}

\begin{exercise}
Derive the rocket equation: for a rocket of instantaneous mass $m(t)$
expelling mass at rate $\dot{m}$ with exhaust velocity $v_e$ relative
to the rocket, in zero gravitational field, show that $m\dot{v} =
\dot{m} v_e$, and that the velocity change after a finite burn is
$\Delta v = v_e \ln(m_0/m_f)$.
\paragraph{Solution sketch.} Conserve momentum over $dt$. At time $t$
the rocket has mass $m$ and velocity $v$; in $dt$ it ejects mass
$-dm>0$ at velocity $v-v_{e}$ (the exhaust speed relative to the
rocket, behind it). Total momentum is conserved (no external force):
$mv=(m+dm)(v+dv)+(-dm)(v-v_{e})$. Expanding and dropping the
second-order term $dm\,dv$ gives $m\,dv=-v_{e}\,dm$, that is
$m\dot{v}=\dot{m}_{\text{loss}}v_{e}$ with $\dot{m}_{\text{loss}}=
-dm/dt$. Separate and integrate from $m_{0}$ to $m_{f}$:
$\int dv=-v_{e}\int dm/m$, so $\Delta v=v_{e}\ln(m_{0}/m_{f})$. The
logarithm is the tyranny of the rocket equation: each fixed increment
of $\Delta v$ costs a fixed mass \emph{ratio}, so high $\Delta v$
demands exponentially large propellant fractions.
\end{exercise}

\begin{exercise}
Derive the work-energy theorem from Newton's second law for a
constant-mass particle. State each step explicitly and identify the
geometric identity $\vec{a}\cdot d\vec{r} = d(\tfrac{1}{2} v^2)$ at
its point of use.
\paragraph{Solution sketch.} Start from $\vec{F}=m\vec{a}$. Dot both
sides with the displacement $d\vec{r}=\vec{v}\,dt$:
$\vec{F}\cdot d\vec{r}=m\vec{a}\cdot\vec{v}\,dt$. The geometric
identity is $\vec{a}\cdot\vec{v}=\tfrac{d}{dt}(\tfrac{1}{2}\vec{v}\cdot
\vec{v})=\tfrac{d}{dt}(\tfrac{1}{2}v^{2})$, since $\tfrac{d}{dt}(\vec{v}
\cdot\vec{v})=2\vec{v}\cdot\dot{\vec{v}}=2\vec{v}\cdot\vec{a}$. Hence
$\vec{F}\cdot d\vec{r}=m\,d(\tfrac{1}{2}v^{2})$. Integrate from $A$ to
$B$: the left side is $W_{\text{net}}$, the right side is $\tfrac{1}{2}
mv_{B}^{2}-\tfrac{1}{2}mv_{A}^{2}=\Delta T$. The identity is used
exactly once, at the conversion of $\vec{a}\cdot\vec{v}$ into a total
time derivative.
\end{exercise}

\begin{exercise}
For a rigid body in plane motion, derive the energy expression $T =
\tfrac{1}{2} m v_G^2 + \tfrac{1}{2} I_G \omega^2$ from the integral
$T = \tfrac{1}{2}\int v^2 \,dm$ over the body.
\paragraph{Solution sketch.} Write the velocity of a mass element as
the centre-of-mass velocity plus the rotational part:
$\vec{v}=\vec{v}_{G}+\vec{\omega}\times\vec{r}'$, where $\vec{r}'$ is
the element's position from $G$. Then $v^{2}=v_{G}^{2}+2\vec{v}_{G}
\cdot(\vec{\omega}\times\vec{r}')+|\vec{\omega}\times\vec{r}'|^{2}$.
Integrate over the body. The first term gives $\tfrac{1}{2}mv_{G}^{2}$.
The cross term integrates to zero because $\int\vec{r}'\,dm=\vec{0}$ by
the definition of the centre of mass. The third term gives $\tfrac{1}{2}
\int\omega^{2}r'^{2}\,dm=\tfrac{1}{2}I_{G}\omega^{2}$ (for planar
motion $|\vec{\omega}\times\vec{r}'|=\omega r'_{\perp}$ and $I_{G}=
\int r'^{2}_{\perp}\,dm$). Summing the surviving terms gives the
result. The vanishing cross term is why the split into translation and
rotation is clean only about the centre of mass.
\end{exercise}

\begin{exercise}
Derive the precession rate $\Omega_p = mgr_G/(I\omega)$ for a
spinning top, using the approximation that the angular momentum
points along the spin axis. State the conditions under which the
approximation is accurate.
\paragraph{Solution sketch.} Take $\vec{L}\approx I\vec{\omega}$ along
the spin axis, magnitude $L=I\omega$. The weight produces a moment
about the support of magnitude $M=mgr_{G}\sin\theta$, horizontal and
perpendicular to the axis. Since $d\vec{L}/dt=\vec{M}$ and $\vec{M}$ is
perpendicular to $\vec{L}$, the moment rotates $\vec{L}$ without
changing its magnitude; the tip of $\vec{L}$ traces a horizontal
circle of radius $L\sin\theta$ at angular rate $\Omega_{p}$, so
$|d\vec{L}/dt|=\Omega_{p}L\sin\theta$. Equate to $M$:
$\Omega_{p}=M/(L\sin\theta)=mgr_{G}/(I\omega)$, independent of
$\theta$. The approximation holds when the spin angular momentum
dominates the precessional and nutational contributions, that is when
$I\omega^{2}\gg mgr_{G}$ (fast spin); otherwise the neglected terms
produce nutation and the simple formula degrades.
\end{exercise}

\begin{exercise}
Derive the equation of motion of a damped harmonic oscillator
$m\ddot{x} + c\dot{x} + kx = 0$ from Newton's second law applied to a
mass on a spring with viscous damping. Identify the damping ratio
$\zeta = c/(2\sqrt{mk})$ and discuss the three damping regimes.
\paragraph{Solution sketch.} Forces on the mass: spring $-kx$, viscous
damper $-c\dot{x}$. Newton's second law $m\ddot{x}=-kx-c\dot{x}$
rearranges to $m\ddot{x}+c\dot{x}+kx=0$. Divide by $m$:
$\ddot{x}+2\zeta\omega_{0}\dot{x}+\omega_{0}^{2}x=0$ with
$\omega_{0}=\sqrt{k/m}$ and $\zeta=c/(2\sqrt{mk})=c/(2m\omega_{0})$.
The characteristic roots are $\lambda=\omega_{0}(-\zeta\pm\sqrt{\zeta^{2}
-1})$. Three regimes follow from the discriminant: underdamped
($\zeta<1$, complex roots, decaying oscillation), critically damped
($\zeta=1$, repeated real root, fastest non-oscillatory return), and
overdamped ($\zeta>1$, two real roots, slow non-oscillatory return).
The closed forms and settling times are computed in
\texttt{damped\_oscillator.py}; the regimes are plotted by the
simulation exercise below.
\end{exercise}

\subsection*{Estimation}\pathtag{core}

\begin{exercise}
Estimate the kinetic energy of a tennis ball served at $200\,\si{
\kilo\meter\per\hour}$. Compare to the energy required to lift the
ball $10\,\si{\meter}$. State the ball's mass assumed.
\paragraph{Solution sketch.} Tennis ball $m\approx 0.057\,\si{
\kilogram}$; $200\,\si{\kilo\meter\per\hour}=55.6\,\si{\meter\per
\second}$. Kinetic energy $T=\tfrac{1}{2}(0.057)(55.6)^{2}\approx
88\,\si{\joule}$. Lifting $10\,\si{\meter}$ needs $mgh=0.057(9.81)(10)
\approx 5.6\,\si{\joule}$. The serve carries about $16$ times the
lift energy: a fast serve is an energetic event, which is why the
return must be timed and braced rather than merely placed.
\end{exercise}

\begin{exercise}
Estimate the impulse delivered by a hammer striking a nail to drive
it $5\,\si{\milli\meter}$ into a wooden plank. State assumptions about
hammer mass, swing speed, and friction force on the nail.
\paragraph{Solution sketch.} Assume hammer head $m\approx 0.5\,\si{
\kilogram}$, impact speed $v\approx 6\,\si{\meter\per\second}$. The
impulse delivered to the nail roughly equals the hammer's momentum
change if it rebounds little: $J\approx mv=0.5(6)=3\,\si{\newton
\second}$. A cross-check on the force: work done equals the nail's
penetration resistance times distance. With $J\approx mv$ and a
contact time $\Delta t\sim 1\,\si{\milli\second}$, the mean force is
$F\approx J/\Delta t\approx 3000\,\si{\newton}$, far above the
hammer's $5\,\si{\newton}$ weight; over $5\,\si{\milli\meter}$ that
force does $\sim 15\,\si{\joule}$, consistent with the hammer's
kinetic energy $\tfrac{1}{2}(0.5)(6)^{2}=9\,\si{\joule}$ to within the
estimate's tolerance. Impulse $\sim 3\,\si{\newton\second}$, peak
force of order kilonewtons.
\end{exercise}

\begin{exercise}
Estimate the thrust required for a $1{,}000\,\si{\kilogram}$ vehicle
to accelerate from $0$ to $100\,\si{\kilo\meter\per\hour}$ in
$10\,\si{\second}$ on a level road, ignoring drag. Compare to the
typical thrust of a passenger-car engine.
\paragraph{Solution sketch.} $100\,\si{\kilo\meter\per\hour}=27.8\,
\si{\meter\per\second}$; mean acceleration $a=27.8/10=2.78\,\si{
\meter\per\second\squared}$. Thrust $F=ma=1000(2.78)=2.78\,\si{\kilo
\newton}$. This is the tractive force at the wheels; the corresponding
power at top speed is $Fv=2780(27.8)\approx 77\,\si{\kilo\watt}$,
about $100$ horsepower, squarely in the range of an ordinary family
car. The estimate also bounds the tyre grip required: $a/g=0.28$, well
within the dry-road friction limit, so traction is not the
constraint here, engine power is.
\end{exercise}

\begin{exercise}
Estimate the time and distance required for a freight train of mass
$5{,}000\,\si{\tonne}$ to brake from $80\,\si{\kilo\meter\per\hour}$
to rest. State assumptions about the friction between wheels and
rails ($\mu \approx 0.1$ to $0.2$).
\paragraph{Solution sketch.} Take $\mu\approx 0.1$ (steel on steel,
the low end for wet or worn rail). Deceleration $a=\mu g\approx
0.1(9.81)=0.98\,\si{\meter\per\second\squared}$, independent of mass.
$80\,\si{\kilo\meter\per\hour}=22.2\,\si{\meter\per\second}$. Time
$t=v/a=22.2/0.98\approx 23\,\si{\second}$; distance $d=v^{2}/(2a)=
(22.2)^{2}/1.96\approx 250\,\si{\meter}$. With $\mu=0.2$ both halve to
$\sim 11\,\si{\second}$ and $\sim 125\,\si{\meter}$. The mass cancels,
so the freight train and a light car stop in similar distances at the
same friction limit; the train's hazard is that signalling distances
must respect the long absolute stopping distance and the low grip.
\end{exercise}

\begin{exercise}
Estimate the angular velocity of a typical desk-mounted ceiling fan
in radians per second and the linear speed at its blade tip. State
assumptions about blade length and RPM.
\paragraph{Solution sketch.} Assume a ceiling fan at $300$ revolutions
per minute with blade tip radius $0.6\,\si{\meter}$. Angular velocity
$\omega=300\times 2\pi/60=31.4\,\si{\radian\per\second}$. Tip speed
$v=\omega R=31.4(0.6)\approx 19\,\si{\meter\per\second}$, about
$68\,\si{\kilo\meter\per\hour}$. The centripetal acceleration at the
tip is $\omega^{2}R=31.4^{2}(0.6)\approx 590\,\si{\meter\per\second
\squared}\approx 60g$, which is why a loose blade or a bug on the tip
is flung off with force; the same $\omega^{2}R$ scaling reappears in
the centrifuge design exercise at much higher speed.
\end{exercise}

\subsection*{Simulation}\pathtag{mastery}

\begin{exercise}
Simulate the trajectory of a projectile launched at $45^{\circ}$ with
speed $30\,\si{\meter\per\second}$, with (a) no air resistance, (b)
linear drag, and (c) quadratic drag. Compare the three trajectories;
identify the speed regime in which linear drag dominates and the
regime in which quadratic drag dominates.
\paragraph{Solution sketch.} The starter code is
\texttt{projectile\_drag.py}: integrate $\ddot{\vec{r}}=-g\hat{\jmath}
-(k/m)f(\vec{v})$ with RK4, where $f=\vec{v}$ for linear drag and
$f=|\vec{v}|\vec{v}$ for quadratic. For the launch given, the
drag-free range is $91.7\,\si{\meter}$, linear-drag $43.3\,\si{
\meter}$, quadratic-drag $38.9\,\si{\meter}$
(Figure~\ref{fig:vol03ch04:projectile-drag}); the quadratic case is
visibly fore-aft asymmetric. The regime is set by the Reynolds number
$\mathrm{Re}=\rho v d/\mu$: linear (Stokes) drag holds for
$\mathrm{Re}\lesssim 1$, the creeping-flow limit of tiny or very slow
bodies; quadratic (form) drag holds for $\mathrm{Re}\gtrsim 10^{3}$,
which covers every ordinary thrown object in air. A good answer
states that linear drag is the wrong model for any ordinary
projectile and is included only to show the contrast.
\end{exercise}

\begin{exercise}
Simulate the motion of a damped harmonic oscillator with $m =
1\,\si{\kilogram}$, $k = 100\,\si{\newton\per\meter}$, and damping
ratios $\zeta = 0.1$, $\zeta = 1.0$, $\zeta = 2.0$. Plot the
displacement versus time for the same initial conditions. Identify
the underdamped, critically damped, and overdamped regimes.
\paragraph{Solution sketch.} Use \texttt{damped\_oscillator.py},
which evaluates the closed-form response in each regime; or integrate
$\ddot{x}+2\zeta\omega_{0}\dot{x}+\omega_{0}^{2}x=0$ numerically with
$\omega_{0}=\sqrt{k/m}=10\,\si{\radian\per\second}$. From $x(0)=x_{0}$,
$\dot{x}(0)=0$: $\zeta=0.1$ rings and decays (underdamped),
$\zeta=1.0$ returns fastest with no overshoot (critical), $\zeta=2.0$
returns slowly with no overshoot (overdamped). The $2\,\%$ settling
times from the script are about $3.8$, $0.6$, and $1.5\,\si{\second}$:
critical damping settles fastest, the design target for instrument
needles and door closers. Chapter 7 of Volume II plots the same three
regimes for the formal second-order linear ODE.
\end{exercise}

\begin{exercise}
Implement a numerical solver for the two-body central-force problem
(Newtonian gravity). Simulate Earth's orbit around the Sun and
verify that the orbit is closed and approximately elliptical.
Compare the simulated orbital period to Kepler's third law.
\paragraph{Solution sketch.} The reference solver is
\texttt{two\_body.py}, which integrates $\ddot{\vec{r}}=-GM\vec{r}/
r^{3}$ with a velocity-Verlet (symplectic) step. Starting from
perihelion with Earth's perihelion distance and speed, the orbit
closes and the relative energy drift over $1.2$ years is of order
$10^{-9}$; an ordinary RK4 step would instead let the energy drift and
the orbit spiral, the lesson of why integrator choice matters for
long runs. Recover the semi-major axis from the energy and check
$T^{2}=4\pi^{2}a^{3}/(GM)$; the simulated period matches Kepler to
four figures (Figure~\ref{fig:vol03ch04:central-orbit}). A complete
answer reports the energy drift, the closure of the orbit, and the
$T^{2}/a^{3}$ check against \texttt{orbit\_period.csv}.
\end{exercise}

\subsection*{Design}\pathtag{standard}

\begin{exercise}
Design a passenger-car braking system to stop a $1{,}500\,\si{
\kilogram}$ vehicle from $130\,\si{\kilo\meter\per\hour}$ in
$60\,\si{\meter}$ on dry asphalt. Specify the required average
deceleration, the average braking force, and the brake-disc
diameter, assuming a peak rotor temperature of $400\,\si{\degreeCelsius}$.
\paragraph{Solution sketch.} A model answer fixes the kinematics
first: $130\,\si{\kilo\meter\per\hour}=36.1\,\si{\meter\per\second}$;
required deceleration $a=v^{2}/(2d)=(36.1)^{2}/(2\times 60)=10.9\,\si{
\meter\per\second\squared}\approx 1.1g$. Average braking force
$F=ma=1500(10.9)=16.3\,\si{\kilo\newton}$. Check the friction limit:
$a/g=1.11>0.8$ exceeds plain dry-asphalt grip, so the specification is
feasible only with high-grip tyres or downforce, a point the answer
should flag. The energy to dissipate is $\tfrac{1}{2}mv^{2}=
\tfrac{1}{2}(1500)(36.1)^{2}\approx 980\,\si{\kilo\joule}$; sizing the
disc to absorb this within a $400\,\si{\degreeCelsius}$ rise sets the
rotor mass through $Q=mc\Delta T$ (cast iron $c\approx 460\,\si{\joule
\per\kilogram\per\kelvin}$), giving a few kilograms per disc and a
diameter in the $300$ to $350\,\si{\milli\meter}$ range. The rubric
rewards the friction-limit check and the energy-to-thermal-mass link.
\end{exercise}

\begin{exercise}
Design a passive vibration isolator to reduce the bearing-loading
from a $50\,\si{\kilogram}$ rotor with residual imbalance $50\,\si{
\gram}$ at a radius of $100\,\si{\milli\meter}$ rotating at
$3{,}000\,\text{RPM}$. State the isolator's natural frequency and
the resulting transmissibility at the rotation rate.
\paragraph{Solution sketch.} Rotation rate $\omega=3000\times 2\pi/60=
314\,\si{\radian\per\second}$ ($50\,\si{\hertz}$). The imbalance force
is $m_{r}e\,\omega^{2}=0.050(0.100)(314)^{2}\approx 494\,\si{\newton}$
(this is $m_{r}=50\,\si{\gram}$ at $r=100\,\si{\milli\meter}$,
matching the $m e$ product). Isolation requires the frequency ratio
$r=\omega/\omega_{n}>\sqrt{2}$; choose $\omega_{n}$ at one quarter of
$\omega$, that is $f_{n}\approx 12.5\,\si{\hertz}$, giving $r=4$. With
light damping ($\zeta=0.05$) the transmissibility is $T=\sqrt{(1+(2
\zeta r)^{2})/((1-r^{2})^{2}+(2\zeta r)^{2})}\approx 0.07$, so the
transmitted force drops to about $35\,\si{\newton}$. The figures are
produced by \texttt{imbalance\_response.py}. The rubric rewards
choosing $\omega_{n}$ well below the forcing frequency and noting the
penalty of passing through resonance during run-up.
\end{exercise}

\begin{exercise}
Design a centrifuge rotor to spin a sample of mass $0.5\,\si{\kilogram}$
at $4{,}000\,\text{RPM}$ at a rotation radius of $0.15\,\si{\meter}$.
State the centripetal acceleration the sample experiences, in
multiples of $g$, and the force on the sample tube.
\paragraph{Solution sketch.} $\omega=4000\times 2\pi/60=419\,\si{
\radian\per\second}$. Centripetal acceleration $a_{c}=\omega^{2}R=
(419)^{2}(0.15)=2.63\times 10^{4}\,\si{\meter\per\second\squared}=
2680g$ (the relative centrifugal field). Force on the tube $F=ma_{c}=
0.5(2.63\times 10^{4})=13.2\,\si{\kilo\newton}$, carried by the rotor
arm in tension. The rubric rewards reporting the field in multiples of
$g$ (the units centrifuge users actually quote) and noting that the
$\omega^{2}R$ tension sets the rotor's material strength requirement,
the governing constraint at high speed.
\end{exercise}

\begin{exercise}
Design a flywheel to store $1\,\text{kWh}$ of kinetic energy.
Choose between a steel disc and an aluminum disc; specify the
geometry and the operating speed; state which design constraint
(material strength, bearing capacity, energy density) governs.
\paragraph{Solution sketch.} $1\,\text{kWh}=3.6\,\si{\mega\joule}$.
For a solid disc, $E=\tfrac{1}{2}I\omega^{2}=\tfrac{1}{4}mR^{2}
\omega^{2}$. The hard limit is material strength: the peak rim stress
in a spinning disc is $\sigma\approx\rho\,(\omega R)^{2}$ up to a
shape factor, so the storable energy per unit mass scales as
$\sigma_{\text{allow}}/\rho$, the specific strength. Steel and
aluminium have similar specific strength, so they store similar
energy per kilogram; the choice turns on cost, fatigue, and burst
containment rather than energy density. A defensible answer picks a
modest rim speed (a few hundred metres per second, below the burst
limit with a safety factor), solves $\tfrac{1}{4}mR^{2}\omega^{2}=
3.6\,\si{\mega\joule}$ for the mass-radius-speed combination, and
states that material strength governs, with bearing capacity a
secondary constraint and energy density a consequence rather than a
free parameter.
\end{exercise}

\subsection*{Diagnosis and reverse engineering}\pathtag{core}

\begin{exercise}
A washing machine vibrates noticeably during spin cycles when its
load is unevenly distributed. State the dynamics responsible for the
vibration and the design feature (active or passive) that the
machine should incorporate to mitigate it.
\paragraph{Solution sketch.} An uneven load displaces the drum's
centre of mass from the spin axis by an eccentricity $e$, producing a
rotating imbalance force $m e\,\omega^{2}$ at the spin frequency, the
same mechanism as machinery imbalance. A strong answer names the
mechanism, notes the $\omega^{2}$ growth that makes the high-speed
spin cycle the worst case, and lists mitigations: passive measures
(soft suspension springs and friction or fluid dampers tuned so the
spin frequency sits well above the suspension natural frequency, into
the isolation regime) and active measures (fluid-filled balance rings
or counterweight chambers, and a controller that detects imbalance at
low speed and ramps only when the load has redistributed).
\end{exercise}

\begin{exercise}
A road vehicle experiences a steering-wheel oscillation at a
particular speed band. State the most probable dynamic cause (tyre
imbalance, suspension geometry, brake-rotor warp), the test that
would distinguish among them, and the corrective action for the
most likely candidate.
\paragraph{Solution sketch.} A speed-dependent steering oscillation
points to a rotating-imbalance source, because the forcing frequency
tracks wheel speed and a resonance shows up in a band. The most
probable cause is wheel or tyre imbalance (a one-times-wheel-speed
shake), with brake-rotor warp (felt under braking, at brake-disc
frequency) and worn suspension or steering joints as alternatives. The
distinguishing test is a road test across the speed band plus an
on-car or off-car wheel balance check, and feeling whether the shake
appears only under braking. The corrective action for the most likely
candidate is to rebalance the wheels (and inspect for tyre
out-of-roundness or a thrown wheel weight). A good answer ties the
band-limited symptom to resonance and the speed tracking to a rotating
source.
\end{exercise}

\begin{exercise}
A small drone hovers stably in still air but accelerates
unpredictably when commanded to translate. State the most probable
dynamic source of the deviation (motor imbalance, propeller
imbalance, control-loop tuning, airframe asymmetry) and the
diagnostic procedure.
\paragraph{Solution sketch.} Stable hover but poor translation
separates a vibration problem from a control problem. Motor or
propeller imbalance shows as a steady vibration at rotor frequency,
present in hover too, so the fact that hover is stable argues against
it. The likeliest source is control-loop tuning: the attitude or
position controller's gains, well behaved at the hover trim, become
sluggish or oscillatory when the demanded tilt changes the thrust
vector and the dynamics. The diagnostic procedure is to log the
inertial-measurement data and motor commands during a commanded
translation, look for control-loop oscillation or lag rather than a
clean rotor-frequency line, and bench-test each propeller for
imbalance to rule out the mechanical cause. A good answer reasons from
``stable in hover'' to demote the imbalance hypotheses.
\end{exercise}

\begin{exercise}
A bicycle wheel rotated by hand wobbles visibly at low speed but
appears stable at high speed. Identify the dynamic phenomenon and
state whether the wobble is dangerous, harmless, or harmless-now-
but-symptomatic of a future issue.
\paragraph{Solution sketch.} The low-speed wobble is the visible
imbalance and runout of the wheel: at low spin rate the angular
momentum $L=I\omega$ is small, so the gravitational and contact
moments deflect the wheel's axis noticeably. At high speed the
gyroscopic stiffness $L=I\omega$ grows, so the same disturbing moments
produce a smaller precession and the wheel appears steady; this is
gyroscopic stabilisation, the same effect that keeps a moving bicycle
upright. A good answer identifies the phenomenon as gyroscopic
stabilisation of a slightly imbalanced or untrued wheel and judges it
harmless-now-but-symptomatic: the underlying imbalance or rim runout
persists, loads the bearings cyclically at speed, and should be trued
and balanced before it wears the hub.
\end{exercise}

\subsection*{Failure analysis}\pathtag{core}

\begin{exercise}
A high-speed turbine fan blade fractures at its hub fillet after
many hours of service. State the dominant failure mechanism
(imbalance-driven fatigue, foreign-object damage, cyclic thermal
stress, or other), the calculation a design review should have
performed, and the inspection procedure that should have detected
the failure precursor.
\paragraph{Solution sketch.} The dominant mechanism is high-cycle
fatigue driven by the cyclic stress at the hub fillet, where the
rotation-frequency bending load from residual imbalance and the
alternating aerodynamic and vibratory loads concentrate. The design
review should have performed a fatigue calculation: the alternating
stress at the fillet (with the stress-concentration factor of the
geometry), placed on the material's stress-life curve at the expected
cycle count, with a margin for resonance crossings. The inspection
that should have caught the precursor is periodic non-destructive
examination of the fillet (dye-penetrant or eddy-current crack
detection) together with vibration trending that would show a rising
one-times line as imbalance grew. A good answer separates the
load-creating mechanism (imbalance and resonance) from the
damage-accumulating mechanism (fatigue at the concentration).
\end{exercise}

\begin{exercise}
A rotating crankshaft fails at a journal-bearing transition fillet
during a routine engine start-up. State, in $200$ to $300$ words, the
likely failure sequence (imbalance amplification through critical
speed; cyclic bending stress at the fillet; fatigue propagation;
final overload) and the design or maintenance discipline that would
have caught the precursor.
\paragraph{Solution sketch.} A model answer narrates the sequence in
order: residual imbalance generates a rotating force $m e\,\omega^{2}$;
during start-up the rotor passes through a critical speed where the
rotor-bearing-foundation system resonates and the imbalance force is
amplified by the system's quality factor, momentarily raising the
cyclic bending stress at the journal-transition fillet; the
stress-concentration there seeds a fatigue crack that propagates a
little on each start-up; after enough cycles the remaining section can
no longer carry the peak load and fails by final overload, often
during a start-up when the transient stress is highest. The discipline
that catches the precursor is balancing to the ISO~1940/1 grade plus a
generous fillet radius at design time, and vibration trending through
the critical-speed crossing plus periodic crack inspection of the
fillet in service. The answer should keep within $200$ to $300$ words
and name both a design and a maintenance control.
\end{exercise}

\begin{exercise}
A bearing in a large industrial fan housing develops audible
rumbling and high vibration over six months of service. State the
most probable failure mechanism, the maintenance discipline (vibration
monitoring, oil analysis, thermography) that should have caught the
trend earlier, and the calculation that translates the bearing's
imbalance amplitude into a residual-life estimate.
\paragraph{Solution sketch.} The audible rumble and rising vibration
over months are the signature of rolling-element bearing fatigue:
sub-surface cracks spall the races and elements, raising the
broadband and bearing-defect-frequency vibration. The maintenance
discipline that should have caught the trend is vibration monitoring
with trending (the defect frequencies, computed from the bearing
geometry, rise above the noise floor weeks to months before failure),
backed by oil analysis (wear-metal particles) and thermography (a hot
bearing). The residual-life calculation is the bearing rating-life
estimate: $L_{10}=(C/P)^{p}$ revolutions, with $C$ the dynamic load
rating, $P$ the equivalent dynamic load (which the imbalance amplitude
raises), and $p=3$ for ball bearings; the rising measured load
shortens the predicted $L_{10}$ and converts the vibration amplitude
into a remaining-revolutions, and hence remaining-time, figure. A good
answer links the measured amplitude to $P$ and $P$ to $L_{10}$.
\end{exercise}

\subsection*{Judgment}\pathtag{standard}

\begin{exercise}
A junior engineer argues that air resistance can be neglected for any
projectile launched at less than $30\,\si{\meter\per\second}$. Write
a one-paragraph reply that takes the argument seriously, states the
conditions under which it is acceptable, and states the conditions
under which it fails.
\paragraph{Solution sketch.} A model reply grants the kernel of truth:
for a dense, compact, slow body over a short range, the drag impulse
is small compared to the gravitational impulse and the parabola is a
good predictor; a thrown stone or a shot put at modest speed fits.
Then it states the speed-blind error in the rule. What matters is not
the launch speed alone but the ratio of drag force to weight, set by
the body's size, density, and the air, summarised by the Reynolds
number and the ballistic coefficient. A light or large body (a foam
ball, a shuttlecock, a beach ball) at the same $30\,\si{\meter\per
\second}$ is dominated by drag and never follows the parabola. The
rule is acceptable only after checking that $\rho_{\text{air}}C_{d}A
v^{2}/2\ll mg$ over the flight; it fails whenever that inequality does
not hold, regardless of the absolute speed.
\end{exercise}

\begin{exercise}
A senior colleague argues that ISO balance grades are conservative
and that a slightly out-of-spec rotor will run fine. Write a
one-paragraph reply that addresses the cost-benefit trade-off and
states a defensible position.
\paragraph{Solution sketch.} A defensible reply concedes that the
grades carry margin and that a marginally out-of-spec rotor will
usually run without immediate failure, then frames the real
trade-off. The imbalance force scales as $\omega^{2}$ and acts
cyclically for the machine's whole life, so the cost of being out of
spec is not a sudden failure but accelerated bearing fatigue, higher
maintenance frequency, and the risk of crossing a resonance with a
larger force; the benefit of skipping a balancing pass is a few
minutes of shop time. For a low-speed, lightly loaded, easily replaced
machine the engineer may accept the deviation with monitoring; for a
high-speed, safety-relevant, or hard-to-service machine the asymmetry
of the trade (small saving now against a long-tailed reliability cost)
argues for holding the spec. The position to state is that the grade
is a design specification, not a suggestion, and deviations are a
documented engineering decision with monitoring, not a default.
\end{exercise}

\begin{exercise}
Identify a piece of rotating machinery in your immediate environment
(fan, motor, household appliance) and state, in $200$ to $300$ words,
the rotation speed, the approximate residual imbalance you can
detect (by feel or by sound), and the failure mode you would expect
if the imbalance trend worsens.
\paragraph{Solution sketch.} An answer in $200$ to $300$ words picks a
concrete machine (a desk fan, a kitchen extractor, a washing-machine
drum) and reasons quantitatively from observation. State the rotation
speed from the nameplate or by counting (a desk fan near $1{,}200$ to
$1{,}500$ revolutions per minute, a drum at $1{,}000$ to $1{,}600$ on
spin). Estimate the residual imbalance from the felt vibration:
relate the perceptible shake to a force $m e\,\omega^{2}$ and back out
$e$, noting that a barely felt vibration corresponds to a small
fraction of a millimetre of eccentricity. Then name the expected
failure mode as the trend worsens: bearing wear first (rumble, play),
then mount loosening or fatigue at the worst stress concentration, and
in a flimsy structure a walking or resonance event. A strong answer
ties the felt magnitude to the $\omega^{2}$ scaling and names a
specific bearing-or-mount failure rather than a generic ``it breaks.''
\end{exercise}
